%%%%%%%%%%%%%%%%%%%%%%%%%%%%%%%%%%%%%%%%%%%%%%%%%%%%%%%%%%%%%%%%%%%%%%%%%%%%%%%%
% preamble
%%%%%%%%%%%%%%%%%%%%%%%%%%%%%%%%%%%%%%%%%%%%%%%%%%%%%%%%%%%%%%%%%%%%%%%%%%%%%%%%

\documentclass[12pt,a4paper]{article}

% character encoding
\usepackage[utf8]{inputenc}

% font encoding (underscores)
\usepackage[T1]{fontenc}

% colors
\usepackage{xcolor}

% check marks
\usepackage{pifont}
\newcommand{\cmark}[1][green]{{\color{#1} \ding{51}}}
\newcommand{\xmark}[1][red]{{\color{#1} \ding{55}}}

% math
\usepackage{amsmath}
\usepackage{amssymb}

% hyperlinks
\usepackage{hyperref}

% tables
\usepackage{tabu}
\usepackage{longtable}
\usepackage{caption}
% \captionsetup[longtable]{skip=0.5\baselineskip}
\usepackage{multirow}

% counters
\usepackage{chngcntr}

% code
\usepackage{listings}
\lstset{
  aboveskip=\bigskipamount,
  belowskip=\bigskipamount,
  captionpos=b,
  abovecaptionskip=\bigskipamount,
  belowcaptionskip=\bigskipamount,
  escapeinside=``,
  keepspaces=true,
  columns=fullflexible,
}
\lstdefinestyle{asm}{
  % belowcaptionskip=0,
  morekeywords={[1]
    LOAD,STORE,FENCE,FLUSH,
    ADD,ADDI,SUB,SUBI,MUL,MULI,
    CMP,JMP,JZ,JNZ,JS,JNS,JNZNS,
    MEM,CAS,
    CHECK,HALT,EXIT
  },
  morecomment=[l]\#,
  basicstyle=\ttfamily,
  keywordstyle=\color{blue},
  commentstyle=\color{gray},
  % keepspaces,
  % columns=fullflexible,
  % tabsize=6,
  % xleftmargin=\parindent,
}
\lstdefinestyle{c++}{
  belowcaptionskip=1\baselineskip,
  breaklines=true,
  % keepspaces=true,
  % columns=fullflexible,
  % frame=L,
  % xleftmargin=\parindent,
  language=C++,
  showstringspaces=false,
  alsoletter={:},
  morekeywords={
    % Instruction,
    % map,
    std::map,
    std::ostringstream,
    std::set,
    % string,
    std::string,
    std::to_string,
    % unordered_map,
    % vector,
    std::vector,
    std::size_t,
    uint,
  },
  % basicstyle=\footnotesize\ttfamily,
  basicstyle=\small\ttfamily,
  keywordstyle=\sffamily\bfseries\color{blue},
  % keywordstyle=\bfseries\color{blue},
  % commentstyle=\itshape\color{purple!40!black},
  commentstyle=\itshape\color{gray},
  % identifierstyle=\color{blue},
  stringstyle=\color{red!60!black},
}
\lstdefinestyle{smtlib}{
  breaklines=true,
  morekeywords={[1]
    declare-fun,
    assert,
    not,
    and,
    or,
    xor,
    =,
    =>,
    ite,
    bvadd,
    bvsub,
    bvmul,
    store,
    select,
    extract,
    push,
    pop
  },
  morekeywords={[2]
    BitVec
  },
  alsoletter={-, =, >},
  morecomment=[l]\;,
  sensitive=false,
  basicstyle=\small\ttfamily,
  commentstyle=\color{gray},
  % keywordstyle=\bfseries
  keywordstyle=[1]\bfseries\color{green!40!black},
  keywordstyle=[2]\bfseries\color{red!40!black},
  % identifierstyle=\color{blue},
}
\lstdefinestyle{btor2}{
  % belowcaptionskip=0,
  morekeywords={[1]
    sort,
    zero,
    one,
    constd,
    write,
    input,
    and,
    or,
    nand,
    implies,
    ite,
    add,
    read,
    eq,
    neq,
  },
  morekeywords={[2]
    array,
    bitvec
  },
  morekeywords={[3]
    state,
    init,
    next,
    constraint,
    bad,
  },
  morecomment=[l]\#,
  basicstyle=\small\ttfamily,
  keywordstyle=[1]\bfseries\color{green!40!black},
  keywordstyle=[2]\bfseries\color{blue!80!black},
  keywordstyle=[3]\bfseries\color{red!40!black},
  commentstyle=\color{gray},
  % keepspaces,
  % columns=fullflexible,
  % tabsize=6,
  % xleftmargin=\parindent,
}

% tikz
\usepackage{tikz}
\usetikzlibrary{positioning,decorations.pathreplacing,calc}
\tikzset{>=latex}

% pgfplots
\usepackage{pgfplots}
\usepackage{pgfplotstable}
\pgfplotsset{compat=1.16}

% csvsimple
\usepackage{csvsimple}

% grammars
\usepackage{syntax}
\setlength{\grammarparsep}{0pt}

% algorithms
\usepackage{algorithm2e}

% string manipulation
\usepackage{xstring}

% todos
\setlength{\marginparwidth}{2cm}
\usepackage{todonotes}
\setuptodonotes{size=\footnotesize, shadow}
\usepackage{marginnote}
\renewcommand{\marginpar}{\marginnote}

%%%%%%%%%%%%%%%%%%%%%%%%%%%%%%%%%%%%%%%%%%%%%%%%%%%%%%%%%%%%%%%%%%%%%%%%%%%%%%%%
% commands
%%%%%%%%%%%%%%%%%%%%%%%%%%%%%%%%%%%%%%%%%%%%%%%%%%%%%%%%%%%%%%%%%%%%%%%%%%%%%%%%

% changes
\newcommand{\CHANGE}[1]{{\color{red}#1}}

% operators
\let\nand\barwedge

% language names
\newcommand{\CPP}{C\texttt{++}17}
\newcommand{\SMTLIB}{SMT-LIB}
\newcommand{\BTOR}{\textsc{Btor2}}

% bound
\newcommand{\BOUND}{\mathcal{K}}

% state variables
\newcommand{\HEAP}{\texttt{heap}}
\newcommand{\EXIT}{\texttt{exit}}
\newcommand{\EXITCODE}{\texttt{exit-code}}

\newcommand{\ACCU}{\texttt{accu}}
\newcommand{\MEM}{\texttt{mem}}
\newcommand{\SBADR}{\CHANGE{\texttt{adr}}}
\newcommand{\SBVAL}{\CHANGE{\texttt{val}}}
\newcommand{\SBFULL}{\CHANGE{\texttt{full}}}
\newcommand{\STMT}{\texttt{stmt}}
\newcommand{\BLOCK}{\texttt{block}}
\newcommand{\HALT}{\texttt{halt}}

% transition variables
\newcommand{\THREAD}{\texttt{thread}}
\newcommand{\EXEC}{\texttt{exec}}
\newcommand{\FLUSH}{\texttt{flush}}
\newcommand{\CHECK}{\texttt{check}}

% SMT-LIB sorts
\newcommand{\ASORT}[1][16]{\mathcal{A}^{#1}}
\newcommand{\BVSORT}[1][16]{\mathcal{B}^{#1}}

% inline code
\newcommand{\lstASM}[1]{\lstinline[style=asm]{#1}}
\newcommand{\lstCPP}[1]{\lstinline[style=c++]{#1}}
\newcommand{\lstSMTLIB}[1]{\lstinline[style=smtlib]{#1}}
\newcommand{\lstBTOR}[1]{\lstinline[style=btor2]{#1}}

% code comments
% \newcommand{\lstCMT}[1]{\color{black!50}{#1}}
\newcommand{\lstCMT}[1]{\color{gray}{#1}}

% todos
% \newcommand{\todo}[1]{
  % \begin{flushleft}
    % \color{red}
    % $\rightarrow$ #1
  % \end{flushleft}
% }

% \newcommand{\todonote}[1]{
  % $^{\color{red}#1}$
% }

%%%%%%%%%%%%%%%%%%%%%%%%%%%%%%%%%%%%%%%%%%%%%%%%%%%%%%%%%%%%%%%%%%%%%%%%%%%%%%%%
% document
%%%%%%%%%%%%%%%%%%%%%%%%%%%%%%%%%%%%%%%%%%%%%%%%%%%%%%%%%%%%%%%%%%%%%%%%%%%%%%%%

\begin{document}

% %%Andreas Neubauer, May 2016
\documentclass[12pt,a4paper]{article}
\usepackage{graphicx}
\graphicspath{{cover/}}
\textwidth 16.7cm
\textheight 25cm
\topmargin -2.7cm
\oddsidemargin 0.25cm
\parindent 0pt
\pagestyle{empty}
\begin{document}
%
% -------- only change entries beginning here ----------------------------
%
% choose language of coverpage: german or english
\newif\ifeng
\engtrue
%\engfalse
%
%
% enter the title of the thesis
%
\def\title{Title of the thesis}
%
%
% choose type of work: 0 ... Dissertation
%                      1 ... Diplomarbeit
%                      2 ... Masterarbeit
\def\type{2}
%
%
% enter name of degree, see examples below (as stated in your curriculum)
%
% e.g. Masterstudium
%\def\degree{Diplom-Ingenieur}
%\def\degree{Diplom-Ingenieurin}
\def\degree{Master of Science}
%
% e.g. Diplomstudium Lehramt
%\def\degree{Magister der Naturwissenschaften}
%\def\degree{Magistra der Naturwissenschaften}
%
% e.g. Doktorratsstudium
%\def\degree{Doktor der technischen Wissenschaften}
%\def\degree{Doktorin der technischen Wissenschaften}
%\def\degree{Doktor der Naturwissenschaften}
%\def\degree{Doktorin der Naturwissenschaften}
%
%
% enter the study (Studienrichtung, as stated in your curriculum)
%
% e.g. Diplomstudium:
%\def\study{Lehramt Mathematik}
%
% e.g. Masterstudium:
\def\study{Computer Science}
%
% e.g. Doktorratsstudium
%\def\study{Technische Wissenschaften}
%\def\study{Naturwissenschaften}
%
%
% If the names for the following entries are too long, break them into several
% lines using \\
%
%
% enter the name of the student
%
\def\name{Florian Schr\"ogendorfer}
%
%
% enter the name of the institute (use translation for english version)
% e.g.
%\def\institute{Industrial Mathematics\\ Institute}
%\def\institute{Institut f\"ur\\ Industriemathematik}
%
\def\institute{Institute for Formal Models and Verification}
%
%
% enter the name of the supervisor and first thesis examiner
% only for type 0 (Dissertation) you need a second thesis examiner
%
% for the german version you also have to enter the sex (male or female)
% of the supervisor/examiners
%
\def\supervisor{Armin Biere}
\newif\ifsupvismale
\supvismaletrue
%\supvismalefalse
%
\def\secondexaminer{Name of 2. examiner}
\newif\ifsecexmale
\secexmaletrue
%\secexmalefalse
%
%
% if there has been assistance by a further person uncomment the following line
% and enter the name.
%
\def\assist{Name of assistant}
%
%
% enter month year
% (the month when you brought it to the Prüfungs- und Anerkennungsservice)
%
\def\date{month year}
%
% do not change anything below this line
% -------------------------------------------------------------------------------
%
\def\ifundefined#1{\expandafter\ifx\csname#1\endcsname\relax}
\DeclareFontShape{OT1}{cmss}{m}{n}
  {<5><6><7><8><9><10><10.95><12><14.4><17.28><20.74><24.88><29.86><35.83>%
   <42.99><51.59><67><77.38>cmss10}{}
\DeclareFontShape{OT1}{cmss}{bx}{n}
  {<5><6><7><8><9><10><10.95><12><14.4><17.28><20.74><24.88><29.86><35.83>%
   <42.99><51.59><67><77.38>cmssbx10}{}
\makeatletter
\def\Huge{\@setfontsize\Huge{29.86pt}{36}}
\makeatother
%
\unitlength 1cm
\sffamily
\begin{picture}(16.7,0)
\ifeng
 \put(11.5,-2.5){\includegraphics[width=5.2cm]{jku_en}}
\else
 \put(11.5,-2.5){\includegraphics[width=5.2cm]{jku_de}}
\fi
\put(12.9,-4.2){\begin{minipage}[t]{3.9cm}\footnotesize%
\ifeng
 Submitted by\\
\else
 Eingereicht von\\
\fi
{\bfseries\name}%
\vskip 4mm%
\ifeng
 Submitted at\\
\else
 Angefertigt am\\
\fi
{\bfseries\institute}%
\vskip 4mm%
\ifcase\type%
 \ifeng
  Supervisor and\\ First Examiner\\
 \else
  \ifsupvismale%
   Betreuer und\\ Erstbeurteiler\\
  \else
   Betreuerin und\\ Erstbeurteilerin\\
  \fi
 \fi
 {\bfseries\supervisor}%
 \vskip 4mm%
 \ifeng
  Second Examiner\\
 \else
  \ifsecexmale%
   Zweitbeurteiler\\
  \else
   Zweitbeurteilerin\\
  \fi
 \fi
 {\bfseries\secondexaminer}%
\else
 \ifeng
  Supervisor\\
 \else
  \ifsupvismale%
   Betreuer\\
  \else
   Betreuerin\\
  \fi
 \fi
 {\bfseries\supervisor}%
\fi
\vskip 4mm%
\ifundefined{assist}\else
 \ifeng
  Co-Supervisor\\
 \else
  Mitbetreuung\\
 \fi
 {\bfseries\assist}%
\vskip 4mm%
\fi
\date
\end{minipage}}
\put(12.9,-25){\begin{minipage}[t]{3.9cm}\footnotesize%
{\bfseries JOHANNES KEPLER\\
\ifeng
 UNIVERSITY
\else
 UNIVERSIT\"AT
\fi
LINZ}\\
Altenbergerstra{\ss}e 69\\
4040 Linz, \"Osterreich\\
www.jku.at\\
DVR 0093696
\end{minipage}}
\put(0,-12.2){\begin{minipage}[b]{12cm}\Huge\bfseries\title\end{minipage}}
\put(0,-17.2){\includegraphics[width=4.4cm]{arr}}
\put(0,-18.3){\begin{minipage}[t]{12cm}%
\ifeng
 {\large\ifcase\type Doctoral \or Diploma \or Master \fi Thesis}%
 \vskip 2mm%
 to obtain the academic degree of%
 \vskip 3mm%
 {\large\degree}
 \vskip 3mm%
 in the \ifcase\type Doctoral \or Diploma \or Master's \fi Program
\else
 {\large\ifcase\type Dissertation\or Diplomarbeit\or Masterarbeit\fi}%
 \vskip 2mm%
 zur Erlangung des akademischen Grades%
 \vskip 3mm%
 {\large\degree}
 \vskip 3mm%
 im \ifcase\type Doktoratsstudium \or Diplomstudium\or Masterstudium\fi
\fi
\vskip 3mm%
{\large\study}
\end{minipage}}
\end{picture}
\end{document}


\tableofcontents

\newpage

\listoffigures

\newpage

\listoftables

\newpage

% \todo[inline]{unify lables of figures and tables}
%
% \bigskip
%
% \todo[inline]{\url{https://arxiv.org/pdf/cs/0108016.pdf}}
% \todo[inline]{\url{https://www.cl.cam.ac.uk/~pes20/weakmemory/cacm.pdf}}
% \todo[inline]{\url{https://mirrors.edge.kernel.org/pub/linux/kernel/people/paulmck/perfbook/perfbook-e2-rc9.pdf}}
% % CBMC
% \todo[inline]{\url{http://www.kroening.com/papers/tcad-sw-2008.pdf}}
% \todo[inline]{\url{http://www.kroening.com/papers/tacas2004.pdf}}
% \todo[inline]{\url{https://link.springer.com/content/pdf/10.1007/978-3-642-39799-8\_9.pdf}}
%
% \newpage

\section{Introduction}

%% sequential consistency:
%% \emph{``the result of any execution is the same as if the operations of all the processors were executed in some sequential order, and the operations of each individual processor appear in this sequence in the order specified by its program''}.

% \todo[inline]{need for software verification}\noindent

Correctness of computer systems is critical in today's information age and although software verification made considerable progress in the last decade, it is still an ongoing research topic.
Testing alone however is not sufficient to validate parallel programs
communicating via shared memory on modern multiprocessor hardware,
% running on current multiprocessor hardware
% and communicating via shared memory,
% running on modern shared memory multiprocessor hardware
as fatal race conditions can have extremely low probabilities of occurrence.
The situation becomes even worse due to counter-intuitive behaviour introduced by certain hardware optimizations.
Without further knowledge about the underlying architecture, inexperienced developers of
concurrent
low-level
systems code like lockless data structures, operating system kernels, synchronization
% libraries,
primitives,
% compilers and so on might assume that the order of access to memory is \emph{sequentially consistent}
% compilers and so on might assume a total order of access to memory for any possible interleaving, where each process is executed in program order.
compilers,
% etc.\
and so on
might assume
% that
% \emph{``the result of any execution is the same as if the operations of all the processors were executed in some sequential order for every possible interleaving and the operations of each individual processor appear in this sequence in the order specified by its program''}
% the result of any execution is the same as if the operations of all processors were executed in some sequential order and the operations of each individual processor appear program order.
\emph{sequential consistency} \cite{ref:Lamport79}.
Unfortunately, none of the major hardware architectures follows this rather restrictive \emph{memory ordering model} and allow
% memory access operations to be reordered
reordering of memory access operations
in various ways.
This opens the door for hard to find bugs caused by unexpected
% reorderings
behaviour
and
% needs to be prevented by
% requires careful use of the target architecture's memory barrier operations to ensure consistency.
requires a deep understanding of the target architecture's memory ordering habits,
% as well as
accompanied by careful use of \emph{memory barrier} operations to ensure consistency.
% referred to as \emph{sequential consistency} \cite{ref:Lamport79}.

% \bigskip

% For example, modern processors are equipped with a store buffer to speed up writes.

% \todo[inline]{counter-intuitive behaviour introduced by hardware optimizations}
% \todo[inline]{memory models}

% \bigskip

\tabulinesep=6pt
\noindent
\begin{table}[!hbt]
  \centering
  \begin{tabu}{|l|c|c|c|c|c|c|c|c|}
    \tabucline{2-}
    \multicolumn{1}{c|}{}
    & \rotatebox{90}{Alpha}
    & \rotatebox{90}{ARM}
    & \rotatebox{90}{Itanium}
    & \rotatebox{90}{MIPS}
    & \rotatebox{90}{POWER}
    & \rotatebox{90}{SPARC-TSO}
    & \rotatebox{90}{x86}
    & \rotatebox{90}{zSystems} \\
    \tabucline{2-}
    \firsthline
    Loads Reordered after Loads/Stores? & \cmark & \cmark & \cmark  & \cmark & \cmark & & & \\
    \hline
    Stores Reordered after Stores? & \cmark & \cmark & \cmark & \cmark & \cmark & & & \\
    \hline
    Stores Reordered after Loads? & \cmark & \cmark & \cmark & \cmark & \cmark & \cmark & \cmark & \cmark \\
    \hline
    Atomic Reordered with Loads/Stores? & \cmark & \cmark & & \cmark & \cmark & & & \\
    % \hline
    % Dependent Loads Reordered? & \cmark & & & & & & & \\
    % \hline
    % Dependent Stores Reordered & & & & & & & & \\
    \lasthline
  \end{tabu}
  % \caption{Summary of Memory Ordering}
  \caption{Memory Ordering on Different Architectures}
  \label{tbl:ordering}
\end{table}

% \subsection{Problem/Motivation}

% \todo[inline]{problems arising due to the unawareness of the specific architectures memory model}
% \todo[inline]{example for counter intuitive behaviour} \noindent

% Table \ref{tbl:ordering} shows that even the most restrictive architectures like the widely used x86

% The Intel-64 and AMD memory ordering models allow a load to be reordered with an earlier store to a different location, thus breaking sequential consistency.
% However, loads are not reordered with stores to the same location.
% Even though this is the only reordering allowed, it is sufficient for introducing counter intuitive behaviour, illustrated by the example given in Listing \ref{fig:intro:code}.

Table \ref{tbl:ordering} shows a rough overview of the memory ordering models used by different architectures.
Even the most restrictive -- like the x86 processors in our desktop computers -- allow loads to be reordered with an earlier store to a different location, thus breaking sequential consistency.
The reason for this particular reordering is a widely used optimization technique commonly referred to as \emph{store buffer}.

While caches improve subsequent repeated loads of a specific variable (after an initial cache miss),
% \todo[noline]{stolen from perfbook}
it is also necessary to accommodate frequent concurrent stores from multiple processors to a set of shared variables.
In cache-coherent systems, if the caches hold multiple copies of a given variable, all the copies of that variable must have the same value.
Each store will therefore invalidate all copies of the old value,
% resulting in even more cache misses
hence slowing down the computation % significantly.
\cite{ref:McKenney17}.

% \todo[inline]{reordering introduced by store buffers}

\begin{figure}[!h]
  \centering
  % styles
\tikzstyle{box} = [draw, text centered, rounded corners]
\tikzstyle{processor} = [box, fill=blue!10]
\tikzstyle{buffer} = [box, fill=red!10]
\tikzstyle{cache} = [box, fill=yellow!10]

\begin{tikzpicture}

  % \draw[step=1cm,gray,very thin] (-5, 0) grid (5, -10);

  % processors
  \path [processor] (-5, 0) rectangle (-1, -2);
  \path (-3, -1) node (p-0) [] {\textbf{Processor 0}};

  \path [processor] (1, 0) rectangle (5, -2);
  \path (3, -1) node (p-n) [] {\textbf{Processor n}};

  % store buffers
  \path [buffer] (-3, -3) rectangle (-1, -5);
  \path (-2, -3.7) node (s-0) [] {\textbf{Store}};
  \path (-2, -4.3) node (b-0) [] {\textbf{Buffer}};

  \path [buffer] (1, -3) rectangle (3, -5);
  \path (2, -3.7) node (s-n) [] {\textbf{Store}};
  \path (2, -4.3) node (b-n) [] {\textbf{Buffer}};

  % caches
  \path [cache] (-5, -6) rectangle (-1, -8);
  \path (-3, -7) node (c-0) [] {\textbf{Cache}};

  \path [cache] (1, -6) rectangle (5, -8);
  \path (3, -7) node (c-n) [] {\textbf{Cache}};

  % heap
  \path [box] (-5, -9) rectangle (5, -11);
  \path (0, -10) node (heap) [] {\textbf{Shared Memory}};

  % arrows
  \path [draw, <-] (-4, -2) -- (-4, -6);  % cache -> processor 0
  \path [draw, <->] (-2, -2) -- (-2, -3); % sb <-> processor 0
  \path [draw, ->] (-2, -5) -- (-2, -6);  % sb -> cache
  \path [draw, <->] (-3, -8) -- (-3, -9); % cache <-> memory

  \path [draw, <-] (4, -2) -- (4, -6);    % cache -> processor n
  \path [draw, <->] (2, -2) -- (2, -3);   % sb -> processor n
  \path [draw, ->] (2, -5) -- (2, -6);    % sb -> cache
  \path [draw, <->] (3, -8) -- (3, -9);   % cache <-> memory

  \path [draw, <-] (-1, -7) -- (-0.5, -7);
  \path (0, -7) node (dots) [] {\large $\ldots$};
  \path [draw, ->] (0.5, -7) -- (1, -7);

  % dots
  \path (0, -1) node (dots) [] {\large $\ldots$};

\end{tikzpicture}

  \caption{System Architecture with Store Buffers}
  \label{fig:intro:store-buffer-architecture}
\end{figure}

To remedy this situation, modern processors come equipped with store buffers, as shown in Figure \ref{fig:intro:store-buffer-architecture}.
When a variable is stored, the new value is placed in the processor's store buffer, which can proceed immediately without having to wait for the store to do something about all the old values of that variable residing in other processors' caches.
% store forwarding
In order to guarantee that each processor at least sees its own operations in program order, it first examines its store buffer before consulting the cache when loading a variable in a process called \emph{store forwarding}.
% store buffer flushes
% flushes and memory barrier operations (explicit and implicit - OS issue a memory barrier before switching contexts)
% arbitrarily delaying visibility of the store to other processors.
% Each store buffer is allowed to voluntarily flush its contents back to memory, or may be forced to do so by certain instructions like memory barriers.
Since each store buffer is allowed to voluntarily flush its contents back to memory, the visibility of the particular stores to other processors might be arbitrarily delayed, leading to the aforementioned memory misordering
illustrated by the store buffer litmus test shown in Listing \ref{lst:intro:store-buffer-litmus}.
% If the processor is not forced to flush its contents back to memory, visibility of stores to other processors is arbitrarily delayed, leading to the aforementioned misordering.
% This leads to the aforementioned misordering, as the particular stores' visibility to other processor might be arbitrarily delayed.
% This simple optimization is sufficient for introducing the aforementioned misordering, as the particular stores' visibility to other processor might be arbitrarily delayed.

\begin{lstlisting}[style=c++, numbers=left, numberstyle=\footnotesize, numberblanklines=false, caption={Store Buffer Litmus Test}, label={lst:intro:store-buffer-litmus}]
#include <assert.h>
#include <pthread.h>

#define ACCESS(x) (*(volatile typeof(x) *) &(x))
#define READ(x) ({typeof(x) TMP = ACCESS(x); TMP;})
#define WRITE(x,v) ({ACCESS(x) = (v);})

static int w0 = 0;
static int w1 = 0;

static int r0 = 0;
static int r1 = 0;

static void * P0 (void * p)
{
  WRITE(w0, 1);
  r0 = READ(w1);
  return p;
}

static void * P1 (void * p)
{
  WRITE(w1, 1);
  r1 = READ(w0);
  return p;
}

int main ()
{
  pthread_t t[2];
  pthread_create (t + 0, 0, P0, 0);
  pthread_create (t + 1, 0, P1, 0);
  pthread_join (t[0], 0);
  pthread_join (t[1], 0);
  assert(r0 + r1);
  return 0;
}
\end{lstlisting}

% \todo[inline]{x86 permits $\texttt{r0 = 0} \land \texttt{r1 = 0}$}
One might assume that the assertion on line 35 never triggers, because any possible interleaving should guarantee that at least one \lstinline[style=c++]{WRITE} must have happened before a subsequent \lstinline[style=c++]{READ} by the other thread in this symmetric example.
This intuition however breaks down with the addition of store buffers.
Consider the interleavings depicted in Figure \ref{fig:intro:store-buffer-litmus-interleaving}.
\begin{figure}[!h]
  \centering
  \begin{tikzpicture}

  \path (0, 0) node (init) [] {$w_0 = w_1 = 0$};
  \path (0, -6) node (join) [] {$r_0 + r_1 = 0$};

  \path (-3, -2) node (w0) [red] {$w_0 = 1$};
  \path [->] (init) edge node [above] {$t_0$} (w0);
  \path (-3, -4) node (r0) [] {$r_0 = w_1 = 0$};
  \path [->] (w0) edge node [left] {$t_0$} (r0);
  \path [->] (r0) edge node [below] {$t_0$} (join);

  \path (3, -2) node (w1) [red] {$w_1 = 1$};
  \path [->] (init) edge node [above] {$t_1$} (w1);
  \path (3, -4) node (r1) [] {$r_1 = w_0 = 0$};
  \path [->] (w1) edge node [right] {$t_1$} (r1);
  \path [->] (r1) edge node [below] {$t_1$} (join);

  \path (0, -2) node (buffered) [red] {buffered};
  \path [draw, dotted, red, <-] (w0) -- (buffered);
  \path [draw, dotted, red, ->] (buffered) -- (w1);

  \path [draw, dotted, ->] (init) -- (r0);
  \path [draw, dotted, ->] (init) -- (r1);

  % forwarding arrows

  % \path [draw, dotted, green, ->] (init) edge node [near end, left] {\cmark} (r0);
  % \path [draw, dotted, green, ->] (init) edge node [near end, right] {$\;$\cmark} (r1);

  % \path [dotted, red, ->] (w1) edge node [near end, below] {\xmark} (r0);
  % \path [dotted, red, ->] (w0) edge node [near end, below] {\xmark} (r1);

  % \path (0, -1.2) node (mem) [green] {memory};
  % \path (0, -3.6) node (sb) [red] {buffered};

  % alternative figure

  % \path (0, -10) node (Xinit) [draw, rounded corners, text centered, inner sep=0pt] {
    % \begin{tabu}{c}
      % \footnotesize memory \\
      % $w_0 = w_1 = r_0 = r_1 = 0$ \\
    % \end{tabu}
  % };
%
  % \path (-3, -13) node (Xw0) [draw, rounded corners, text centered, inner sep=0pt] {
    % \begin{tabu}{c}
      % \textbf{thread 0} \\
      % \hline
      % \footnotesize memory \\
      % $w_1 = r_0 = r_1 = 0$ \\
      % % \tabucline[on 1pt off 2pt]{-}
      % \hline
      % \footnotesize store buffer \\
      % % \tabucline[on 1pt off 2pt]{-}
      % $w_0 = 1$
    % \end{tabu}
  % };
%
  % \path (-3, -17) node (Xr0) [draw, rounded corners, text centered, inner sep=0pt] {
    % \begin{tabu}{c}
      % \textbf{thread 0} \\
      % \hline
      % \footnotesize memory \\
      % $w_1 = r_0 = r_1 = 0$ \\
      % % \tabucline[on 1pt off 2pt]{-}
      % \hline
      % \footnotesize store buffer \\
      % % \tabucline[on 1pt off 2pt]{-}
      % $w_0 = 1\;$
    % \end{tabu}
  % };
%
  % \path [->] (Xinit) edge node [left] {$w_0 = 1$} (Xw0);
  % \path [->] (Xw0) edge node [left] {$r_0 = w_1$} (Xr0);
\end{tikzpicture}

  \caption{Store Buffer Litmus Test Interleavings}
  \label{fig:intro:store-buffer-litmus-interleaving}
\end{figure}
After initializing the shared variables with zero, %($w_0 = w_1 = 0$)
% both
writes to $w_0$ and $w_1$ will be placed into the respective processor's store buffer.
If neither of the store buffers has been flushed prior to the subsequent reads $r_0 = w_1$ and $r_1 = w_0$, both processors must resort to the variable's initial value in their caches,
% which still contains the initial value,
hence triggering the assertion.
% Since it still contains their initial value,
% the threads will erroneously load zero,
% hence triggering the assertion.
% \todo[inline]{\texttt{experiments/demo/run.sh}: 177169 out of 1000000 failed = $\sim 17.7$ \% ($\sim 12$ min)}
% \todo[inline]{only 82 out of 1000000 returned the perfectly legal outcome of 2}
This willingness to reorder can be validated by running the store buffer litmus test given in Listing \ref{lst:intro:store-buffer-litmus}.
On our Intel i7-3770K this counter-intuitive reordering happened 177169 out of 1000000 test runs, while the perfectly legal outcome of $r_0 + r_1 = 2$ on the other hand only occurred 82 times.

% \todo[inline]{lack of concrete, formal specifications what programmers can rely on}
Programmers must therefore make careful use of memory barriers,
forcing the store buffer to be flushed back to memory after writing to a shared variable whenever necessary.
Unfortunately, many hardware vendors don't supply concrete formal specifications programmers can rely on.
For example, the memory ordering models of Intel's \cite{ref:Intel} as well as AMD's \cite{ref:AMD} x86 implementations are specified in an informal prose together with a set of litmus tests, having the potential for leading to widespread confusion.

% \subsection{Toolchain}

% \bigskip
\newpage
% \noindent

In this thesis we tried to address this issue by
% exploring encodings of an abstract machine model including memory ordering.
means of \emph{bounded model checking} \cite{ref:BMC}
based on the architectural view of a simple virtual machine model using store buffers,
reassambling the memory ordering habits of
the industry's leading
% current
x86 implementations.
% on the basis of a simple virtual machine model including store buffers.
% and explored encodings of an abstract machine model including memory ordering.
In order to evaluate our approach we developed
ConcuBinE\footnote{\url{https://github.com/phlo/concubine}} (short for \textbf{Concu}rrent \textbf{Bin}ary \textbf{E}valuator),
a tool for simulating and
% verifying
automated reasoning about
random memory access sequences by
concurrent programs.
% on a virtual machine using store buffers,
% dubbed ConcuBinE\footnote{\url{https://github.com/phlo/concubine}} (\textbf{Concu}rrent \textbf{Bin}ary \textbf{E}valuator).
It offers three modes of execution:
\begin{itemize}
  \item \texttt{simulate}:
    simulates the execution of given programs and returns the corresponding execution trace in an easy to read format.
  \item \texttt{solve}:
    automatically generates an encoding of given programs for a specific upper bound and evaluates the resulting bounded model checking problem with the help of state-of-the-art SMT solvers.
    If the problem is satisfiable, the corresponding execution trace (model) will be returned.
  \item \texttt{replay}:
    % used to validate a given trace via simulation.
    reevaluates a given trace via simulation.
    Used to validate the traces found by the \texttt{solve} mode.
\end{itemize}
Figure \ref{fig:intro:toolchain} shows a basic overview of its usage,
where red nodes symbolize in and output files,
% blue nodes the components according to the mode it is used
blue nodes the mode of opration
and yellow nodes external software (SMT solvers).

% In this thesis we developed a toolchain for simulating and bounded model checking of random memory access sequences by concurrent programs.

% \bigskip
% In this thesis we tried to address this issue by developing a toolchain for simulating and
% automated reasoning about
% % random memory access sequences of
% concurrent programs by means of \emph{bounded model checking} \cite{ref:BMC}.
% Instead of real hardware, it is based on the architectural view of a simple virtual machine model abstracting caches but

% \bigskip

% \todo[inline]{address the problem by the means of bounded model checking}
% \todo[inline]{ConcuBinE
% /ˈkɒŋ.kjə.baɪn/
% \textbf{Concu}rrent \textbf{Bin}ary \textbf{E}valuator - a toolchain for simulating and bounded model checking of random memory access sequences by concurrent programs on a simple virtual machine using store buffers.}
% \todo[inline]{components structured according to the mode they are used}

\begin{figure}[h]
  \centering
  % styles
\tikzstyle{module} = [draw, fill=blue!10, text centered, rounded corners, minimum height=1cm, minimum width=2.5cm]
\tikzstyle{file} = [module, fill=red!10]
\tikzstyle{external} = [module, fill=yellow!10]

\begin{tikzpicture}[auto]

  \node [file] (program) {program};

  \node [module] (simulate) [above right=0cm and 1cm of program] {\texttt{simulate}};
  \path [draw, ->] (program.east) -- ++(0.5, 0) |- (simulate);

  \node [module] (solve) [below right=0cm and 1cm of program] {\texttt{solve}};
  \path [draw, ->] (program.east) -- ++(0.5, 0) |- (solve);

  \node [external] (solver) [below= of solve] {SMT solver};
  \path [draw, ->, transform canvas={xshift=-0.5cm}] (solve) -- (solver);
  \path [draw, <-, transform canvas={xshift=0.5cm}] (solve) -- (solver);

  \node [file] (trace) [below right=0cm and 1cm of simulate] {trace};
  \path [draw, ->] (simulate.east) -- ++(0.5, 0) |- (trace);
  \path [draw, ->] (solve.east) -- ++(0.5, 0) |- (trace);

  \node [module, dashed] (replay) [right= of trace] {\texttt{replay}} edge [<-, dashed] (trace);
  \path [draw, dashed, ->] (replay) -- (replay.north |- simulate) -| (trace.north);

\end{tikzpicture}

  \caption{ConcuBinE's Various Modes of Operation}
  \label{fig:intro:toolchain}
\end{figure}


\newpage

% \todo[inline]{minimal model abstracting caches, preemption and store buffer flushes independent from execution} \noindent

\section{Machine Model}

% intro

% Our abstract machine model of a x86 multiprocessor system is illustrated in Figure \ref{fig:machine:overview}.

%We will start by defining an abstract machine model of a multiprocessor system that allows stores to be reordered after loads.
%Since we are only concerned with the machine's behaviour as observed by assembly programs, the internal structure of any real processor's microarchitecture is a highly abstracted.
%To keep the state space of the resulting model checking problems as small as possible, we defined our model on the basis of a 16 bit, 1 register machine.

% consistent with amd/intel litmus tests

% as observed by assembly programs.

% to reduce the state space, while emulating the behaviour of a x86 multiprocessor system.

% * abstract machine model to investigate the interaction of parallel programs through shared memory.

We will start by defining a \CHANGE{minimal virtual} machine model of a multiprocessor system as observed by assembly programs.
% \CHANGE{
% abstracting caches (assuming cache-coherence) and interleaving the execution of
% }
To keep the state space of the resulting model checking problems as small as possible, it is based on a 16 bit architecture, using only a minimal set of registers and a radically reduced instruction set.

\begin{figure}[!h]
  \centering
  % layers
\pgfdeclarelayer{background}
\pgfdeclarelayer{foreground}
\pgfsetlayers{background, main, foreground}

% styles
\tikzstyle{box} = [draw, text centered, rounded corners]
\tikzstyle{register} = [box, fill=white, minimum height=2em, minimum width=4em]
\tikzstyle{processor} = [draw, fill=blue!10, rounded corners]
\tikzstyle{buffer} = [processor, fill=red!10]%, dashed]

% distances
\def\blockdist{3}
\def\borderdist{0.2}

\begin{tikzpicture}

  % processor 0 %%%%%%%%%%%%%%%%%%%%%%%%%%%%%%%%%%%%%%%%%%%%%%%%%%%%%%%%%%%%%%%%
  \path (0, 0) node (title-0) [] {\textbf{Thread 0}};

  % store buffer
  \path (title-0)+(0, -1) node (sb-0) [] {\textbf{Store Buffer}};
  \path (sb-0)+(0, -1) node (full-0) [register] {\texttt{full}};
  \path (full-0)+(0, -1) node (adr-0) [register] {\texttt{adr}};
  \path (adr-0)+(0, -1) node (val-0) [register] {\texttt{val}};

  % registers
  \path (val-0)+(-1.1, -1.4) node (accu-0) [register] {\texttt{accu}};
  \path (val-0)+(1.1, -1.4) node (mem-0) [register, dashed] {\texttt{mem}};

  % paths
  \path [draw, <->] (val-0.west) -- (val-0.west -| accu-0.north) -- (accu-0.north);
  \path [draw, dashed, ->] (val-0.east) -- (val-0.east -| mem-0.north) -- (mem-0.north);

  \begin{pgfonlayer}{background}
    % bounding box
    \path (accu-0.west |- title-0.north)+(-\borderdist, \borderdist) node (a) {};
    \path (mem-0.south east)+(\borderdist, -\borderdist) node (b) {};
    \path [processor] (a) rectangle (b);

    % top left heap coordinate
    \path (a |- b)+(0, -1) node (tlh) {};

    % store buffer box
    \path (sb-0.north west)+(-\borderdist, \borderdist) node (a) {};
    \path (val-0.south -| sb-0.east)+(\borderdist, -\borderdist) node (b) {};
    \path [buffer] (a) rectangle (b);
  \end{pgfonlayer}

  % dots %%%%%%%%%%%%%%%%%%%%%%%%%%%%%%%%%%%%%%%%%%%%%%%%%%%%%%%%%%%%%%%%%%%%%%%
  \path (title-0)+(\blockdist, 0) node (dots) [] {\large $\ldots$};

  % processor n %%%%%%%%%%%%%%%%%%%%%%%%%%%%%%%%%%%%%%%%%%%%%%%%%%%%%%%%%%%%%%%%
  \path (dots)+(\blockdist, 0) node (title-n) [] {\textbf{Thread n}};

  % store buffer
  \path (title-n)+(0, -1) node (sb-n) [] {\textbf{Store Buffer}};
  \path (sb-n)+(0, -1) node (full-n) [register] {\texttt{full}};
  \path (full-n)+(0, -1) node (adr-n) [register] {\texttt{adr}};
  \path (adr-n)+(0, -1) node (val-n) [register] {\texttt{val}};

  % registers
  \path (val-n)+(-1.1, -1.4) node (accu-n) [register] {\texttt{accu}};
  \path (val-n)+(1.1, -1.4) node (mem-n) [register, dashed] {\texttt{mem}};

  % paths
  \path [draw, <->] (val-n.west) -- (val-n.west -| accu-n.north) -- (accu-n.north);
  \path [draw, dashed, ->] (val-n.east) -- (val-n.east -| mem-n.north) -- (mem-n.north);

  \begin{pgfonlayer}{background}
    % bounding box
    \path (accu-n.west |- title-n.north)+(-\borderdist, \borderdist) node (a) {};
    \path (mem-n.south east)+(\borderdist, -\borderdist) node (b) {};
    \path [processor] (a) rectangle (b);

    % bottom right heap coordinate
    \path (b)+(0, -\blockdist) node (brh) {};

    % store buffer box
    \path (sb-n.north west)+(-\borderdist, \borderdist) node (a) {};
    \path (val-n.south -| sb-n.east)+(\borderdist, -\borderdist) node (b) {};
    \path [buffer] (a) rectangle (b);
  \end{pgfonlayer}

  % heap %%%%%%%%%%%%%%%%%%%%%%%%%%%%%%%%%%%%%%%%%%%%%%%%%%%%%%%%%%%%%%%%%%%%%%%
  \path [box] (tlh) rectangle node (heap) [minimum height=2cm] {\textbf{Shared Memory}} (brh);

  % paths
  \path [draw, ->] (val-0) -- (heap.north -| val-0);
  \path [draw, <-, transform canvas={xshift=-3mm}] (accu-0) -- (heap.north -| accu-0);
  \path [draw, dashed, ->, transform canvas={xshift=3mm}] (accu-0) -- (heap.north -| accu-0);
  \path [draw, dashed, <-] (mem-0) -- (heap.north -| mem-0);

  \path [draw, ->] (val-n) -- (heap.north -| val-n);
  \path [draw, <-, transform canvas={xshift=-3mm}] (accu-n) -- (heap.north -| accu-n);
  \path [draw, dashed, ->, transform canvas={xshift=3mm}] (accu-n) -- (heap.north -| accu-n);
  \path [draw, dashed, <-] (mem-n) -- (heap.north -| mem-n);

\end{tikzpicture}

  \caption{Abstract Machine Model}
  \label{fig:machine:overview}
\end{figure}

A schematic overview is illustrated in Figure \ref{fig:machine:overview}, \CHANGE{where dashed lines depict components and data paths solely used by our \emph{compare and swap} mechanism}.
At the top of the figure are an arbitrary number of \CHANGE{logical} processors, each running a single \CHANGE{abstract} thread containing:

\begin{itemize}
  \item \ACCU: a single 16 bit accumulator register
  \item \MEM: a special purpose 16 bit register, storing the expected value required by a unary \emph{compare and swap} instruction
  % \item a \emph{store buffer} to break sequential consistency by delaying a single write, consisting of:
  \item a single element \emph{store buffer}, consisting of:
    \begin{itemize}
      \item \SBFULL: a one bit flag register, signaling that it contains a value and may be flushed
      \item \SBADR: a 16 bit address register
      \item \SBVAL: a 16 bit value register
    \end{itemize}
\end{itemize}

All \CHANGE{threads} are directly connected to the machine's shared memory, referred to as $\HEAP$, which will be uninitialized \CHANGE{unless any} eventual input data.
\CHANGE{Caches are abstracted as they don't influence
% our virtual machine's behaviour
the memory subsystem's basic behaviour
if they are coherent.}
In terms of memory ordering, the addition of a \emph{store buffer} allows stores to be reordered after loads, making our model consistent with Intel's or AMD's x86 memory ordering models \cite{ref:Intel, ref:AMD}.

% \bigskip
% \todo[inline]{model based on interleaving the execution of instructions and store buffer flushes}

\subsection*{Scheduling}

At each step \CHANGE{in time} a \CHANGE{single thread} either executes an instruction or flushes its store buffer back to memory.
Scheduling is generally performed non-deterministically under the following constraints.

\begin{enumerate}
  \item A \CHANGE{thread} can execute a read, modify or control operation at any time.
  \item A \CHANGE{thread} can voluntarily flush its store buffer to memory only if it is full.
  \item A \CHANGE{thread} can execute a write or barrier operation only if its store buffer is empty.
  % \item \CHANGE{A thread can only execute an instruction if it is not suspended} %(by a previous \lstASM{CHECK} statement).}
  % \item \CHANGE{A thread cannot execute any more instructions if it has been stopped.} % by a \lstASM{HALT}.}
\end{enumerate}

\CHANGE{
We argue that this interleaving semantic is sufficient for simulating the effects of real hardware, due to the execution of instructions and store buffer flushes being seen as independent.
A similar argument applies regarding preemption, assuming the presence of at least as many physical processors as running threads.
Since preemptive operating system kernels contain memory barriers when switching contexts, our thread centric view remains valid
% as the architectural state remains the same and the move to another physical processor just manifests as a store buffer flush.
% and the switch to another processor
and preemption
is just hidden behind a store buffer flush.
}

% \todo[inline]{checkpoints?}
% \todo[inline]{halted threads?}

\subsection*{Instructions}

Our machine uses a radically reduced instruction set that contains only the most substantial operations.
Instructions are stored separately for each thread and are therefore not contained in memory.
This abstraction allows the program counter to address instructions by their index, starting from zero.
To simplify the definition of operational semantics, instructions are labelled using the following attributes:

\begin{itemize}
  \item \textbf{modify} -- Modifies a register's content.
  \item \textbf{read} -- Reads from memory using \emph{store forwarding}: if $\SBFULL$ is set and $\SBADR$ equals the given target address, the value contained in $\SBVAL$ is read instead of the corresponding shared memory location.
  \item \textbf{write} -- Writes to the store buffer by setting $\SBFULL$ to true, $\SBADR$ to the given target address and $\SBVAL$ to the value contained in $\ACCU$.
  \item \textbf{barrier} -- Blocks execution if the store buffer is full ($\SBFULL$ is set). %memory barrier% - requires the store buffer to be flushed
  \item \textbf{atomic} -- Multiple operations performed as a single, \CHANGE{indivisible} instruction.% (implies barrier)
  \item \textbf{control} -- Modifies the order in which instructions are executed.
\end{itemize}

Due to the single register architecture, all instructions have at most one operand.
The following list shall give an overview of our machine's instruction set.

\newcommand{\defop}[3]{
  \paragraph{#1} \hfill #2
  \rule[0.5\baselineskip]{\textwidth}{0.1pt}\vspace{-0.5\baselineskip}\par\noindent
  #3
}

% \noindent
% \begin{tabu} to \textwidth {X[l]X[l]X[r]}
% \texttt{LOAD adr} & \texttt{accu = heap[adr]} & accu, read \\
% \texttt{STORE adr} & \texttt{heap[adr] = accu} & write
% \end{tabu}

\subsubsection*{Memory}

\CHANGE{
Two addressing modes are supported: direct and indirect denoted by square brackets (e.g. \lstASM{LOAD [arg]}).
}

\defop
{\lstASM{LOAD arg}}
{modify, read}
{Loads the value found at address \texttt{arg} into \ACCU.}

\defop
{\lstASM{STORE arg}}
{write}
{Stores the value found in $\ACCU$ at address \texttt{arg}.}

\defop
{\lstASM{FENCE}}
{barrier}
{Memory barrier.}

\subsubsection*{Arithmetic}

\CHANGE{
% \todo[inline]{say something about that all are unsigned and two's-compliment}
% Interpretation of a value's signedness is left to the programmmer.
% All values may be interpreted as unsigned or signed according to two's compliment.
Since the interpretation of a value's signedness is left to the programmer, we only included the most basic arithmetic instructions, which are equivalent for unsigned and two's compliment representations.
Division is therefore left out due to its different implementations for signed and unsigned integers.
% For example, division is left out on purpose, since it would require different
}

\defop
{\lstASM{ADD arg}}
{modify, read}
{Adds the value found at address \texttt{arg} to \ACCU.}

\defop
{\lstASM{ADDI arg}}
{modify}
{Adds the immediate value \texttt{arg} to \ACCU.}

\defop
{\lstASM{SUB arg}}
{modify, read}
{Subtracts the value found at address \texttt{arg} from \ACCU.}

\defop
{\lstASM{SUBI arg}}
{modify}
{Subtracts the immediate value \texttt{arg} from \ACCU.}

\defop
{\lstASM{MUL arg}}
{modify, read}
{Multiplies $\ACCU$ with the value found at address \texttt{arg}.}

\defop
{\lstASM{MULI arg}}
{modify}
{Multiplies $\ACCU$ with the immediate value \texttt{arg}.}

% ~\\\noindent
% \todo[inline]{\texttt{DIV} left because it's varying semantics in terms of signedness}
% Division instructions are omitted due to the varying semantics in terms of signedness.

\subsubsection*{Control Flow}

% \defop
% {\lstASM{CMP arg}}
% {modify, read}
% {Compares $\ACCU$ to the value found at address \texttt{arg} by performing an unsigned subtraction.}

\defop
{\lstASM{JMP arg}}
{control}
{Jumps to the statement at \texttt{arg} unconditionally.}

\defop
{\lstASM{JZ arg}}
{control}
{Jumps to the statement at \texttt{arg} if $\ACCU$ is zero.}

\defop
{\lstASM{JNZ arg}}
{control}
{Jumps to the statement at \texttt{arg} if $\ACCU$ is non-zero.}

\defop
{\lstASM{JS arg}}
{control}
{Jumps to the statement at \texttt{arg} if $\ACCU$ is negative according to two's compliment (most significant bit is set).}

\defop
{\lstASM{JNS arg}}
{control}
{Jumps to the statement at \texttt{arg} if $\ACCU$ is zero or positive according to two's compliment (most significant bit is unset).}

\defop
{\lstASM{JNZNS arg}}
{control}
{Jumps to the statement at \texttt{arg} if $\ACCU$ is positive according to two's compliment (non-zero and most significant bit is unset).}

\subsubsection*{Atomic}

% \todo[inline]{motivate}
% \noindent
\CHANGE{
We included a basic atomic \emph{compare and swap} operation, commonly used to implement synchronization primitives like semaphores and mutexes, as well as lockless data structures.
Since this particular command requires at least two operands (the target address and expected previous value), it had to be split up into two instructions due to our machine's unary input language.
% Due to the unary nature of our instruction format,
}

\defop
{\lstASM{MEM arg}}
{modify, read}
{Loads the value \CHANGE{from} address \texttt{arg} into $\ACCU$ and $\MEM$ as the expectation during a latter \emph{compare and swap} operation.}

\defop
{\lstASM{CAS arg}}
{modify, read, barrier, atomic}
{Atomically compares the expected value in $\MEM$ to the actual value found at address \texttt{arg} and only writes the value found in $\ACCU$ back to address \texttt{adr} if they are equal.
Acts like a memory barrier.
% \todo{explain why non-blocking}}
}

\subsubsection*{Termination}

% \todo[inline]{motivate}
\CHANGE{
% Explicit termination of a single thread or the whole machine can be issued by the following commands.
Explicitly stopping a single thread or terminating the whole machine can be achieved by the following commands.
}

\defop
{\lstASM{HALT}}
{barrier, control}
{Stops the current thread.}

\defop
{\lstASM{EXIT arg}}
{control}
{Stops the machine with exit code \texttt{arg}.}

\subsubsection*{Meta}

\CHANGE{
The following high-level meta instruction, mimicking a reverse semaphore (and therefore not available on real hardware),
shall simplify the implementation of so called \emph{checker threads}
% by mimicking common synchronization primitives
used to programmatically validate machine states at runtime.
% validate machine states at runtime.
}

\defop
{\lstASM{CHECK arg}}
{control}
{Synchronize on checkpoint \texttt{arg} \CHANGE{(rendezvous)}.
Suspends execution until all threads, containing a call to checkpoint \texttt{arg}, reached the corresponding \lstASM{CHECK} statement.}

% \newpage

\subsection*{Programs}

Each thread is programmed using an assembly style language, defined by the following syntax.

\begin{figure}[!h]
\begin{grammar}
\small

<int> ::= an integer number

<label> ::= a sequence of printable characters without "#"

<string> ::= a sequence of whitespace and printable characters without "#" and "\\n"

<comment> ::= "#" <string>

<nullary> ::= "FENCE" | "HALT"

<unary> ::= "ADDI" | "SUBI" | "MULI" | "EXIT" | "CHECK"

<memory> ::= "LOAD" | "STORE" | "ADD" | "SUB" | "MUL" | "CMP" | "MEM" | "CAS"

<jump> ::= "JMP" | "JZ" | "JNZ" | "JS" | "JNS" | "JNZNS"

<instruction> ::= <nullary>
\alt <unary> <int>
\alt <memory> ( <int> | "["<int>"]" )
\alt <jump> ( <int> | <label> )

<statement> ::= <label>":" <instruction> | <instruction>

<line> ::= <statement> | <statement> <comment> | <comment>

<program> ::= <line> | <line> "\\n" <program>
\end{grammar}
\caption{Program Syntax}
\label{fig:syntax:program}
\end{figure}

% If the final statement in a given program is no control instruction, a final \texttt{HALT} is inserted implicitly.
If the final statement in a given program is not an \lstASM{EXIT} instruction or unconditional \lstASM{JMP}, \CHANGE{an additional} \lstASM{HALT} is inserted implicitly.


\newpage

\section{Execution}

% \begin{verbatim}
% usage: concubine simulate [options] <program> ...
%
% options:
  % -c         run until an exit code > 0 is encountered
  % -k bound   execute a specific number of steps
  % -m mmap    read initial heap contents from file
  % -o name    output file name (default: sim.{trace,mmap})
  % -s seed    random number generator's seed
  % program    one ore more source files, each being executed as a separate thread
% \end{verbatim}

Execution can be simulated using the \texttt{simulate} submodule, implementing a virtual machine for the previously defined model.
It takes an arbitrary number of program files, each being run on a separate thread and produces an execution trace using a pseudo random number generator to determine the schedule.

\begin{minipage}{.45\textwidth}
  \lstinputlisting[language={[concubine]Assembler}, caption=\texttt{thread.0.asm}, label={lst:program:thread.0.asm}, xleftmargin=0.39\textwidth]{../../examples/demo/processor.0.asm}
\end{minipage}
\begin{minipage}{.45\textwidth}
  \lstinputlisting[language={[concubine]Assembler}, caption=\texttt{thread.1.asm}, label={lst:program:thread.1.asm}, xleftmargin=0.39\textwidth]{../../examples/demo/processor.1.asm}
\end{minipage}

An implementation of the store buffer litmus test for our machine is given in Listings \ref{lst:program:thread.0.asm} and \ref{lst:program:thread.1.asm}.
The following command line can now be used to simulate it's execution.

\begin{verbatim}
$ concubine simulate thread.0.asm thread.1.asm
\end{verbatim}

After the simulation has been finished, the execution trace will be saved to the current directory as \texttt{sim.trace} per default.
% To change the path and output file names, the \texttt{-o} command line parameter can be used.

\subsection{Traces}

The execution traces are simple text files, capturing the basic environment and machine state transitions.
A trace of the previous simulation could look as follows.

\begin{lstlisting}[language={[concubine]Assembler}, caption={Simple Output Trace}, label={lst:trace:simple}, xleftmargin=\parindent]
thread.0.asm
thread.1.asm
.
# tid pc  cmd   arg accu  mem adr val full  heap
0     0   ADDI  1   0     0   0   0   0     {}      # 0
1     0   ADDI  1   0     0   0   0   0     {}      # 1
0     1   STORE 0   1     0   0   0   0     {}      # 2
1     1   STORE 1   1     0   0   0   0     {}      # 3
0     2   FLUSH -   1     0   0   1   1     {}      # 4
1     2   FLUSH -   1     0   1   1   1     {(0,1)} # 5
0     2   LOAD  1   1     0   0   1   0     {(1,1)} # 6
1     2   LOAD  0   1     0   1   1   0     {}      # 7
0     3   CHECK 0   1     0   0   1   0     {}      # 8
1     3   CHECK 0   1     0   1   1   0     {}      # 9
0     4   HALT  -   1     0   0   1   0     {}      # 10
1     4   HALT  -   1     0   1   1   0     {}      # 11
\end{lstlisting}

At the beginning is a list of paths to the programs involved in creating the specific trace, followed by a delimiter.
The rest of the file contains the actual execution steps, one per line, starting with the scheduled thread's identifier, followed by details about the state prior to the current statement's execution, consisting of:

\begin{itemize}
  \item program counter of the instruction about to be executed
  \item instruction symbol, or \lstinline[language={[concubine]Assembler}]{FLUSH} if a thread is flushing its store buffer
  \item instruction argument, or \texttt{-} for unary instructions
  \item accumulator register
  \item \lstinline[language={[concubine]Assembler}]{CAS} memory register
  \item store buffer address register
  \item store buffer value register
  \item store buffer full flag
  \item memory cell updated in the previous step, represented as a single pair in set notation \texttt{\string{(index,value)\string}}, or an empty set \texttt{\string{\string}} if the memory state didn't change
\end{itemize}

The result of the instruction's execution will be visible, the next time the specific thread is scheduled.
A better visualization of a particular thread's state transitions can be achieved by only considering e.g. thread 0 in the trace given in Listing \ref{lst:trace:simple}:

\begin{lstlisting}[language={[concubine]Assembler}, xleftmargin=\parindent]
# tid pc  cmd   arg accu  mem adr val full  heap
0     0   ADDI  1   0     0   0   0   0     {}      # 0
0     1   STORE 0   1     0   0   0   0     {}      # 2
0     2   FLUSH -   1     0   0   1   1     {}      # 4
0     2   LOAD  1   1     0   0   1   0     {(1,1)} # 6
0     3   CHECK 0   1     0   0   1   0     {}      # 8
0     4   HALT  -   1     0   0   1   0     {}      # 10
\end{lstlisting}

A more formal definition of the trace file syntax is given below.

\begin{figure}[h]
\begin{grammar}
\small

<int> ::= an integer number

<path> ::= a unix file path

<label> ::= a sequence of printable characters without "#"

<string> ::= a sequence of whitespace and printable characters without "#" and "\\n"

<comment> ::= "#" <string>

<program> ::= <path> | <path> <comment> | <comment>

<programs> ::= <program> | <program> "\\n" <programs>

<mmap> ::= <path> | <path> <comment>

<header> ::= <programs> "\\n" "." | <programs> "\\n" "." <mmap>

<pc> ::= <int> | <label>

<op> ::= "FLUSH"
\alt "LOAD" | "STORE" | "FENCE"
\alt "ADD" | "ADDI" | "SUB" | "SUBI" | "MUL" | "MULI"
\alt "CMP" | "JMP" | "JZ" | "JNZ" | "JS" | "JNS" | "JNZNS"
\alt "MEM" | "CAS"
\alt "HALT" | "EXIT" | "CHECK"

<arg> ::= <int> | "["<int>"]" | <label> | "-"

<bool> ::= "0" | "1"

<heap> ::= "{}" | "{("<adr>","<val>")}"

<state> ::= <int> <pc> <op> <arg> <int> <int> <int> <int> <bool> <heap>

<step> ::= <state> | <state> <comment> | <comment>

<steps> ::= <state> | <state> "\\n" <steps>

<trace> ::= <header> "\\n" <steps>
\end{grammar}
\caption{Trace Syntax}
\label{fig:syntax:trace}
\end{figure}

\subsection{Memory Maps}

Access to uninitialized memory cells is captured in so called memory maps, again a simple text file format, containing one address value pair per line, delimited by a whitespace.

\begin{lstlisting}[language={[concubine]Assembler}, caption={Output Trace Accessing Uninitialized Memory}, label={lst:trace:uninitialized}, xleftmargin=\parindent]
thread.0.asm
thread.1.asm
. sim.mmap
# tid pc  cmd   arg accu  mem adr val full  heap
0     0   ADDI  1   0     0   0   0   0     {}      # 0
1     0   ADDI  1   0     0   0   0   0     {}      # 1
0     1   STORE 0   1     0   0   0   0     {}      # 2
0     2   LOAD  1   1     0   0   1   1     {}      # 3
0     3   CHECK 0   57647 0   0   1   1     {}      # 4
1     1   STORE 1   1     0   0   0   0     {}      # 5
1     2   LOAD  0   1     0   1   1   1     {}      # 6
1     3   CHECK 0   34446 0   1   1   1     {}      # 7
0     4   FLUSH -   57647 0   0   1   1     {}      # 8
1     4   FLUSH -   34446 0   1   1   1     {(0,1)} # 9
0     4   HALT  -   57647 0   0   1   0     {(1,1)} # 10
1     4   HALT  -   34446 0   1   1   0     {}      # 11
\end{lstlisting}

Consider another random simulation, given in Listing \ref{lst:trace:uninitialized}.
In this simulation, both threads are reading uninitialized memory locations at addresses 0 and 1 respectively.
To ensure reproducibility, a memory map file \texttt{sim.mmap} is stored together with the trace file \texttt{sim.trace} per default and is referenced just after the delimiter at the end of the program list in the trace file's header.

\begin{lstlisting}[language={[concubine]Assembler}, caption={Output Memory Map of the Trace in Listing \ref{lst:trace:uninitialized}}, label={lst:mmap:uninitialized}, xleftmargin=0.39\textwidth]
0 34446
1 57647
\end{lstlisting}

Memory maps can also be used to initialize shared memory with input data by supplying the \texttt{-m} command line parameter, followed by a path to an existing memory map file.
For completeness, the syntax of such trace files is given below.

\begin{figure}[h]
\begin{grammar}
\small

<int> ::= an integer number

<string> ::= a sequence of whitespace and printable characters without "#" and "\\n"

<comment> ::= "#" <string>

<cell> ::= <int> <int> | <int> <int> <comment> | <comment>

<mmap> ::= <cell> | <cell> "\\n" <mmap>
\end{grammar}
\caption{Memory Map Syntax}
\label{fig:syntax:mmap}
\end{figure}

% \subsection{Finding Problematic Traces}
% \subsection{Searching Problematic Traces Exhaustively}
\subsection{Finding Specific Machine States}

It also possible to search for specific machine states via simulation.
The command line switch \texttt{-c} simulates the given programs exhaustively until a trace yielding an exit code greater than zero is encountered.

\lstinputlisting[language={[concubine]Assembler}, caption=\texttt{checker.asm}, label={lst:program:checker.asm}, xleftmargin=0.39\textwidth]{../../examples/demo/checker.asm}

To find traces exhibiting a problematic reordering of writes in our store buffer litmus test example, we make use of a so called \emph{checker thread}, given in Listing \ref{lst:program:checker.asm}, that examines the machine state at a specific point in time by using the \lstinline[language={[concubine]Assembler}]{CHECK} instruction and stops execution with an \lstinline[language={[concubine]Assembler}]{EXIT 1} if none of the store buffers has been flushed, leading to both memory locations still being zero.

\lstinputlisting[language={[concubine]Assembler}, caption=\texttt{init.mmap}, label={lst:mmap:init.mmap}, xleftmargin=0.39\textwidth]{../../examples/demo/init.mmap}

This requires that the relevant memory cells are set to zero in order to prevent uninitialized reads from influencing the checking procedure and is achieved by using the memory map given in Listing \ref{lst:mmap:init.mmap}.
The search can now be started by issuing the following command line:

\begin{verbatim}
$ concubine simulate -c -m init.mmap \
                     checker.asm \
                     thread.0.asm \
                     thread.1.asm
\end{verbatim}

After a couple of simulations, a trace like the one shown in Listing \ref{lst:trace:error}, where none of the store buffers has been flushed prior to the \emph{checker threads} \lstinline[language={[concubine]Assembler}]{ADD} instructions, will be found.

\begin{lstlisting}[language={[concubine]Assembler}, caption={Sequentially Inconsistent Trace}, label={lst:trace:error}, xleftmargin=\parindent]
checker.asm
thread.0.asm
thread.1.asm
. sim.mmap
# tid pc    cmd   arg   accu  mem adr val full  heap
1     0     ADDI  1     0     0   0   0   0     {}  # 0
2     0     ADDI  1     0     0   0   0   0     {}  # 1
1     1     STORE 0     1     0   0   0   0     {}  # 2
2     1     STORE 1     1     0   0   0   0     {}  # 3
1     2     LOAD  1     1     0   0   1   1     {}  # 4
2     2     LOAD  0     1     0   1   1   1     {}  # 5
1     3     CHECK 0     0     0   0   1   1     {}  # 6
2     3     CHECK 0     0     0   1   1   1     {}  # 7
0     0     CHECK 0     0     0   0   0   0     {}  # 8
0     1     ADD   0     0     0   0   0   0     {}  # 9
0     2     ADD   1     0     0   0   0   0     {}  # 10
0     3     JZ    error 0     0   0   0   0     {}  # 11
0     error EXIT  1     0     0   0   0   0     {}  # 12
\end{lstlisting}


\newpage

\section{Verification}

\todo{intro
  \begin{itemize}
    \item submodule \texttt{solve}
    \item verification of assembly programs by bounded model checking based on register states
    \item using satisfiability modulo theories (arrays and bitvectors)
    \item problem defined by state and transition variables
  \end{itemize}
}

\todo{states
  \begin{itemize}
    \item $\texttt{accu}_t^k$ - accumulator of thread $t$
    \item mem
    \item sb-adr
    \item sb-val
    \item sb-full
    \item stmt
    \item block
    \item halt
    \item heap
    \item exit-flag
    \item exit-code
  \end{itemize}
}

\paragraph{$\accu^k_t$} accumulator register state

The following variables are used to encode the machine's state at a particular step $k \in [0, bound] \subseteq \mathbb{N}$, where $k = 0$ is the machine's initial state.

\cite{btor2}

Let $\bvsort[n]$ be the fixed size bitvector sort of width $n$ and $\mathcal{A}$ the array sort with index and element sorts $\bvsort$.

\renewcommand{\arraystretch}{1.25}
\begin{table}[h!]
\begin{tabular}{lll}
  \hline
  \textbf{global} &&\\
  \hline
  $\texttt{heap}^k$ & $\mathcal{A}$ & shared memory state in step $k$ \\
  $\texttt{exit}^k$ & $\bvsort[1]$ & exit flag in step $k$ \\
  $\texttt{exit-code}^k$ & $\bvsort$ & exit code in step $k$ \\
  \hline
  \textbf{thread local} &&\\
  \hline
  $\accu^k_t$ & $\bvsort$ & accumulator register state of thread $t$ in step $k$ \\
  $\texttt{mem}^k_t$ & $\bvsort$ & CAS memory register state of thread $t$ in step $k$ \\
  $\texttt{sb-adr}^k_t$ & $\bvsort$ & store buffer address register of thread $t$ in step $k$ \\
  $\texttt{sb-val}^k_t$ & $\bvsort$ & store buffer value register of thread $t$ in step $k$ \\
  $\texttt{sb-full}^k_t$ & $\bvsort[1]$ & store buffer full flag of thread $t$ in step $k$ \\
  $\texttt{stmt}^k_{t, pc}$ & $\bvsort[1]$ & statement activation for $pc$ of thread $t$ in step $k$ \\
  $\texttt{block}^k_{t, id}$ & $\bvsort[1]$ & block flag for checkpoint $id$ of thread $t$ in step $k$ \\
  $\texttt{halt}^k_t$ & $\bvsort[1]$ & halt flag of thread $t$ in step $k$ \\
  \hline
\end{tabular}
\caption{State Variables}
\end{table}

\todo{transitions
  \begin{itemize}
    \item thread
    \item flush
    \item check
  \end{itemize}
}

\todo{memory access
  \begin{itemize}
    \item $load$ predicate
  \end{itemize}
}

Let $ite: \bvsort[1] \times \bvsort[n] \times \bvsort[n] \to \bvsort[n]$ be a functional if-then-else, returning the value $a \in \bvsort[n]$ if $x \in \bvsort[1]$ is \emph{true}, else the value $b \in \bvsort[n]$.
\[
  ite(x, a, b) =
  \begin{cases}
    a \text{ if } x \text{ is } true \\
    b \text{ otherwise}
  \end{cases}
\]

Let $read^k: \bvsort \to \bvsort$ be the function returning the element $e \in \bvsort$ with index $i \in \bvsort$ from array $a \in \asort$.

Let $load^k_t: \asort \times \bvsort \times \bvsort[1] \to \bvsort$ be the function for loading a particular shared memory element $e \in \bvsort$ with index $i \in \bvsort$ from array $a \in \asort$ and the indirection flag $indirect \in \bvsort[1]$, defined as follows.

\begin{align}
  load(a, i, indirect) = ite( & indirect, \\
                              & ite(\sbfull^k, \\
                              & \;\;\;\;\;\, ite(\sbadr^k = i, \\
                              & \;\;\;\;\;\, \;\;\;\;\;\, ite(\sbval^k_t = i, \\
                              & \;\;\;\;\;\, \;\;\;\;\;\, \;\;\;\;\;\, \sbval^k_t, \\
                              & \;\;\;\;\;\, \;\;\;\;\;\, \;\;\;\;\;\, read^k(i)), \\
                              & \;\;\;\;\;\, \;\;\;\;\;\, ite(\sbadr^k_t = read^k(i), \\
                              & \;\;\;\;\;\, \;\;\;\;\;\, \;\;\;\;\;\, \sbval^k_t, \\
                              & \;\;\;\;\;\, \;\;\;\;\;\, \;\;\;\;\;\, read(read^k(i)))), \\
                              & \;\;\;\;\;\, read^k(read^k(i))), \\
                              & ite(\sbfull^k_t \land \sbadr^k_t = i, \\
                              & \;\;\;\;\;\, \sbval, \\
                              & \;\;\;\;\;\, read^k(a, i)))
\end{align}

\paragraph{Store Forwarding:} $\sbfull \; \land \; \sbadr = adr \land \accu = \sbval \lor \lnot (\sbfull \land \sbadr = adr)$

\todo{frame axioms
  \begin{itemize}
    \item instructions
    \item \texttt{ADDI val}\\ $\texttt{exec}^k_{t, pc} \land \accu^{k + 1}_t = \accu^k_t + \texttt{val} \land \mem^{k + 1}_t = \mem^k_t$
  \end{itemize}
}

\paragraph{\texttt{LOAD adr}}

\todo{constraints}


\newpage

\section{{\SMTLIB} Encoding}

% We will now outline the actual encoding process for \SMTLIB.
% We will now outline the process of generating \SMTLIB formulas for our model checking problems, following the previously defined encoding scheme.
% We will now outline a simplistic {\CPP} implementation for generating {\SMTLIB} \cite{ref:SMT-Lib} formulas of our model checking problems, following the previously defined encoding scheme and use the store buffer litmus test as an exemplary input for demonstration.
We will now outline a simplistic {\CPP} implementation for \CHANGE{encoding our bounded model checking problems in the {\SMTLIB} \cite{ref:SMT-Lib} format according to} the previously defined scheme and use the store buffer litmus test as an exemplary input for demonstration.
% We will now outline a simplistic {\CPP} implementation for generating our model checking problems in the \SMTLIB format, following the previously defined encoding scheme.
% Our store buffer litmus test is used as an example to demonstrate the output
\CHANGE{
% ConCuBinE's actual source code is available at \url{https://github.com/phlo/concubine}.
%
{\SMTLIB} is the result of an ongoing international initiative with the goal of developing a rigorously standardized language for describing SMT theories as well as formulas and is supported by nearly all SMT solvers. % (we used Boolector, Z3 and CVC4).
Based on a Lisp like syntax, it covers a wide variety of theories like the theory of integers, fixed size bit-vectors, arrays and many more.
In order to apply specialized and more efficient satisfiability techniques, theories are grouped into so called logics and every formula must explicitly specify the logic it is based on.
See \url{www.smt-lib.org} for a full list of theories and logics available.
% We used \mbox{\texttt{QF_AUFBV}}, allowing closed quantifier-free formulas over the theory of bit-vectors and bit-vector arrays extended with free sort and function symbols.
In our case \mbox{\texttt{QF_AUFBV}} was used, allowing closed quantifier-free formulas over the theory of bit-vectors and bit-vector arrays extended with free sort and function symbols.

% Even though incremental solving would have been the best choice for our bounded model checking problem, it just started to be supported by the main solvers at the time we started development and we therefore chose to manually unroll our problem for every step $k \leq \BOUND$.
% Even though {\SMTLIB} also offers \lstSMTLIB{push} and \lstSMTLIB{pop} commands for incremental solving and would have been the best choice to keep the size of our generated formulas as small as possible.
{\SMTLIB} also offers incremental solving through \lstSMTLIB{push} and \lstSMTLIB{pop} commands.
Because of our encoding's iterative nature, it would have been the best choice in order to keep the size of our formulas as small as possible.
Unfortunately, incremental solving just started to be supported by the main SMT solvers at the time we began development and therefore chose to manually unroll our bounded model checking problem for every step $k \leq \BOUND$.
% Unfortunately, support for incremental solving by the main SMT solvers just started at the time we began development and therefore chose to manually unroll our bounded model checking problem for every step $k \leq \BOUND$.
}

\subsection*{Types}
% \todo[inline]{types} %%%%%%%%%%%%%%%%%%%%%%%%%%%%%%%%%%%%%%%%%%%%%%%%%%%%%%%%%%%

To differentiate among particular machine states, we start by introducing:% an \lstCPP{enum} \lstCPP{State}.
% A \lstCPP{enum} \lstCPP{State} is used to differentiate among machine states

% \begin{lstlisting}[style=c++]
% typedef uint16_t word_t;
% \end{lstlisting}

\begin{lstlisting}[style=c++]
enum State
{
  heap,
  accu,
  mem,
  adr,
  val,
  full,
  stmt,
  block,
  halt,
  exit,
  exit_code
};
\end{lstlisting}
% All possible values contained in our machine's 16 bit registers are captured by the corresponding unsigned integer type \lstCPP{word_t}.
% \begin{lstlisting}[style=c++]
% using uint = unsigned;
% \end{lstlisting}
% Each available instruction is represented by it's own type, derived from an abstract class \lstCPP{Instruction}, capturing arguments and abolishing the need for lengthy case splits to determine the actual instruction at runtime by utilizing dynamic dispatch of member functions.
Each available instruction is represented by it's own type, derived from an abstract class \lstCPP{Instruction}, capturing arguments and abolishing the need for lengthy case splits by utilizing dynamic dispatch of virtual member functions for taking the appropriate action at runtime.%determine the type of instruction at runtime.

\begin{lstlisting}[style=c++]
struct Instruction
{
  uint arg;
  bool indirect;

  Instruction (uint a, bool i = false) : arg(a), indirect(i) {};

  virtual std::string encode (uint k, uint t, State state) = 0;
};
\end{lstlisting}
% virtual std::string encode (uint k, uint t, uint pc) = 0;

\subsection*{Globals}
% \todo[inline]{gobal variables} %%%%%%%%%%%%%%%%%%%%%%%%%%%%%%%%%%%%%%%%%%%%%%%%%

% Inputs and utility variables are given as globals to keep things simple and function signatures as small as possible.
Commonly used variables are given as globals to keep things simple and function signatures as small as possible.

\begin{table}[!h]
\small
\centering
% \begin{center}
\begin{tabu}{|c|c|c|}
  \firsthline
  % \begin{tabu}{c} \textbf{Thread 0} \end{tabu} &
  % \begin{tabu}{c} \textbf{Thread 1} \end{tabu} &
  % \begin{tabu}{c} \textbf{Thread 2} \end{tabu} \\
  \textbf{Thread 0} & \textbf{Thread 1} & \textbf{Thread 2} \\
  \hline
  \hline
  \begin{tabu}{ll}
    \lstASM{ADDI 1}   & $\STMT^k_{0, 0}$ \\
    \lstASM{STORE 0}  & $\STMT^k_{0, 1}$ \\
    \lstASM{LOAD 1}   & $\STMT^k_{0, 2}$ \\
    \lstASM{CHECK 0}  & $\STMT^k_{0, 3}$ \\
    \lstASM{HALT}     & $\STMT^k_{0, 4}$ \\
  \end{tabu}
  &
  \begin{tabu}{ll}
    \lstASM{ADDI 1}   & $\STMT^k_{1, 0}$ \\
    \lstASM{STORE 1}  & $\STMT^k_{1, 1}$ \\
    \lstASM{LOAD 0}   & $\STMT^k_{1, 2}$ \\
    \lstASM{CHECK 0}  & $\STMT^k_{1, 3}$ \\
    \lstASM{HALT}     & $\STMT^k_{1, 4}$ \\
  \end{tabu}
  &
  \begin{tabu}{ll}
    & \\
    & \\
    & \\
    & \\
    \lstASM{CHECK 0}        & $\STMT^k_{2, 0}$ \\
    \lstASM{ADD 0}          & $\STMT^k_{2, 1}$ \\
    \lstASM{ADD 1}          & $\STMT^k_{2, 2}$ \\
    \lstASM{JZ error}       & $\STMT^k_{2, 3}$ \\
    \lstASM{EXIT 0}         & $\STMT^k_{2, 4}$ \\
    \lstASM{error: EXIT 1}  & $\STMT^k_{2, 5}$ \\
  \end{tabu} \\
  \lasthline
\end{tabu}
\caption{Store buffer litmus test programs and their activation variables.}
\label{tbl:encoding:programs}
% \end{center}
\end{table}
% \vspace{-\baselineskip}
\noindent
Input programs are expressed as lists of instruction \CHANGE{pointers} and initialized according to our store buffer litmus test, given in Table \ref{tbl:encoding:programs}.

\newpage

% \begin{lstlisting}[style=c++]
% std::vector<std::vector<Instruction *>> programs;
% \end{lstlisting}
% \vspace{-\baselineskip}
% \begin{tabu} {XXX}
  % \texttt{\{\{} & & \\
  % \lstCPP{new Addi(1)}  & \lstCPP{new Addi(1)}  & \lstCPP{new Check(0)} \\
  % \lstCPP{new Store(0)} & \lstCPP{new Store(1)} & \lstCPP{new Add(0)} \\
  % \lstCPP{new Load(1)}  & \lstCPP{new Load(0)}  & \lstCPP{new Add(1)} \\
  % \lstCPP{new Check(0)} & \lstCPP{new Check(0)} & \lstCPP{new Jz(5)} \\
  % \texttt{\},} && \lstCPP{new Exit(0)} \\
  % && \lstCPP{new Exit(1)} \\
% \end{tabu}
% \vspace{.5\baselineskip}
  % {               {               {
    % new Store(0),   new Store(1),   new Add(0),
    % new Load(1),    new Load(0),    new Add(1),
    % new Check(0),   new Check(0),   new Jz(5),
    % new Halt()      new Halt()      new Exit(0),
  % },              },                new Exit(1)
                                  % }};
\begin{lstlisting}[style=c++]
std::vector<std::vector<Instruction *>> programs = {
  {
    new Addi(1),
    new Store(0),
    new Load(1),
    new Check(0),
    new Halt()
  },
  {
    new Addi(1),
    new Store(1),
    new Load(0),
    new Check(0),
    new Halt()
  },
  {
    new Check(0),
    new Add(0),
    new Add(1),
    new Jz(5),
    new Exit(0),
    new Exit(1)
  }
};
\end{lstlisting}
The initial memory layout in our example, given by the memory map in Listing \ref{lst:mmap:init.mmap}, is captured by the following variable.
\begin{lstlisting}[style=c++]
std::map<uint, uint> mmap = {{0, 0}, {1, 0}};
\end{lstlisting}
Last but not least, we define the upper bound \CHANGE{$\BOUND$} of the resulting model checking problem as the final input variable.
\begin{lstlisting}[style=c++]
uint bound = 17;
\end{lstlisting}
% To increase readability by the use of compound statements, we use a string stream to store the generated formula.
The resulting formula is stored in a string stream to increase readability by the use of compound statements.
\begin{lstlisting}[style=c++]
std::ostringstream formula;
\end{lstlisting}
% The following utility variables are generated during a preprocessing step by analyzing the input programs and used to simplify the encoding process.
To simplify the encoding process, \CHANGE{the} following utility variables are generated during a preprocessing step by analyzing the input programs.
\CHANGE{
  % Although they might look a bit complex, we included their values after our store buffer litmus test example, since they are the basis of
  % Although they might look a bit complex, we omit further details on how they are created as they just capture frequently used properties and initialize them according to our store buffer litmus test example in order to enable the reader to reproduce the output of encoding routines.
Although they might seem a bit complex, we omit further details on how they are created as they just identify certain types of instructions and initialize them according to our store buffer litmus test example for completeness.
% to prevent repeated iteration of input programs required for their detection.
% and initialize them according to our store buffer litmus test example.
}
% % by reducing the number of tests.
% \\~\\
% \noindent
% % Program counters of instructions, modifying a particular state are collected in a map to prevent repeated iteration of input programs for their detection.
% % A map, containing the program counters of instructions modifying a particular state.
% We use a map to collect the program counters of instructions, modifying a particular state, in order to reduce the number of tests inside our encoding loops.
\begin{itemize}
\item Program counters of instructions modifying a particular state.
% \vspace{-.5\baselineskip}
% \item \CHANGE{A nested mapping, containing program counters of instructions modifying a particular state for every thread.}
\begin{lstlisting}[style=c++]
std::map<State, std::vector<std::vector<uint>>> updates = {
  {State::accu, {{0, 2},    // `\lstCMT{$\ACCU_0 \leftarrow$ ADDI 1, LOAD 0}`
                 {0, 2},    // `\lstCMT{$\ACCU_1 \leftarrow$ ADDI 1, LOAD 1}`
                 {1, 2}}},  // `\lstCMT{$\ACCU_2 \leftarrow$ ADD 0, ADD 1}`
  {State::adr,  {{1},       // `\lstCMT{$\SBADR_0 \leftarrow$ STORE 0}`
                 {1},       // `\lstCMT{$\SBADR_1 \leftarrow$ STORE 1}`
                 {}}},      // `\lstCMT{$\SBADR_2 \leftarrow$ $\varnothing$}`
  {State::val,  {{1},       // `\lstCMT{$\SBVAL_0 \leftarrow$ STORE 0}`
                 {1},       // `\lstCMT{$\SBVAL_1 \leftarrow$ STORE 1}`
                 {}}},      // `\lstCMT{$\SBVAL_2 \leftarrow$ $\varnothing$}`
  {State::full, {{1},       // `\lstCMT{$\SBFULL_0 \leftarrow$ STORE 0}`
                 {1},       // `\lstCMT{$\SBFULL_1 \leftarrow$ STORE 1}`
                 {}}},      // `\lstCMT{$\SBFULL_2 \leftarrow$ $\varnothing$}`
  {State::halt, {{4},       // `\lstCMT{$\HALT_0 \leftarrow$ HALT}`
                 {4},       // `\lstCMT{$\HALT_1 \leftarrow$ HALT}`
                 {}}},      // `\lstCMT{$\HALT_2 \leftarrow$ $\varnothing$}`
  {State::exit, {{},        // `\lstCMT{$\EXIT \leftarrow$ $\varnothing$}`
                 {},        // `\lstCMT{$\EXIT \leftarrow$ $\varnothing$}`
                 {4, 5}}}}; // `\lstCMT{$\EXIT \leftarrow$ EXIT 0, EXIT 1}`
\end{lstlisting}
% \vspace{-\baselineskip}
% Similarly, predecessors of each program statement are identified.
%, required for statement activation variable updates.
\item Predecessors of each statement.
% \vspace{-.5\baselineskip}
% \item \CHANGE{A list of mappings containing the predecessors of every statement in the corresponding thread's program.}
\begin{lstlisting}[style=c++]
std::vector<std::vector<std::vector<uint>>> predecessors = {
  // thread 0
  {{},    // `\lstCMT{ADDI 1}`    `\lstCMT{$\leftarrow$ $\varnothing$}`
   {0},   // `\lstCMT{STORE 0}`   `\lstCMT{$\leftarrow$ ADDI 1}`
   {1},   // `\lstCMT{LOAD 1}`    `\lstCMT{$\leftarrow$ STORE 0}`
   {2},   // `\lstCMT{CHECK 0}`   `\lstCMT{$\leftarrow$ LOAD 1}`
   {3}},  // `\lstCMT{HALT}`      `\lstCMT{$\leftarrow$ CHECK 0}`
  // thread 1
  {{},    // `\lstCMT{ADDI 1}`    `\lstCMT{$\leftarrow$ $\varnothing$}`
   {0},   // `\lstCMT{STORE 1}`   `\lstCMT{$\leftarrow$ ADDI 1}`
   {1},   // `\lstCMT{LOAD 0}`    `\lstCMT{$\leftarrow$ STORE 1}`
   {2},   // `\lstCMT{CHECK 0}`   `\lstCMT{$\leftarrow$ LOAD 0}`
   {3}},  // `\lstCMT{HALT}`      `\lstCMT{$\leftarrow$ CHECK 0}`
  // thread 2
  {{},    // `\lstCMT{CHECK 0}`   `\lstCMT{$\leftarrow$ $\varnothing$}`
   {0},   // `\lstCMT{ADD 0}`     `\lstCMT{$\leftarrow$ CHECK 0}`
   {1},   // `\lstCMT{ADD 1}`     `\lstCMT{$\leftarrow$ ADD 0}`
   {2},   // `\lstCMT{JZ error}`  `\lstCMT{$\leftarrow$ ADD 1}`
   {3},   // `\lstCMT{EXIT 0}`    `\lstCMT{$\leftarrow$ JZ error}`
   {3}}}; // `\lstCMT{error:}` `\lstCMT{EXIT 1}` `\lstCMT{$\leftarrow$ JZ error}`
\end{lstlisting}

% Program counters of barrier instructions, requiring an empty store buffer.
\item Program counters of memory barriers.
\begin{lstlisting}[style=c++]
std::vector<std::vector<uint>> barriers = {
  {1, 4}, // thread 0: `\lstCMT{STORE 0, HALT}`
  {1, 4}, // thread 1: `\lstCMT{STORE 1, HALT}`
  {}};    // thread 2: `\lstCMT{$\varnothing$}`
\end{lstlisting}

% Program counters of \lstASM{CHECK} instructions per checkpoint and thread.
\item Program counters of \lstASM{CHECK} instructions per checkpoint and thread.
\begin{lstlisting}[style=c++]
std::map<uint, std::map<uint, std::vector<uint>>> checkpoints = {
  // checkpoint 0
  {0, {{0, {3}},    // thread 0: `\lstCMT{CHECK 0}`
       {1, {3}},    // thread 1: `\lstCMT{CHECK 0}`
       {2, {0}}}}}; // thread 2: `\lstCMT{CHECK 0}`
\end{lstlisting}
\end{itemize}

% \begin{lstlisting}[style=c++]
% // pcs of halt statements
% //
% // thread -> list of program counters
% //
% std::map<uint, std::vector<uint>> halts = {{0, {4}}, {1, {4}}};
% \end{lstlisting}

% \begin{lstlisting}[style=c++]
% // pcs of exit calls
% //
% // thread -> list of program counters
% //
% std::map<uint, std::vector<uint>> exits = {{2, {4, 5}}};
% \end{lstlisting}

% \newpage
\subsection*{{\SMTLIB} Generator Functions}
% \todo[inline]{smtlib generator functions} %%%%%%%%%%%%%%%%%%%%%%%%%%%%%%%%%%%%%%

%Valid \SMTLIB expressions are generated by a set of wrapper functions for each command used in our encoding.
%> \SMTLIB expressions are generated by a set of functions, based on two templates implementing the basic expressions generators.
% using parameter packs to define variadic template functions, implementing the basic expressions generators.
% We start by defining a variadic function template to implement generators for a fixed number of arguments.
% The first for generating functions with an arbitrary but fixed number of arguments.

%Based on the previously defined template functions, we are now able to define a set of wrapper functions for generating the commands used in our encoding.%, shown in Table \ref{tbl:encoding:generators}.

% \SMTLIB commands used during encoding are generated by a set of functions, shown in Table \ref{tbl:encoding:generators}.

% Frequently used \SMTLIB expressions are generated by a set of functions, each exposing overloads for every possible number of arguments supported by the command.
% Frequently used \SMTLIB expressions are generated by a set of functions, each exposing a number of overloads for commands accepting a variable number of
% Frequently used \SMTLIB expressions are generated by a set of functions, including overloads for commands accepting a variable number of arguments.
%providing overloads to match a specific arity.

Frequently used {\SMTLIB} expressions are generated by a set of functions, based on the following variadic function template as the basic expression generator.

\begin{lstlisting}[style=c++]
template <class ... T>
std::string expr (const std::string & op, const T & ... args)
{
  std::string e = '(' + op;
  (((e += ' ') += args), ...);
  return e += ')';
}
\end{lstlisting}

\noindent
% \SMTLIB commands supporting a variable number of arguments by being defined as either \lstSMTLIB{:left-assoc} or \lstSMTLIB{:right-assoc} (like \lstSMTLIB{and}, \lstSMTLIB{or}, etc.), also include two additional overloads.
Generator functions for {\SMTLIB} commands supporting a variable number of arguments by being defined as either \lstSMTLIB{:left-assoc}, \lstSMTLIB{:right-assoc} or \lstSMTLIB{:chainable} %(like conjunction, disjunction, etc.),
also include two additional overloads:
% A variadic version to support an arbitrary but fixed number of arguments, defined as a template identically to the previous expression generatore \lstCPP{expr}.
a variadic version for an arbitrary but fixed number of arguments, using the same template definition as the basic expression generator \lstCPP{expr}
% and another version, accepting arguments contained in an arbitrary STL container for handling a variable number of arguments at runtime, based on the following template.
% and another for handling a varying number of arguments at runtime by accepting an arbitrary STL container, based on the following template.
and another for handling a varying number of arguments at runtime, based on the following template accepting an arbitrary STL container.
% Both overloads handle calls containing only a single argument by directly returning it in order to emit syntactically correct \SMTLIB expressions.
In order to emit syntactically correct {\SMTLIB} expressions, both handle calls containing only a single argument by directly returning it.

% A variadic version as a wrapper to the identically defined \lstCPP{expr} generator function for
% They also expose a variadic version, defined identically to the basic expression generator \lstCPP{expr} for an arbitrary but fixed number of arguments.
% One for generating commands with an arbitrary but fixed number of arguments,
% In order to generate commands with an arbitrary but fixed number of arguments, the according generator functions

% \todo[inline]{variable number of arguments}

\newpage

\begin{lstlisting}[style=c++]
template <template <class, class ...> class C>
std::string expr (const std::string & op, const C<std::string> & args)
{
  std::string e = '(' + op;
  for (const auto & a : args)
    (e += ' ') += a;
  return e += ')';
}
\end{lstlisting}

\noindent
Table \ref{tbl:encoding:generators} shows all generator function names and the corresponding {\SMTLIB} command they are creating.
% arity

\begin{table}[!h]
% \noindent
\begin{tabu} to \textwidth {|X|l|X|}
  \firsthline
  \lstCPP{declare_bool}   & declare boolean   & \lstSMTLIB{(declare-fun ...)} \\
  \lstCPP{declare_bv}     & declare bit-vector & \\%\lstSMTLIB{(declare-fun ... () (_ BitVec 16))} \\
  \lstCPP{declare_array}  & declare array     & \\%\lstSMTLIB{(declare-fun ... () (Array (_ BitVec 16) (_ BitVec 16)))} \\
  \hline
  \hline
  \lstCPP{assertion} & assertion & \lstSMTLIB{(assert ...)} \\
  \hline
  \hline
  \lstCPP{lnot}   & negation                  & \lstSMTLIB{(not ...)} \\
  \lstCPP{land}   & conjunction               & \lstSMTLIB{(and ...)} \\
  \lstCPP{lor}    & disjunction               & \lstSMTLIB{(or ...)} \\
  \lstCPP{lxor}   & exclusive disjunction     & \lstSMTLIB{(xor ...)} \\
  \lstCPP{imply}  & implication               & \lstSMTLIB{(=> ...)} \\
  \lstCPP{equal}  & equivalence               & \lstSMTLIB{(= ...)} \\
  \lstCPP{ite}    & functional if-then-else   & \lstSMTLIB{(ite ...)} \\
  \hline
  \hline
  \lstCPP{bvadd}  & bit-vector addition        & \lstSMTLIB{(bvadd ...)} \\
  \lstCPP{bvsub}  & bit-vector subtraction     & \lstSMTLIB{(bvsub ...)} \\
  \lstCPP{bvmul}  & bit-vector multiplication  & \lstSMTLIB{(bvmul ...)} \\
  \hline
  \hline
  \lstCPP{select}   & array read            & \lstSMTLIB{(select ...)} \\
  \lstCPP{store}    & array store           & \lstSMTLIB{(store ...)} \\
  \lstCPP{extract}  & bit-vector extraction  & \lstSMTLIB{(extract ...)} \\
  \lasthline
\end{tabu}
\caption{{\SMTLIB} \cite{ref:SMT-Lib} expression generator functions.}
\label{tbl:encoding:generators}
\end{table}

% \begin{longtabu}{|l|X|}
  % \firsthline
  % \lstCPP{assertion (std::string)} & \lstSMTLIB{(assert x)} \\
  % % lnot (std::string x) & foo \\
  % % land (std::string) & bar \\
  % \lstCPP{land (std::string)} & \lstSMTLIB{x} \\
  % % land (std::string, std::string ...args) & bar \\
  % \lstCPP{land (std::string, ...)} & \lstSMTLIB{(and x ...)} \\
  % % land (C<std::string>) & bar \\
  % \lstCPP{land (C<std::string>)} & \lstSMTLIB{(and x ...)} \\
  % \lasthline
% \end{longtabu}
% Functions for fixed number of arguments like \lstSMTLIB{or} use the following variadic function template, based on
% make use of \emph{parameter packs} to create a variadic function template
% \todo[inline]{arbitrary but fixed number of arguments}

% \begin{lstlisting}[style=c++]
% std::string assertion (std::string arg)
% {
  % return expr("assert", arg);
% }
% \end{lstlisting}
%
% \begin{lstlisting}[style=c++]
% std::string lnot (std::string arg)
% {
  % return expr("not", arg);
% }
% \end{lstlisting}
%
% \begin{lstlisting}[style=c++]
% std::string land (const std::string & arg) { return arg; }
%
% template <class ... T>
% std::string land (const std::string & arg, const T & ... args)
% {
  % return expr("and", arg, args...);
% }
%
% template <template<class, class...> class C>
% std::string land (const C<std::string> & args)
% {
  % return args.size() < 2 ? *args.begin() : expr("and", args);
% }
% \end{lstlisting}
%
% \begin{lstlisting}[style=c++]
% std::string lor (const std::string & arg) { return arg; }
%
% template <class ... T>
% std::string lor (const std::string & arg, const T & ... args)
% {
  % return expr("or", arg, args...);
% }
%
% template <template<class, class...> class C>
% std::string lor (const C<std::string> & args)
% {
  % return args.size() < 2 ? *args.begin() : expr("or", args);
% }
% \end{lstlisting}
%
% \begin{lstlisting}[style=c++]
% std::string lxor (const std::string & arg) { return arg; }
%
% template <class ... T>
% std::string lxor (const std::string & arg, const T & ... args)
% {
  % return expr("xor", arg, args...);
% }
%
% template <template<class, class...> class C>
% std::string lxor (const C<std::string> & args)
% {
  % return args.size() < 2 ? *args.begin() : expr("xor", args);
% }
% \end{lstlisting}
%
% \begin{lstlisting}[style=c++]
% std::string implication (std::string antecedent, std::string consequent)
% {
  % return expr("=>", antecedent, consequent);
% }
% \end{lstlisting}
%
% \begin{lstlisting}[style=c++]
% std::string equality (std::string var, std::string val)
% {
  % return expr("=", var, val);
% }
% \end{lstlisting}
%
% \begin{lstlisting}[style=c++]
% std::string ite (std::string condition, std::string t, std::string f)
% {
  % return expr("ite", condition, t, f);
% }
% \end{lstlisting}
%
% \begin{lstlisting}[style=c++]
% std::string bvadd (std::string a, std::string b)
% {
  % return expr("bvadd", a, b);
% }
% \end{lstlisting}
%
% \begin{lstlisting}[style=c++]
% std::string bvsub (std::string a, std::string b)
% {
  % return expr("bvsub", a, b);
% }
% \end{lstlisting}
%
% \begin{lstlisting}[style=c++]
% std::string bvmul (std::string a, std::string b)
% {
  % return expr("bvmul", a, b);
% }
% \end{lstlisting}
%
% \begin{lstlisting}[style=c++]
% std::string select (std::string var, std::string idx)
% {
  % return expr("select", var, idx);
% }
% \end{lstlisting}

% \begin{lstlisting}[style=c++]
% std::string store (std::string var, std::string idx, std::string val)
% {
  % return expr("store", var, idx, val);
% }
% \end{lstlisting}
%
% \begin{lstlisting}[style=c++]
% std::string extract (std::string msb, std::string lsb, std::string var)
% {
  % return expr(expr("_ extract", msb, lsb).c_str(), var);
% }
% \end{lstlisting}
%
% \begin{lstlisting}[style=c++]
% std::string declare_bool_var (std::string name)
% {
  % // return "(declare-fun " + name + " () Bool)";
  % return expr("declare-fun", {name, "()", "Bool"});
% }
% \end{lstlisting}
%
% \begin{lstlisting}[style=c++]
% std::string declare_bv_var (std::string name)
% {
  % // return "(declare-fun " + name + " () (_ BitVec 16))";
  % return expr("declare-fun", {name, "()", "(_ BitVec 16)"});
% }
% \end{lstlisting}
%
% \begin{lstlisting}[style=c++]
% std::string declare_array_var (std::string name)
% {
  % // return "(declare-fun "
  % //        + name
  % //        + " (Array (_ BitVec 16) (_ BitVec 16)))";
  % static const char * sort = "(_ BitVec 16)";
  % return expr("declare-fun", name, expr("Array", sort, sort));
% }
% \end{lstlisting}

\noindent
Finally, we introduce a function \lstCPP{std::string consth (uint val)} for generating hexadecimal bit-vector constants and a helper function to simplify variable assignments.

% \noindent
% Hexadecimal bit-vector constants are generated by the following function.
%
% \begin{lstlisting}[style=c++]
% std::string consth (uint val)
% {
  % std::ostringstream s;
  % s << "#x" << std::setfill('0') << std::setw(4) << std::hex << val;
  % return s.str();
% }
% \end{lstlisting}
%
% \noindent
% Finally, we also define a helper for assigning states to a specific value.

\begin{lstlisting}[style=c++]
std::string assign (std::string var, std::string val)
{
  return assertion(equal(var, val));
}
\end{lstlisting}

% \newpage

\subsection*{Common Encoding Functions}
% \todo[inline]{main encoding function (smtlib)} %%%%%%%%%%%%%%%%%%%%%%%%%%%%%%%%%

We will now give a top-down overview of the actual encoding process \CHANGE{implemented in \texttt{src/encoder_smtlib.cc}} and start with the main encoding function defined below.
% \lstCPP{encode}.

\begin{lstlisting}[style=c++]
void encode ()
{
  formula << "(set-logic QF_AUFBV)\n";

  for (uint k = 0; k <= bound; k++)
    {
      declare_states(k);
      declare_transitions(k);
      define_transitions (k);

      if (k)
        define_states(k);
      else
        init_states();
    }
}
\end{lstlisting}

\noindent
After setting the appropriate logic (\lstSMTLIB{QF_AUFBV}), each individual step $k$ is encoded iteratively, starting with the declaration of state variables.
% At every step $k$, we start by declaring  variables

% \todo[inline]{declare states} %%%%%%%%%%%%%%%%%%%%%%%%%%%%%%%%%%%%%%%%%%%%%%%%%%

\begin{lstlisting}[style=c++]
void declare_states (uint k)
{
  // thread states
  declare_accu(k);
  declare_mem(k);
  declare_adr(k);
  declare_val(k);
  declare_full(k);
  declare_stmt(k);
  declare_block(k);
  declare_halt(k);

  // machine states
  declare_heap(k);
  declare_exit(k);

  if (!k) declare_exit_code();
}
\end{lstlisting}

% \todo[inline]{declare accu} %%%%%%%%%%%%%%%%%%%%%%%%%%%%%%%%%%%%%%%%%%%%%%%%%%%%

% \begin{lstlisting}[style=c++]
% template <template<class, class...> class C>
% std::string var (std::string name, C<uint> attributes)
% {
  % for (const auto & attr : attributes)
    % (name += '_') += std::to_string(attr);
%
  % return name;
% }
% \end{lstlisting}

\noindent
Similarly to the generation of {\SMTLIB} commands, the following variadic function template serves as the basis for generating symbols, matching our established variable naming scheme.

\begin{lstlisting}[style=c++]
template <class ... T>
std::string var (const std::string & name, const T & ... attributes)
{
  return (((name += '_') += std::to_string(attributes)), ...);
}
\end{lstlisting}

\noindent
% The generator function for our accumulator register symbols $\ACCU^k_t$, for example, is given below and
This allows us to define a generator function for each variable as a simple wrapper, % to \lstCPP{var}.
like the one of our accumulator register variables $\ACCU^k_t$ given below.
% and for our accumulator register variable $\ACCU^k_t$, for example, created

% As an example, the generator function for accumulator register variables $\ACCU^k_t$ is given below.
% For example, the generator function of our accumulator register variables $\ACCU^k_t$ is given below.

% To further increase readability, we introduce a generator function for each variable.

% return "accu_" + std::to_string(k) + '_' + std::to_string(t);
\begin{lstlisting}[style=c++]
std::string accu (uint k, uint t) { return var("accu", k, t); }
\end{lstlisting}

\noindent
The declaration of accumulator register states for each thread $t$ in step $k$ is now appended to the formula by using previously defined generator functions.
% The accumulator register state variable declaration for each thread $t$ in step $k$ is now appended to the formula by using the previously defined generator functions.
% The declaration of accumulator register state variables is now appended to the formula, using the previously defined generator expressions.
% Using the previously defined generator functions, the accumulator register state declarations  of a specific step $k$ are now appended to the formula.
% Declaration of all accumulator register variables for a specific step $k$

\begin{lstlisting}[style=c++]
void declare_accu (uint k)
{
  for (uint t = 0; t < programs.size(); t++)
    formula << declare_bv(accu(k, t)) << '\n';
}
\end{lstlisting}

\noindent
% Declaring accumulator register states for the initial step, required in our demo a, leads to the following output being produced.
% Executing \lstCPP{declare_accu} in the initial step $k = 0$ of our demo example would lead to the following declarations being appended to the formula.
% Calling \lstCPP{declare_accu} for the initial step $k = 0$ in our demo example would lead to the following register state variables being declared.
% Example output: \lstCPP{declare_accu(0)}.
\emph{Example:} \lstCPP{declare_accu(0)}

% \lstinputlisting[style=smtlib, linerange={9-12}]{../../examples/demo/formula.functional.smt2}
% \begin{lstlisting}[style=smtlib, caption={Example output \lstCPP{declare_accu(0)}}]
\begin{lstlisting}[style=smtlib]
(declare-fun accu_0_0 () (_ BitVec 16))
(declare-fun accu_0_1 () (_ BitVec 16))
(declare-fun accu_0_2 () (_ BitVec 16))
\end{lstlisting}

\noindent
% We skip the declaration functions of the remaining register states, as they are defined identically.
While the declarator functions of the remaining register states
% The declarator functions of the remaining register states
$\MEM^k_t$, $\SBADR^k_t$, $\SBVAL^k_t$ and $\SBFULL^k_t$
% \lstCPP{declare_mem},
% \lstCPP{declare_adr},
% \lstCPP{declare_val}
% and \lstCPP{declare_full}
are defined almost identically, except for the latter
% being declared as \lstSMTLIB{Bool},
changing the sort to \lstSMTLIB{Bool},
% we need to additionally iterate the program of each thread to declare the required number of statement activation variables.
% an additional nested loop is needed to declare the required number of statement activation variables for each thread.
% we need a nested loop to declare the required number of statement activation variables for each thread.
% we additionally need to iterate the program of each thread to declare the required number of statement activation variables.
% an additional nested loop is needed to declare the required number of statement activation variables for a given thread's program.
an additional nested loop is needed to declare an activation variable for each statement in a thread's program.
% a nested loop is required to declare an activation variable for each statement of the

% \todo[inline]{declare mem} %%%%%%%%%%%%%%%%%%%%%%%%%%%%%%%%%%%%%%%%%%%%%%%%%%%%%

% \begin{lstlisting}[style=c++]
% std::string mem_var (uint k, uint t)
% {
  % return var("mem", k, t);
% }
% \end{lstlisting}

% \begin{lstlisting}[style=c++]
% void declare_mem (uint k)
% {
  % for (uint t = 0; t < programs.size(); t++)
    % formula << declare_bv_var(mem_var(k, t)) << eol;
  % formula << eol;
% }
% \end{lstlisting}

% \begin{lstlisting}[style=smtlib]
% ; mem variables - mem_<step>_<thread>
% (declare-fun mem_0_0 () (_ BitVec 16))
% (declare-fun mem_0_1 () (_ BitVec 16))
% (declare-fun mem_0_2 () (_ BitVec 16))
% \end{lstlisting}

% \todo[inline]{declare adr} %%%%%%%%%%%%%%%%%%%%%%%%%%%%%%%%%%%%%%%%%%%%%%%%%%

% \begin{lstlisting}[style=c++]
% std::string adr_var (uint k, uint t)
% {
  % return var("adr", k, t);
% }
% \end{lstlisting}

% \begin{lstlisting}[style=c++]
% void declare_adr (uint k)
% {
  % for (uint t = 0; t < programs.size(); t++)
    % formula << declare_bv_var(adr_var(k, t)) << eol;
  % formula << eol;
% }
% \end{lstlisting}

% \begin{lstlisting}[style=smtlib]
% ; store buffer address variables - adr_<step>_<thread>
% (declare-fun adr_0_0 () (_ BitVec 16))
% (declare-fun adr_0_1 () (_ BitVec 16))
% (declare-fun adr_0_2 () (_ BitVec 16))
% \end{lstlisting}

% \todo[inline]{declare val} %%%%%%%%%%%%%%%%%%%%%%%%%%%%%%%%%%%%%%%%%%%%%%%%%%

% \begin{lstlisting}[style=c++]
% std::string val_var (uint k, uint t)
% {
  % return var("val", k, t);
% }
% \end{lstlisting}

% \begin{lstlisting}[style=c++]
% void declare_val (uint k)
% {
  % for (uint t = 0; t < programs.size(); t++)
    % formula << declare_bv_var(val_var(k, t)) << eol;
  % formula << eol;
% }
% \end{lstlisting}

% \begin{lstlisting}[style=smtlib]
% ; store buffer value variables - val_<step>_<thread>
% (declare-fun val_0_0 () (_ BitVec 16))
% (declare-fun val_0_1 () (_ BitVec 16))
% (declare-fun val_0_2 () (_ BitVec 16))
% \end{lstlisting}

% \todo[inline]{declare full} %%%%%%%%%%%%%%%%%%%%%%%%%%%%%%%%%%%%%%%%%%%%%%%%%

% \begin{lstlisting}[style=c++]
% std::string full_var (uint k, uint t)
% {
  % return var("full", k, t);
% }
% \end{lstlisting}

% \begin{lstlisting}[style=c++]
% void declare_full (uint k)
% {
  % for (uint t = 0; t < programs.size(); t++)
    % formula << declare_bool_var(full_var(k, t)) << eol;
  % formula << eol;
% }
% \end{lstlisting}

% \begin{lstlisting}[style=smtlib]
% ; store buffer full variables - full_<step>_<thread>
% (declare-fun full_0_0 () Bool)
% (declare-fun full_0_1 () Bool)
% (declare-fun full_0_2 () Bool)
% \end{lstlisting}

% \todo[inline]{declare stmt} %%%%%%%%%%%%%%%%%%%%%%%%%%%%%%%%%%%%%%%%%%%%%%%%%%%%

% \begin{lstlisting}[style=c++]
% std::string stmt_var (uint k, uint t, uint pc)
% {
  % return var("stmt", k, t, pc);
% }
% \end{lstlisting}

% \noindent
% according to Table \ref{tbl:encoding:programs}

\begin{lstlisting}[style=c++]
void declare_stmt (uint k)
{
  for (uint t = 0; t < programs.size(); t++)
    for (uint pc = 0; pc < programs[t].size(); pc++)
      formula << declare_bool(stmt(k, t, pc)) << '\n';
}
\end{lstlisting}

\newpage

\noindent
% According to our demo example programs, given in Table \ref{tbl:encoding:programs}, calling \lstCPP{declare_stmt} for the initial step $k = 0$ would lead to the following statement activation variables being declared.
% Example output: \lstCPP{declare_stmt(0)}.
\emph{Example:} \lstCPP{declare_stmt(0)}

\begin{lstlisting}[style=smtlib]
(declare-fun stmt_0_0_0 () Bool)
(declare-fun stmt_0_0_1 () Bool)
(declare-fun stmt_0_0_2 () Bool)
(declare-fun stmt_0_0_3 () Bool)
(declare-fun stmt_0_0_4 () Bool)

(declare-fun stmt_0_1_0 () Bool)
(declare-fun stmt_0_1_1 () Bool)
(declare-fun stmt_0_1_2 () Bool)
(declare-fun stmt_0_1_3 () Bool)
(declare-fun stmt_0_1_4 () Bool)

(declare-fun stmt_0_2_0 () Bool)
(declare-fun stmt_0_2_1 () Bool)
(declare-fun stmt_0_2_2 () Bool)
(declare-fun stmt_0_2_3 () Bool)
(declare-fun stmt_0_2_4 () Bool)
(declare-fun stmt_0_2_5 () Bool)
\end{lstlisting}

% \todo[inline]{declare block} %%%%%%%%%%%%%%%%%%%%%%%%%%%%%%%%%%%%%%%%%%%%%%%%%%%

% \begin{lstlisting}[style=c++]
% std::string block_var (uint k, uint id, uint t)
% {
  % return var("block", k, id, t);
% }
% \end{lstlisting}

\noindent
% Declaration of block flags uses the \lstCPP{checkpoints} map to
% Block flags are declared on the basis of the \lstCPP{checkpoints} map, created during
In order to \CHANGE{avoid} repeated iteration of input programs,
the precomputed utility variables are used to indicate the presence of certain instructions and other control flow related information.
% processing of variables  on certain instructions
% the precomputed utility variables are used for processing variables related to a certain type of instruction.
% Declaration of block flags therefore uses the \lstCPP{checkpoints} map to
% Declaration of block flags therefore boils down to simply iterating the
% Declaration of block flags can therefore be achieved by simply iterating the \lstCPP{checkpoints} map.
% Declaration of block flags, for example, can therefore be simplified by relying on the \lstCPP{checkpoints} map,
Declaration of block flags $\BLOCK^k_{id, t}$ can therefore be simplified by relying on the \lstCPP{checkpoints} map,
% containing the program counters of \lstASM{CHECK} instructions.
containing the list of threads eventually waiting for a specific checkpoint.

\begin{lstlisting}[style=c++]
void declare_block (uint k)
{
  for (const auto & [id, threads] : checkpoints)
    for (const auto & [t, _] : threads)
      formula << declare_bool(block(k, id, t)) << '\n';
}
\end{lstlisting}

\noindent
% Example output: \lstCPP{declare_block(0)}.
\emph{Example:} \lstCPP{declare_block(0)}

\begin{lstlisting}[style=smtlib]
(declare-fun block_0_0_0 () Bool)
(declare-fun block_0_0_1 () Bool)
(declare-fun block_0_0_2 () Bool)
\end{lstlisting}

% \todo[inline]{declare halt} %%%%%%%%%%%%%%%%%%%%%%%%%%%%%%%%%%%%%%%%%%%%%%%%%%%%

% \begin{lstlisting}[style=c++]
% std::string halt_var (uint k, uint t)
% {
  % return var("halt", k, t);
% }
% \end{lstlisting}

\newpage

\noindent
% Halt flags $\HALT^k_t$ are declared in a similar manner, based on the threads contained in the \lstCPP{updates} map's \lstCPP{State::halt} entry.
Halt flags $\HALT^k_t$ are again declared \CHANGE{by the appropriate generator function.}

\begin{lstlisting}[style=c++]
void declare_halt (uint k)
{
  for (uint t = 0; t < programs.size(); t++)
    formula << declare_bool(halt(k, t)) << '\n';
}
\end{lstlisting}

\noindent
% \textbf{Output:} \lstCPP{declare_halt(0)}
\emph{Example:} \lstCPP{declare_halt(0)}

\begin{lstlisting}[style=smtlib]
(declare-fun halt_0_0 () Bool)
(declare-fun halt_0_1 () Bool)
(declare-fun halt_0_2 () Bool)
\end{lstlisting}

% \todo[inline]{declare heap} %%%%%%%%%%%%%%%%%%%%%%%%%%%%%%%%%%%%%%%%%%%%%%%%%%%%

% \begin{lstlisting}[style=c++]
% std::string heap_var (uint k)
% {
  % return var("heap", k);
% }
% \end{lstlisting}

% \noindent
% With all thread states being declared, we continue with machine states, starting with our shared memory array $\HEAP^k$.

% \begin{lstlisting}[style=c++]
% void declare_heap (uint k)
% {
  % formula << declare_array(heap_var(k)) << "\n\n";
% }
% \end{lstlisting}

% \noindent
% Example output: \lstCPP{declare_heap(0)}.

% \begin{lstlisting}[style=smtlib]
% (declare-fun heap_0 () (Array (_ BitVec 16) (_ BitVec 16)))
% \end{lstlisting}

% \todo[inline]{declare exit-flag} %%%%%%%%%%%%%%%%%%%%%%%%%%%%%%%%%%%%%%%%%%%%%%%

% \begin{lstlisting}[style=c++]
% std::string exit_flag_var (uint k)
% {
  % return var("exit", k);
% }
% \end{lstlisting}

% Followed by the exit flag $\EXIT^k$, which is only added if the programs contain an uninterruptible infinite loop.

% \begin{lstlisting}[style=c++]
% void declare_exit_flag (uint k)
% {
  % if (halts.empty() && exits.empty())
    % return;
%
  % formula << declare_bool(exit_flag_var(k)) << eol << eol;
% }
% \end{lstlisting}

% \noindent
% Followed by the exit flag $\EXIT^k$.

% \begin{lstlisting}[style=c++]
% void declare_exit (uint k)
% {
  % formula << declare_bool(exit_flag_var(k)) << "\n\n";
% }
% \end{lstlisting}

% \noindent
% Example output: \lstCPP{declare_exit(0)}.

% \begin{lstlisting}[style=smtlib]
% (declare-fun exit_0 () Bool)
% \end{lstlisting}

% \todo[inline]{declare exit-code} %%%%%%%%%%%%%%%%%%%%%%%%%%%%%%%%%%%%%%%%%%%%%%%

\noindent
% With all thread state variables for the specific step being declared, machine states $\HEAP^k$, $\EXIT^k$ are next, appending the corresponding variable using the required sort.
% With all thread state variables for the specific step being declared, machine states $\HEAP^k$, $\EXIT^k$ are next.
% We omit the definition of the corresponding declarator functions as they are trivial and just add a single variable for each step using the appropriate sort.
% With all thread state variables being declared, machine states $\HEAP^k$, $\EXIT^k$ and $\EXITCODE^k$ are next.
Machine states $\HEAP^k$, $\EXIT^k$ and $\EXITCODE^k$ are declared next.
In contrast to our basic encoding scheme, only a single bit-vector variable representing the machine's final exit code is used and declared during the initial step.
We omit the definition of the corresponding declarator functions as they just add a single variable with the appropriate sort.

~\\
\noindent
\emph{Example:} \lstCPP{declare_heap(0)}

\begin{lstlisting}[style=smtlib]
(declare-fun heap_0 () (Array (_ BitVec 16) (_ BitVec 16)))
\end{lstlisting}

\noindent
\emph{Example:} \lstCPP{declare_exit(0)}.
% \textbf{Output:} \lstCPP{declare_exit(0)}

\begin{lstlisting}[style=smtlib]
(declare-fun exit_0 () Bool)
\end{lstlisting}

\noindent
\emph{Example:} \lstCPP{declare_exit_code()}

\begin{lstlisting}[style=smtlib]
(declare-fun exit-code () (_ BitVec 16))
\end{lstlisting}

% \noindent
% In contrast to our basic encoding scheme, we include only a single bit-vector variable representing the machine's final exit code, declared during the initial step.% $k = 0$.

% \begin{lstlisting}[style=c++]
% std::string exit_code_var = "exit-code";
%
% void declare_exit_code ()
% {
  % formula << declare_bv(exit_code_var) << "\n\n";
% }
% \end{lstlisting}
%
% \noindent
% \emph{Example:} \lstCPP{declare_exit_code()}
%
% \begin{lstlisting}[style=smtlib]
% (declare-fun exit-code () (_ BitVec 16))
% \end{lstlisting}

% ~\\
% \noindent
% \emph{Example:} \lstCPP{declare_heap(0), declare_exit(0), declare_exit_code()}
%
% \begin{lstlisting}[style=smtlib]
% (declare-fun heap_0 () (Array (_ BitVec 16) (_ BitVec 16)))
% (declare-fun exit_0 () Bool)
% (declare-fun exit-code () (_ BitVec 16))
% \end{lstlisting}

% \todo[inline]{declare transitions} %%%%%%%%%%%%%%%%%%%%%%%%%%%%%%%%%%%%%%%%%%%%%

\noindent
% With all state variables being declared, we continue with transition variables.
% Declaration of the remaining transition and helper variables is performed by the following functions.
% Declaration of the remaining transition and helper variables is consolidated in the following function.
% Declaration of the remaining transition and helper variables is consolidated in the following function.
% The remaining variables
% Finally, the remaining transition and helper variables are declared.
Variable declaration is concluded by adding the required transition variables.

\begin{lstlisting}[style=c++]
void declare_transitions (uint k)
{
  declare_thread(k);
  declare_flush(k);
  declare_exec(k);
  declare_check(k);
}
\end{lstlisting}

% \todo[inline]{declare thread} %%%%%%%%%%%%%%%%%%%%%%%%%%%%%%%%%%%%%%%%%%%%%%%%%%

% \begin{lstlisting}[style=c++]
% std::string thread_var (uint k, uint t)
% {
  % return var("thread", k, t);
% }
% \end{lstlisting}

\noindent
% Again, we skip further details about the declarator functions for $\THREAD^k$, $\EXEC^k_{t, pc}$ and $\FLUSH^k_t$, as they are based on the same principles used for declaring thread states.
% We skip further details about the declarator functions for $\THREAD^k$, $\EXEC^k_{t, pc}$ and $\FLUSH^k_t$, as they are again defined almost identically to those of register states and statement activation variables.
% We skip further details about the declarator functions for $\THREAD^k$, $\EXEC^k_{t, pc}$ and $\FLUSH^k_t$, as they are based on the same principles used for the declaration of register states and statement activation variables.
We skip further details about their declarator functions, as they are based on the same principles used for register states and statement activation variables, except $\CHECK^k_{id}$ which also refers to the identifiers stored in the \lstCPP{checkpoints} map.

% \begin{lstlisting}[style=c++]
% void declare_thread (uint k)
% {
  % for (uint t = 0; t < programs.size(); t++)
    % formula << declare_bool(thread_var(k, t)) << '\n';
  % formula << eol;
% }
% \end{lstlisting}
%
% \begin{lstlisting}[style=smtlib]
% ; thread activation variables - thread_<step>_<thread>
% (declare-fun thread_0_0 () Bool)
% (declare-fun thread_0_1 () Bool)
% (declare-fun thread_0_2 () Bool)
% \end{lstlisting}

% \todo[inline]{declare exec} %%%%%%%%%%%%%%%%%%%%%%%%%%%%%%%%%%%%%%%%%%%%%%%%%%%%

% \begin{lstlisting}[style=c++]
% std::string exec_var (uint k, uint t, uint pc)
% {
  % return var("exec", k, t, pc);
% }
% \end{lstlisting}
%
% \begin{lstlisting}[style=c++]
% void declare_exec (uint k)
% {
  % for (uint t = 0; t < programs.size(); t++)
    % for (uint pc = 0; pc < programs[t].size(); pc++)
      % formula << declare_bool_var(exec_var(k, t, pc)) << eol;
  % formula << eol;
% }
% \end{lstlisting}
%
% \begin{lstlisting}[style=smtlib]
% ; statement execution variables - exec_<step>_<thread>_<pc>
% (declare-fun exec_0_0_0 () Bool)
% (declare-fun exec_0_0_1 () Bool)
% (declare-fun exec_0_0_2 () Bool)
% (declare-fun exec_0_0_3 () Bool)
% (declare-fun exec_0_0_4 () Bool)
%
% (declare-fun exec_0_1_0 () Bool)
% (declare-fun exec_0_1_1 () Bool)
% (declare-fun exec_0_1_2 () Bool)
% (declare-fun exec_0_1_3 () Bool)
% (declare-fun exec_0_1_4 () Bool)
%
% (declare-fun exec_0_2_0 () Bool)
% (declare-fun exec_0_2_1 () Bool)
% (declare-fun exec_0_2_2 () Bool)
% (declare-fun exec_0_2_3 () Bool)
% (declare-fun exec_0_2_4 () Bool)
% (declare-fun exec_0_2_5 () Bool)
% \end{lstlisting}

% \todo[inline]{declare flush} %%%%%%%%%%%%%%%%%%%%%%%%%%%%%%%%%%%%%%%%%%%%%%%%%%%
%
% \begin{lstlisting}[style=c++]
% std::string flush_var (uint k, uint t)
% {
  % return var("flush", k, t);
% }
% \end{lstlisting}
%
% \begin{lstlisting}[style=c++]
% void declare_flush (uint k)
% {
  % for (uint t = 0; t < programs.size(); t++)
    % formula << declare_bool_var(flush_var(k, t)) << eol;
  % formula << eol;
% }
% \end{lstlisting}
%
% \begin{lstlisting}[style=smtlib]
% ; store buffer flush variables - flush_<step>_<thread>
% (declare-fun flush_0_0 () Bool)
% (declare-fun flush_0_1 () Bool)
% (declare-fun flush_0_2 () Bool)
% \end{lstlisting}

% \todo[inline]{declare check} %%%%%%%%%%%%%%%%%%%%%%%%%%%%%%%%%%%%%%%%%%%%%%%%%%%

% \begin{lstlisting}[style=c++]
% std::string check_var (uint k, uint id)
% {
  % return var("check", k, id);
% }
% \end{lstlisting}

\begin{lstlisting}[style=c++]
void declare_check (uint k)
{
  for (const auto & [id, _] : checkpoints)
    formula << declare_bool(check(k, id)) << '\n';
}
\end{lstlisting}

\noindent
\emph{Example:} \lstCPP{declare_check(0)}

\begin{lstlisting}[style=smtlib]
(declare-fun check_0_0 () Bool)
\end{lstlisting}

% \todo[inline]{define transitions} %%%%%%%%%%%%%%%%%%%%%%%%%%%%%%%%%%%%%%%%%%%%%%
%
% \noindent
% Since the helper variables $\EXEC^k_{t, pc}$ and $\CHECK^k_{id}$ only depend on the transition variables, non-deterministically determined by our scheduling constraint, we are able to define them immediately.
% With all variables being declared, the actual encoding starts with the definition of helper variables.
% With all variables being declared, we are now able to start the actual encoding process by defining our helper variables.
%
% \begin{lstlisting}[style=c++]
% void define_helper (uint k)
% {
  % define_exec(k);
  % define_check(k);
% }
% \end{lstlisting}

% \todo[inline]{define constraints} %%%%%%%%%%%%%%%%%%%%%%%%%%%%%%%%%%%%%%%%%%%%%%

\noindent
% With all variables being declared, we are now able to start the actual encoding process by defining transition variables and scheduling constraints.
% With all variables being declared, we are now able to define transition variables and scheduling constraints.
% With all variables being declared, the actual encoding process is started by defining the transition variables according to our machine model's scheduling constraints.
With all variables being declared, the actual encoding process is started by defining helper variables and scheduling constraints according to our machine model.

\begin{lstlisting}[style=c++]
void define_transitions (uint k)
{
  define_exec(k);
  define_check(k);
  define_cardinality_constraints(k);
  define_store_buffer_constraints(k);
  define_checkpoint_constraints(k);
  define_halt_constraints(k);
}
\end{lstlisting}

% \todo[inline]{define exec} %%%%%%%%%%%%%%%%%%%%%%%%%%%%%%%%%%%%%%%%%%%%%%%%%%%%%

\noindent
Execution variables $\EXEC^k_{t, pc}$ are defined as a conjunction of the corresponding statement and thread activation variables.

\begin{lstlisting}[style=c++]
void define_exec (uint k)
{
  for (uint t = 0; t < programs.size(); t++)
    for (uint pc = 0; pc < programs[t].size(); pc++)
      formula << assign(exec(k, t, pc),
                        land(stmt(k, t, pc),
                             thread(k, t))) << '\n';
}
\end{lstlisting}

\noindent
\emph{Example:} \lstCPP{define_exec(0)}

\begin{lstlisting}[style=smtlib]
(assert (= exec_0_0_0 (and stmt_0_0_0 thread_0_0)))
(assert (= exec_0_0_1 (and stmt_0_0_1 thread_0_0)))
(assert (= exec_0_0_2 (and stmt_0_0_2 thread_0_0)))
(assert (= exec_0_0_3 (and stmt_0_0_3 thread_0_0)))
(assert (= exec_0_0_4 (and stmt_0_0_4 thread_0_0)))

(assert (= exec_0_1_0 (and stmt_0_1_0 thread_0_1)))
(assert (= exec_0_1_1 (and stmt_0_1_1 thread_0_1)))
(assert (= exec_0_1_2 (and stmt_0_1_2 thread_0_1)))
(assert (= exec_0_1_3 (and stmt_0_1_3 thread_0_1)))
(assert (= exec_0_1_4 (and stmt_0_1_4 thread_0_1)))

(assert (= exec_0_2_0 (and stmt_0_2_0 thread_0_2)))
(assert (= exec_0_2_1 (and stmt_0_2_1 thread_0_2)))
(assert (= exec_0_2_2 (and stmt_0_2_2 thread_0_2)))
(assert (= exec_0_2_3 (and stmt_0_2_3 thread_0_2)))
(assert (= exec_0_2_4 (and stmt_0_2_4 thread_0_2)))
(assert (= exec_0_2_5 (and stmt_0_2_5 thread_0_2)))
\end{lstlisting}

% \todo[inline]{define check} %%%%%%%%%%%%%%%%%%%%%%%%%%%%%%%%%%%%%%%%%%%%%%%%%%%%

\noindent
% Definition of $\CHECK^k_{id}$, signalling that all participating threads synchronized upon a given checkpoint, uses the lists of threads contained in the \lstCPP{checkpoint} map to create the required conjunction of block variables for every checkpoint.
Definition of checkpoint variables $\CHECK^k_{id}$ is again based on the threads stored in \lstCPP{checkpoints} to generate the list of corresponding block variables $\BLOCK^k_{id, t}$ for creating the required conjunction.

\begin{lstlisting}[style=c++]
void define_check (uint k)
{
  for (const auto & [id, threads] : checkpoints)
    {
      std::vector<std::string> args;
      for (const auto & [t, _] : threads)
        args.push_back(block(k, id, t));
      formula << assign(check(k, id), land(args)) << '\n';
    }
}
\end{lstlisting}

% \begin{lstlisting}[style=smtlib]
% ; checkpoint variables - check_<step>_<id>
% (assert (not check_0_0))
% \end{lstlisting}

\noindent
\emph{Example:} \lstCPP{define_check(0)}

\begin{lstlisting}[style=smtlib]
(assert (= check_0_0 (and block_0_0_0 block_0_0_1 block_0_0_2)))
\end{lstlisting}

% \todo[inline]{define scheduling constraints} %%%%%%%%%%%%%%%%%%%%%%%%%%%%%%%%%%%

% \begin{lstlisting}[style=c++]
% std::string card_constraint_naive (std::vector<std::string> vars)
% {
%   switch (vars.size())
%     {
%     case 1: formula << assertion(vars.front()) << '\n'; return;
%     case 2: formula << assertion(lxor(vars)) << '\n'; return;
%     }
%
%   // >= 1 constraint
%   formula << assertion(lor(vars)) << '\n';
%
%   // <= 1 constraint
%   for (auto it1 = vars.begin(); it1 != vars.end(); ++it1)
%     for (auto it2 = it1 + 1; it2 != vars.end(); ++it2)
%       formula << assertion(lor(lnot(*it1), lnot(*it2))) << '\n';
% }
% \end{lstlisting}

\noindent
% We continue with the definition of scheduling constraints, restricting possible transitions from any given state.
% We continue with the definition of scheduling constraints, restricting possible transitions from any given state.
We continue with the cardinality constraint, based on the following functions implementing the different at-most-one predicates.
% Since the cardinality constraint, responsible for non-deterministic scheduling is based on
% If not more than two threads are present, the na\"ive combinatorial constraint is used if

\begin{lstlisting}[style=c++]
void cardinality_naive (const std::vector<std::string> & vars)
{
  for (auto i = vars.begin(); i != vars.end(); ++i)
    for (auto j = i + 1; j != vars.end(); ++j)
      formula << assertion(lor(lnot(*i), lnot(*j))) << '\n';
}
\end{lstlisting}

% \noindent
% Otherwise, if the number of variables is greater than five, Carsten Sinz's sequential counter constraint \cite{ref:Sinz} is used.
% Otherwise, we use Carsten Sinz's sequential counter constraint \cite{ref:Sinz} due to its superiority with respect to the number of clauses.

\begin{lstlisting}[style=c++]
void cardinality_sequential (const std::vector<std::string> & vars)
{
  std::vector<std::string> auxs;
  const auto end = --vars.end();
  for (auto i = vars.begin(); i != end; ++i)
    formula << declare_bool(auxs.emplace_back(*i + "_aux")) << '\n';
  auto var = vars.begin();
  auto aux = auxs.begin();
  formula << assertion(lor(lnot(*var), *aux)) << '\n';
  while (++var != end)
    {
      const std::string & prev = *aux++;
      formula << assertion(lor(lnot(*var), *aux)) << '\n'
              << assertion(lor(lnot(prev), *aux)) << '\n'
              << assertion(lor(lnot(*var), lnot(prev))) << '\n';
    }
  formula << assertion(lor(lnot(*var), lnot(*aux))) << '\n';
}
\end{lstlisting}

% \begin{lstlisting}[style=c++]
% void define_scheduling_constraints (uint k)
% {
%   std::vector<std::string> vars;
%
%   for (uint t = 0; t < programs.size(); t++)
%     {
%       vars.push_back(thread_var(k, t));
%       vars.push_back(flush_var(k, t));
%     }
%
%   if (!halts.empty() || !exits.empty())
%     vars.push_back(exit_flag_var(k));
%
%   formula << (use_sinz_constraint ? card_constraint_sinz(vars)
%                                   : card_constraint_naive(vars))
%           << eol;
% }
% \end{lstlisting}

% vars.insert(vars.end(), {thread_var(k, t), flush_var(k, t)});

\noindent
% adding the conjunction for expressing at-least-one constraint, followed by the .. depending on the number of variables.
The required exactly-one constraint can now be defined by collecting the relevant variables and selecting the appropriate at-most-one predicate after adding a disjunction expressing that at-least-one needs to be true.
% The exactly-one constraint can now be defined by collecting the relevant variables and adding the conjunction for expressing that at least one variable is required to be true, followed by selecting the appropriate at-most-one predicate depending on the number of variables.

\begin{lstlisting}[style=c++]
void define_cardinality_constraints (uint k)
{
  std::vector<std::string> vars;
  for (uint t = 0; t < programs.size(); t++)
    {
      vars.push_back(thread(k, t));
      vars.push_back(flush(k, t));
    }
  vars.push_back(exit(k));
  // >= 1 constraint
  formula << assertion(lor(vars)) << '\n';
  // <= 1 constraint
  if (vars.size() > 5)
    cardinality_sequential(vars);
  else
    cardinality_naive(vars);
}
\end{lstlisting}

\newpage

\noindent
\emph{Example:} \lstCPP{define_cardinality_constraints(0)}

\begin{lstlisting}[style=smtlib]
(assert (or thread_0_0 flush_0_0
            thread_0_1 flush_0_1
            thread_0_2 flush_0_2
            exit_0))

(declare-fun thread_0_0_aux () Bool)
(declare-fun flush_0_0_aux () Bool)
(declare-fun thread_0_1_aux () Bool)
(declare-fun flush_0_1_aux () Bool)
(declare-fun thread_0_2_aux () Bool)
(declare-fun flush_0_2_aux () Bool)

(assert (or (not thread_0_0) thread_0_0_aux))
(assert (or (not flush_0_0) flush_0_0_aux))
(assert (or (not thread_0_0_aux) flush_0_0_aux))
(assert (or (not flush_0_0) (not thread_0_0_aux)))
(assert (or (not thread_0_1) thread_0_1_aux))
(assert (or (not flush_0_0_aux) thread_0_1_aux))
(assert (or (not thread_0_1) (not flush_0_0_aux)))
(assert (or (not flush_0_1) flush_0_1_aux))
(assert (or (not thread_0_1_aux) flush_0_1_aux))
(assert (or (not flush_0_1) (not thread_0_1_aux)))
(assert (or (not thread_0_2) thread_0_2_aux))
(assert (or (not flush_0_1_aux) thread_0_2_aux))
(assert (or (not thread_0_2) (not flush_0_1_aux)))
(assert (or (not flush_0_2) flush_0_2_aux))
(assert (or (not thread_0_2_aux) flush_0_2_aux))
(assert (or (not flush_0_2) (not thread_0_2_aux)))
(assert (or (not exit_0) (not flush_0_2_aux)))
\end{lstlisting}

% \todo[inline]{define store buffer constraints} %%%%%%%%%%%%%%%%%%%%%%%%%%%%%%%%%

\noindent
% Since the store buffer related constraints of a specific thread depend on the presence of barrier instructions,
% if a thread does not contain any barrier instruction and only
% Flushing an empty store buffer or execution of a barrier instruction while it is full is prevented on the basis
% we resort to their programm counters, given in the \lstCPP{barriers} map,
% for creating the list of corresponding statement activation variables to generate the implication preventing execution of a barrier instruction if its store buffer is full.
% Since the store buffer related constraints of each thread depend on the presence of memory barriers,
% Since the store buffer related constraints differ depending on the presence of memory barriers
% If the \lstCPP{barrier} map does not contain an entry for the given thread .
% Otherwise we use the list of program counters for creating the list of corresponding statement activation variables to generate the implication preventing execution of a barrier instruction if its store buffer is full.
% if the thread does not contain any memory barrier, flushing an empty store buffer is prevented by the according implication.
% Otherwise, we use the list of program counters store in the \lstCPP{barriers} map to create the list of corresponding statement activation variables for generating the implication preventing execution of a barrier instruction if its store buffer is full
% Due to the different store buffer constraints depending on the presence of memory barriers, we use the program counters stored in \lstCPP{barriers} to generate the required list of corresponding statement activation variables for delaying their execution while the store buffer is full.
% Depending on the presence of memory barriers, different store buffer constraints are defined
% Since the store buffer related constraints of each thread depend on the presence of memory barriers,
% \lstCPP{barriers} containing the program counters of memory barrier instructions to generate the list of corresponding statement activation variables required by the conjunction for preventing their execution while the store buffer is full.
% Since a thread's store buffer related constraints depend on the presence of memory barriers,
% an empty store buffer is not allowed to be flushed,
%
\CHANGE{Since a program might not even contain a single memory barrier,}
% Depending on the presence of memory barriers,
different store buffer related constraints are added for each thread.
% Otherwise, if the thread's programm contains a \emph{barrier} instruction,
Beside the common restriction of flushing an empty store buffer, %, common for each thread,
% we use the program counters of barrier instructions stored in \lstCPP{barriers} for generating the list of corresponding statement activation variables required to additionally delay their execution while the store buffer is full.
we resort to \lstCPP{barriers} containing the program counters of barrier instructions for generating the list of corresponding statement activation variables required to additionally delay their execution while the store buffer is full.
% we use \lstCPP{barriers} containing the program counters of barrier instructions for generating the list of corresponding statement activation variables required to additionally delay their execution while the store buffer is full.
% we resort to the relevant program counters stored in \lstCPP{barriers} for generating the list of corresponding statement activation variables required to additionally delay their execution while the store buffer is full.

\newpage

\begin{lstlisting}[style=c++]
void define_store_buffer_constraints (uint k)
{
  for (uint t = 0; t < programs.size(); t++)
    if (barriers[t].empty())
      formula << assertion(imply(flush(k, t), full(k, t))) << '\n';
    else
      {
        std::vector<std::string> stmts;
        for (uint pc : barriers[t])
          stmts.push_back(stmt(k, t, pc));
        formula << assertion(ite(full(k, t),
                                 imply(lor(stmts),
                                       lnot(thread(k, t))),
                                 lnot(flush(k, t)))) << '\n';
      }
}
\end{lstlisting}

\noindent
\emph{Example:} \lstCPP{define_store_buffer_constraints(0)}

% \begin{lstlisting}[style=smtlib]
% (assert (ite full_1_0
             % (=> (or stmt_1_0_1 stmt_1_0_4) (not thread_1_0))
             % (not flush_1_0)))
% (assert (ite full_1_1
             % (=> (or stmt_1_1_1 stmt_1_1_4) (not thread_1_1))
             % (not flush_1_1)))
% (assert (=> (not full_1_2) (not flush_1_2)))
% \end{lstlisting}

\begin{lstlisting}[style=smtlib]
(assert (ite full_0_0
             (=> (or stmt_0_0_1 stmt_0_0_4) (not thread_0_0))
             (not flush_0_0)))
(assert (ite full_0_1
             (=> (or stmt_0_1_1 stmt_0_1_4) (not thread_0_1))
             (not flush_0_1)))
(assert (=> flush_0_2 full_0_2))
\end{lstlisting}

% \todo[inline]{define checkpoint constraints} %%%%%%%%%%%%%%%%%%%%%%%%%%%%%%%%%%%

% \begin{lstlisting}[style=c++]
% void define_checkpoint_constraints (uint k)
% {
  % for (const auto & [id, threads] : checkpoints)
    % {
      % for (const auto & [t, _] : threads)
        % {
          % formula
            % << assertion(implication(land(block_var(k, id, t),
                                          % lnot(check_var(k, id))),
                                     % lnot(thread_var(k, t))))
            % << eol;
        % }
      % formula << eol;
    % }
% }
% \end{lstlisting}

\noindent
% Checkpoint constraints are again based on the \tettt{checkpoints} map for generating the relevant $\BLOCK^k^{id, t}$ and $\CHECK^k_{id}$ variables to explicitly disable a thread's activation while it is waiting until all participants reached the checkpoint.
Checkpoint constraints are defined by using \lstCPP{checkpoints} to generate the relevant $\BLOCK^k_{id, t}$ and $\CHECK^k_{id}$ variables for explicitly disabling a thread's activation while it is waiting for all other participants to synchronize upon checkpoint $id$.

\begin{lstlisting}[style=c++]
void define_checkpoint_constraints (uint k)
{
  for (const auto & [id, threads] : checkpoints)
    for (const auto & [t, _] : threads)
      formula << assertion(imply(land(block(k, id, t),
                                      lnot(check(k, id))),
                                 lnot(thread(k, t)))) << '\n';
}
\end{lstlisting}

% \begin{lstlisting}[style=smtlib]
% (assert (=> (and block_1_0_0 (not check_1_0)) (not thread_1_0)))
% (assert (=> (and block_1_0_1 (not check_1_0)) (not thread_1_1)))
% (assert (=> (and block_1_0_2 (not check_1_0)) (not thread_1_2)))
% \end{lstlisting}

\newpage
\noindent
\emph{Example:} \lstCPP{define_checkpoint_constraints(0)}

\begin{lstlisting}[style=smtlib]
(assert (=> (and block_0_0_0 (not check_0_0)) (not thread_0_0)))
(assert (=> (and block_0_0_1 (not check_0_0)) (not thread_0_1)))
(assert (=> (and block_0_0_2 (not check_0_0)) (not thread_0_2)))
\end{lstlisting}

% \todo[inline]{define halt constraints} %%%%%%%%%%%%%%%%%%%%%%%%%%%%%%%%%%%%%%%%%

% \begin{lstlisting}[style=c++]
% void define_halt_constraints (uint k)
% {
  % if (halts.empty())
    % return;
%
  % for (const auto & [t, _] : halts)
    % formula << assertion(implication(halt_var(k, t),
                                     % lnot(thread_var(k, t))))
            % << eol;
%
  % formula << eol;
% }
% \end{lstlisting}

\noindent
% The definition of scheduling constraints is concluded by preventing halted threads from being scheduled.
% The definition of scheduling constraints is concluded by stopping halted threads, based on the entries contained in the \lstCPP{updates} map.
The definition of scheduling constraints is concluded by \CHANGE{preventing halted threads from being executed.}

\begin{lstlisting}[style=c++]
void define_halt_constraints (uint k)
{
  for (uint t = 0; t < programs.size(); t++)
    formula << assertion(imply(halt(k, t),
                               lnot(thread(k, t)))) << '\n';
}
\end{lstlisting}

% \begin{lstlisting}[style=smtlib]
% (assert (=> halt_1_0 (not thread_1_0)))
% (assert (=> halt_1_1 (not thread_1_1)))
% \end{lstlisting}

\noindent
\emph{Example:} \lstCPP{define_halt_constraints(0)}

\begin{lstlisting}[style=smtlib]
(assert (=> halt_0_0 (not thread_0_0)))
(assert (=> halt_0_1 (not thread_0_1)))
(assert (=> halt_0_2 (not thread_0_2)))
\end{lstlisting}

% \todo[inline]{init states} %%%%%%%%%%%%%%%%%%%%%%%%%%%%%%%%%%%%%%%%%%%%%%%%%%%%%

\noindent
Encoding of each particular step ends with the definition of machine states,
% Finally, encoding of each step if finalized by the definition of machine states,
starting with their initialization in the first step $k = 0$.

\begin{lstlisting}[style=c++]
void init_states ()
{
  init_accu();
  init_mem();
  init_adr();
  init_val();
  init_full();
  init_stmt();
  init_block();
  init_halt();
  init_heap();
  init_exit();
}
\end{lstlisting}

% \begin{lstlisting}[style=smtlib]
% \end{lstlisting}

% \todo[inline]{init accu} %%%%%%%%%%%%%%%%%%%%%%%%%%%%%%%%%%%%%%%%%%%%%%%%%%%%%%%

\noindent
% Similar to their declaration, we start by requiring all initial accumulator register states to be zero.
We begin by setting each thread's initial accumulator register state to zero.

\begin{lstlisting}[style=c++]
void init_accu ()
{
  for (uint t = 0; t < programs.size(); t++)
    formula << assign(accu(0, t), consth(0)) << '\n';
}
\end{lstlisting}

\noindent
\emph{Example:} \lstCPP{init_accu()}

\begin{lstlisting}[style=smtlib]
(assert (= accu_0_0 #x0000))
(assert (= accu_0_1 #x0000))
(assert (= accu_0_2 #x0000))
\end{lstlisting}

% \todo[inline]{init mem} %%%%%%%%%%%%%%%%%%%%%%%%%%%%%%%%%%%%%%%%%%%%%%%%%%%%%%%%

% \begin{lstlisting}[style=c++]
% void init_mem ()
% {
  % for (uint t = 0; t < programs.size(); t++)
    % formula << assign(mem_var(k, t), to_hex(0)) << '\n';
% }
% \end{lstlisting}

% \begin{lstlisting}[style=smtlib]
% ; mem variables - mem_<step>_<thread>
% (assert (= mem_0_0 #x0000))
% (assert (= mem_0_1 #x0000))
% (assert (= mem_0_2 #x0000))
% \end{lstlisting}

% \todo[inline]{init adr} %%%%%%%%%%%%%%%%%%%%%%%%%%%%%%%%%%%%%%%%%%%%%%%%%%%%%

% \begin{lstlisting}[style=c++]
% void init_adr ()
% {
  % for (uint t = 0; t < programs.size(); t++)
    % formula << assign(adr_var(k, t), to_hex(0)) << eol;
  % formula << eol;
% }
% \end{lstlisting}
%
% \begin{lstlisting}[style=smtlib]
% ; store buffer address variables - adr_<step>_<thread>
% (assert (= adr_0_0 #x0000))
% (assert (= adr_0_1 #x0000))
% (assert (= adr_0_2 #x0000))
% \end{lstlisting}

% \todo[inline]{init val} %%%%%%%%%%%%%%%%%%%%%%%%%%%%%%%%%%%%%%%%%%%%%%%%%%%%%

% \begin{lstlisting}[style=c++]
% void init_val ()
% {
  % for (uint t = 0; t < programs.size(); t++)
    % formula << assign(val_var(k, t), to_hex(0)) << eol;
  % formula << eol;
% }
% \end{lstlisting}

% \begin{lstlisting}[style=smtlib]
% ; store buffer value variables - val_<step>_<thread>
% (assert (= val_0_0 #x0000))
% (assert (= val_0_1 #x0000))
% (assert (= val_0_2 #x0000))
% \end{lstlisting}

% \todo[inline]{init full} %%%%%%%%%%%%%%%%%%%%%%%%%%%%%%%%%%%%%%%%%%%%%%%%%%%%

% \begin{lstlisting}[style=c++]
% void init_full ()
% {
  % for (uint t = 0; t < programs.size(); t++)
    % formula << assertion(lnot(full_var(k, t))) << eol;
  % formula << eol;
% }
% \end{lstlisting}

% \begin{lstlisting}[style=smtlib]
% ; store buffer full variables - full_<step>_<thread>
% (assert (not full_0_0))
% (assert (not full_0_1))
% (assert (not full_0_2))
% \end{lstlisting}

% \todo[inline]{init stmt} %%%%%%%%%%%%%%%%%%%%%%%%%%%%%%%%%%%%%%%%%%%%%%%%%%%%%%%
% \hline
% ~\\
% \noindent
% \texttt{init_stmt}
% \vspace{.3\baselineskip}
% \hline
% ~\\
\noindent
Similar to their declaration, we skip the initialization functions of the remaining register states as they are again defined almost identical and continue with the activation of every thread's initial statement $\STMT^0_{t, 0}$.
% Similar to their declaration, we skip the initialization functions of the remaining register states as they are again defined almost identical and continue by activating the initial statement $\STMT^0_{t, 0}$ of every thread $t$.
% \\
% \hline
\begin{lstlisting}[style=c++]
void init_stmt ()
{
  for (uint t = 0; t < programs.size(); t++)
    for (uint pc = 0; t < programs[t].size(); t++)
      formula << assertion(pc ? lnot(stmt(0, t, pc))
                              : stmt(0, t, pc)) << '\n';
}
\end{lstlisting}

\noindent
\emph{Example:} \lstCPP{init_stmt()}

\begin{lstlisting}[style=smtlib]
(assert stmt_0_0_0)
(assert (not stmt_0_0_1))
(assert (not stmt_0_0_2))
(assert (not stmt_0_0_3))
(assert (not stmt_0_0_4))

(assert stmt_0_1_0)
(assert (not stmt_0_1_1))
(assert (not stmt_0_1_2))
(assert (not stmt_0_1_3))
(assert (not stmt_0_1_4))

(assert stmt_0_2_0)
(assert (not stmt_0_2_1))
(assert (not stmt_0_2_2))
(assert (not stmt_0_2_3))
(assert (not stmt_0_2_4))
(assert (not stmt_0_2_5))
\end{lstlisting}
% \hline

% \todo[inline]{init block} %%%%%%%%%%%%%%%%%%%%%%%%%%%%%%%%%%%%%%%%%%%%%%%%%%%%%%
% ~\\
\noindent
% Initialization of block flags $\BLOCK^k_{0, t}$ is uses the
% Block flags $\BLOCK^0_{id, t}$ are initially disabled on the basis of the \lstCPP{checkpoints} map.
% Block flags $\BLOCK^0_{id, t}$ of every checkpoint $id$ and thread $t$ are initially disabled by simply asserting their negation.
% Block flags $\BLOCK^0_{id, t}$ for every checkpoint $id$ and thread $t$ are initially disabled by iterating the corresponding \lstCPP{checkpoints} map.
% Block flags $\BLOCK^0_{id, t}$ are initially disabled by
% Scheduling related thread states $\BLOCK^0_{id, t}$, $\HALT^0_t$ and $\EXIT_0$ are initially disabled.
% Block flags $\BLOCK^0_{id, t}$ are initially disabled based on the entries given in the \lstCPP{checkpoints} map.
Block flags $\BLOCK^0_{id, t}$ are initially disabled by relying on the entries given in the \lstCPP{checkpoints} map.
% Scheduling related thread states $\BLOCK^0_{id, t}$ and $\HALT^0_t$ are initialized by iterating the corresponding utility variables.

\newpage

\begin{lstlisting}[style=c++]
void init_block ()
{
  for (const auto & [id, threads] : checkpoints)
    for (const auto & [t, _] : threads)
      formula << assertion(lnot(block(0, id, t))) << '\n';
}
\end{lstlisting}

\noindent
\emph{Example:} \lstCPP{init_block()}

\begin{lstlisting}[style=smtlib]
(assert (not block_0_0_0))
(assert (not block_0_0_1))
(assert (not block_0_0_2))
\end{lstlisting}

% \todo[inline]{init halt} %%%%%%%%%%%%%%%%%%%%%%%%%%%%%%%%%%%%%%%%%%%%%%%%%%%%%%%

\noindent
% Halt flags $\HALT^0_t$ are initially disabled on the basis of the corresponding entry in the \lstCPP{updates} map.
% Halt flags $\HALT^0_t$ are initialized in a similar manner, using the relevant entry in the \lstCPP{updates} map.
Halt flags $\HALT^0_t$ \CHANGE{of each thread are initialized in a similar manner.}

\begin{lstlisting}[style=c++]
void init_halt ()
{
  for (uint t = 0; t < programs.size(); t++)
    formula << assertion(lnot(halt(0, t))) << '\n';
}
\end{lstlisting}

\noindent
\emph{Example:} \lstCPP{init_halt()}

\begin{lstlisting}[style=smtlib]
(assert (not halt_0_0))
(assert (not halt_0_1))
(assert (not halt_0_2))
\end{lstlisting}

% \todo[inline]{init heap} %%%%%%%%%%%%%%%%%%%%%%%%%%%%%%%%%%%%%%%%%%%%%%%%%%%%%%%

\noindent
% Although the machine's memory is considered to be uninitialized in general, relevant cells are initialized according to the presence of a memory map.
Although the machine's memory is considered to be uninitialized in general, input data is assigned according to the entries of a given memory map.

\begin{lstlisting}[style=c++]
void init_heap ()
{
  for (const auto & [adr, val] : mmap)
    formula << assign(select(heap(0), consth(adr)),
                      consth(val)) << '\n';
}
\end{lstlisting}

\noindent
\emph{Example:} \lstCPP{init_heap()}

\begin{lstlisting}[style=smtlib]
(assert (= (select heap_0 #x0000) #x0000))
(assert (= (select heap_0 #x0001) #x0000))
\end{lstlisting}

% \todo[inline]{init exit-flag} %%%%%%%%%%%%%%%%%%%%%%%%%%%%%%%%%%%%%%%%%%%%%%%%%%

\noindent
Finally, the initial exit flag $\EXIT_0$ is disabled to enforce execution.
% Finally, execution of is enforced by disabling the initial exit flag $\EXIT_0$.

\begin{lstlisting}[style=c++]
void init_exit()
{
  formula << assertion(lnot(exit(0))) << '\n';
}
\end{lstlisting}

\noindent
\emph{Example:} \lstCPP{init_exit()}

\begin{lstlisting}[style=smtlib]
(assert (not exit_0))
\end{lstlisting}

% \todo[inline]{encoding instructions} %%%%%%%%%%%%%%%%%%%%%%%%%%%%%%%%%%%%%%%%%%%

% \todo[inline]{encode LOAD} %%%%%%%%%%%%%%%%%%%%%%%%%%%%%%%%%%%%%%%%%%%%%%%%%%%%%

\noindent
% Memory access is implemented by the following helper according to the previously defined $\LOAD^k_t$ function.
Memory access according to the previously defined $\LOAD^k_t$ is implemented by the following helper function.

\begin{lstlisting}[style=c++]
std::string load (uint k, uint t, uint adr, bool indirect = false)
{
  std::string address = consth(adr);
  std::string adr_var = adr(k, t);
  std::string val_var = val(k, t);
  std::string full_var = full(k, t);
  std::string heap_var = heap(k);

  if (indirect)
    return
      ite(full_var,
          ite(equal(adr_var, address),
              ite(equal(val_var, address),
                  val_var,
                  select(heap_var, val_var)),
              ite(equal(adr_var, select(heap_var, address)),
                  val_var,
                  select(heap_var, select(heap_var, address)))),
          select(heap_var, select(heap_var, address)));
  else
    return
      ite(land(full_var, equal(adr_var, address)),
          val_var,
          select(heap_var, address));
}
\end{lstlisting}

\noindent
% Definition of common encoding functions is concluded by introducing the implementation of virtual member function \lstCPP{Instruction::encode(uint, uint, State)} for all instructions, based on the frame axioms.
% Common encoding functions are concluded by introducing the virtual member functions implementing \lstinline[style=c++, basicstyle=\normalsize\ttfamily]{Instruction::encode(uint, uint, State)}, returning the successor of a given state after executing the specific instruction, expect for the statement activation variables $\STMT^k_{t, pc}$, which are handeled seperately.
% Common encoding functions are concluded by implementing the virtual member function \lstinline[style=c++, basicstyle=\normalsize\ttfamily]{Instruction::encode(uint, uint, State)} of all descendants, returning the successor of a given state after executing the specific instruction, expect for the statement activation variables $\STMT^k_{t, pc}$, which are handeled seperately.
Common encoding functions are concluded by introducing implementations for \lstCPP{Instruction::encode(uint, uint, State)}, returning the successor of a given state after executing the specific instruction expect for the statement activation variables $\STMT^k_{t, pc}$, which are handeled seperately and only jump conditions being generated this way.
% Statement activation variables $\STMT^k_{t, pc}$ are handeled seperately and only the positive
Implementations for \lstASM{FENCE}, \lstASM{JMP} and \lstASM{HALT} instructions are omitted, as they just return an empty \lstCPP{std::string}.

% \newcommand{\defencode}[1]{
  % \par\noindent
  % #1
  % \par\noindent
  % \rule[0.5\baselineskip]{\textwidth}{0.1pt}\vspace{-1.3\baselineskip}%\par%\noindent
  % % #2
% }

\newcommand{\defencode}[1]{
  \vfill
  \par\noindent
  % {\small #1}
  #1
  \par\noindent
  \rule[0.5\baselineskip]{\textwidth}{0.1pt}%\par%\noindent
  % #2
}

\newlength{\encodeskip}
\setlength{\encodeskip}{.5\baselineskip}

\lstdefinestyle{encode}{aboveskip=-\encodeskip, belowskip=\encodeskip}

% \newcommand{\defencode}[1]{
  % \par\noindent
  % #1
  % \par\noindent\makebox[\textwidth]{\rule[.5\baselineskip]{\textwidth}{0.1pt}}\par\noindent
  % \vspace{-2.4\baselineskip}
  % % \rule[0.5\baselineskip]{\textwidth}{0.1pt}\vspace{-1.3\baselineskip}%\par%\noindent
  % % #2
% }

% \lstset{belowskip=\baselineskip}

\defencode{\lstASM{LOAD arg}}
% \begin{lstlisting}[style=c++,aboveskip=-\encodeskip,belowskip=\encodeskip]
\begin{lstlisting}[style=c++, style=encode]
std::string Load::encode (uint k, uint t, State state)
{
  return load(k, t, arg, indirect);
}
\end{lstlisting}

% \todo[inline]{encode STORE} %%%%%%%%%%%%%%%%%%%%%%%%%%%%%%%%%%%%%%%%%%%%%%%%%%%%

\defencode{\lstASM{STORE arg}}
\begin{lstlisting}[style=c++, style=encode]
std::string Store::encode (uint k, uint t, State state)
{
  switch (state)
    {
    case State::adr: return indirect ? load(k, arg) : consth(arg);
    case State::val: return accu(k, t);
    }
}
\end{lstlisting}

% \todo[inline]{encode ADD} %%%%%%%%%%%%%%%%%%%%%%%%%%%%%%%%%%%%%%%%%%%%%%%%%%%%%%

\defencode{\lstASM{ADD arg}}
\begin{lstlisting}[style=c++, style=encode]
std::string Add::encode (uint k, uint t, State state)
{
  return bvadd(accu(k, t), load(k, arg, indirect));
}
\end{lstlisting}

% \todo[inline]{encode ADDI} %%%%%%%%%%%%%%%%%%%%%%%%%%%%%%%%%%%%%%%%%%%%%%%%%%%%%

\defencode{\lstASM{ADDI arg}}
\begin{lstlisting}[style=c++, style=encode]
std::string Addi::encode (uint k, uint t, State state)
{
  return bvadd(accu(k, t), consth(arg));
}
\end{lstlisting}

% \todo[inline]{encode SUB} %%%%%%%%%%%%%%%%%%%%%%%%%%%%%%%%%%%%%%%%%%%%%%%%%%%%%%

\defencode{\lstASM{SUB arg}}
\begin{lstlisting}[style=c++, style=encode]
std::string Sub::encode (uint k, uint t, State state)
{
  return bvsub(accu(k, t), load(k, arg, indirect));
}
\end{lstlisting}

% \todo[inline]{encode SUBI} %%%%%%%%%%%%%%%%%%%%%%%%%%%%%%%%%%%%%%%%%%%%%%%%%%%%%

\defencode{\lstASM{SUBI arg}}
\begin{lstlisting}[style=c++, style=encode]
std::string Subi::encode (uint k, uint t, State state)
{
  return bvsub(accu(k, t), consth(arg));
}
\end{lstlisting}

% \todo[inline]{encode MUL} %%%%%%%%%%%%%%%%%%%%%%%%%%%%%%%%%%%%%%%%%%%%%%%%%%%%%%

\newpage

\defencode{\lstASM{MUL arg}}
\begin{lstlisting}[style=c++, style=encode]
std::string Mul::encode (uint k, uint t, State state)
{
  return bvmul(accu(k, t), load(k, arg, indirect));
}
\end{lstlisting}

% \todo[inline]{encode MULI} %%%%%%%%%%%%%%%%%%%%%%%%%%%%%%%%%%%%%%%%%%%%%%%%%%%%%

\defencode{\lstASM{MULI arg}}
\begin{lstlisting}[style=c++, style=encode]
std::string Muli::encode (uint k, uint t, State state)
{
  return bvmul(accu(k, t), consth(arg));
}
\end{lstlisting}

% \todo[inline]{encode CMP} %%%%%%%%%%%%%%%%%%%%%%%%%%%%%%%%%%%%%%%%%%%%%%%%%%%%%%

\defencode{\lstASM{CMP arg}}
\begin{lstlisting}[style=c++, style=encode]
std::string Cmp::encode (uint k, uint t, State state)
{
  return bvsub(accu(k, t), load(k, arg, indirect));
}
\end{lstlisting}

% \todo[inline]{encode JMP} %%%%%%%%%%%%%%%%%%%%%%%%%%%%%%%%%%%%%%%%%%%%%%%%%%%%%%

% \defencode{\lstASM{JMP arg}}
% \begin{lstlisting}[style=c++, style=encode]
% std::string Jmp::encode (uint k, uint t, State state) { return ""; }
% \end{lstlisting}

% \todo[inline]{encode JZ} %%%%%%%%%%%%%%%%%%%%%%%%%%%%%%%%%%%%%%%%%%%%%%%%%%%%%%%

\defencode{\lstASM{JZ arg}}
\begin{lstlisting}[style=c++, style=encode]
std::string Jz::encode (uint k, uint t, State state)
{
  return equal(accu(k, t), consth(0));
}
\end{lstlisting}

% \todo[inline]{encode JNZ} %%%%%%%%%%%%%%%%%%%%%%%%%%%%%%%%%%%%%%%%%%%%%%%%%%%%%%

\defencode{\lstASM{JNZ arg}}
\begin{lstlisting}[style=c++, style=encode]
std::string Jnz::encode (uint k, uint t, State state)
{
  return lnot(equal(accu(k, t), consth(0)));
}
\end{lstlisting}

% \todo[inline]{encode JS} %%%%%%%%%%%%%%%%%%%%%%%%%%%%%%%%%%%%%%%%%%%%%%%%%%%%%%%

% \todo[noline]{instroduce global std::string msb = "15"?}
\defencode{\lstASM{JS arg}}
\begin{lstlisting}[style=c++, style=encode]
std::string Js::encode (uint k, uint t, State state)
{
  return equal("#b1", extract("15", "15", accu(k, t)));
}
\end{lstlisting}

% \todo[inline]{encode JNS} %%%%%%%%%%%%%%%%%%%%%%%%%%%%%%%%%%%%%%%%%%%%%%%%%%%%%%

\newpage

\defencode{\lstASM{JNS arg}}
\begin{lstlisting}[style=c++, style=encode]
std::string Jns::encode (uint k, uint t, State state)
{
  return equal("#b0", extract("15", "15", accu(k, t)));
}
\end{lstlisting}

% \todo[inline]{encode JNZNS} %%%%%%%%%%%%%%%%%%%%%%%%%%%%%%%%%%%%%%%%%%%%%%%%%%%%

\defencode{\lstASM{JNZNS arg}}
\begin{lstlisting}[style=c++, style=encode]
std::string Jnzns::encode (uint k, uint t, State state)
{
  return land(lnot(equal(accu(k, t), consth(0))),
              equal("#b0", extract("15", "15", accu(k, t))));
}
\end{lstlisting}

% \todo[inline]{encode MEM} %%%%%%%%%%%%%%%%%%%%%%%%%%%%%%%%%%%%%%%%%%%%%%%%%%%%%%

\defencode{\lstASM{MEM arg}}
\begin{lstlisting}[style=c++, style=encode]
std::string Mem::encode (uint k, uint t, State state)
{
  return load(k, arg, indirect);
}
\end{lstlisting}

% \todo[inline]{encode CAS} %%%%%%%%%%%%%%%%%%%%%%%%%%%%%%%%%%%%%%%%%%%%%%%%%%%%%%

\defencode{\lstASM{CAS arg}}
\begin{lstlisting}[style=c++, style=encode]
std::string Cas::encode (uint k, uint t, State state)
{
  auto heap_var = heap(k);
  auto address = indirect ? select(heap_var, consth(arg))
                          : consth(arg);
  auto condition = equal(mem(k, t), select(heap_var, address));
  switch (state)
    {
    case State::accu: return ite(condition, consth(1), consth(0));
    case State::heap: return ite(condition,
                                store(heap_var,
                                      address,
                                      accu(k, t)),
                                heap_var);
    }
}
\end{lstlisting}

% \todo[inline]{encode HALT} %%%%%%%%%%%%%%%%%%%%%%%%%%%%%%%%%%%%%%%%%%%%%%%%%%%

% \defencode{\lstASM{HALT}}
% \begin{lstlisting}[style=c++, style=encode]
% std::string Halt::encode (uint k, uint t, State state) { return ""; }
% \end{lstlisting}

% \todo[inline]{encode EXIT} %%%%%%%%%%%%%%%%%%%%%%%%%%%%%%%%%%%%%%%%%%%%%%%%%%%%%

\defencode{\lstASM{EXIT arg}}
\begin{lstlisting}[style=c++, style=encode]
std::string Exit::encode (uint k, uint t, State state)
{
  return consth(arg);
}
\end{lstlisting}

% \newpage


\newpage

\section{Functional Next State Logic}

% \todo[inline]{define states (functional)} %%%%%%%%%%%%%%%%%%%%%%%%%%%%%%%%%%%%%%%

\CHANGE{
In principle there are two ways to encode state updates: a functional, relying on \lstCPP{ite} cascades to perform \emph{static single assignments} for determining the successor of each state explicitly and a relational, based on implying the next state for each possible transition.
While the relational version is easier to encode since it is closer to our model's semantics and results in a simpler structure, the functional encoding has the benefit of being more compact and also turned out to be more efficient for SMT solving. % due to its strictly hierarchical nature.
% , which can be radically simplified by clever term rewriting due to its strictly hierarchical nature. %structure.
}%
% deterministically
% functional encoding scheme, static single assignment to determine the successor of each state individually
% functional encoding scheme, based on static single assignments for determining the successor of each state individually.
% We will now introduce the main state update scheme according to our model's frame axioms, generated using the \texttt{-e smtlib} command line parameter.
We will now introduce our \CHANGE{functional} state update scheme, generated by using the \mbox{\texttt{-e smtlib}} command line parameter.
% This functional encoding relies on \lstCPP{ite} cascades to perform \emph{static single assignments} for determining the successor of each state individually and aids the solver by simplifying the deduction process.
% relying on \lstCPP{ite} cascades to perform \emph{static single assignments} for determining the successor of each state explicitly in a functional manner.% and aids the solver by simplifying the deduction process.
% in order to aid the solver by simplifying the deduction process and minimizing the size of the generated formula.

\begin{lstlisting}[style=c++]
void define_states (uint k)
{
  define_accu(k);
  define_mem(k);
  define_adr(k);
  define_val(k);
  define_full(k);
  define_stmt(k);
  define_block(k);
  define_halt(k);
  define_heap(k);
  define_exit(k);

  if (k == bound)
    define_exit_code();
}
\end{lstlisting}

% \todo[inline]{define accu} %%%%%%%%%%%%%%%%%%%%%%%%%%%%%%%%%%%%%%%%%%%%%%%%%%%%%

\noindent
% \begin{itemize}
  % \item recursively
  % \item starts with base case of preserving its value
  % \item for each statement altering $\ACCU^k_t$
% \end{itemize}
% starting with the preservation of it's previous state $\ACCU^{k - 1}_t$, the
% starting with the base case of preserving it's previous state $\ACCU^{k - 1}_t$, the
% State update expressions for every thread's accumulator register $\ACCU^k_t$ are initialized with the predecessor $\ACCU^{k - 1}_t$
% Starting with the accumulator registers $\ACCU^k_t$, each update expression is generated by adding an \lstCPP{ite} for every program statement altering the accumulator
% Starting with the accumulator registers $\ACCU^k_t$, each thread's update expression is initialized with it's predecessor $\ACCU^{k - 1}_t$, forming the base case of preserving the previous state.
% Each program statement altering the accumulator then adds an \lstCPP{ite}, returning the successor generated by the \lstinline[style=c++, basicstyle=\normalsize\ttfamily]{Instruction::encode(uint, uint, State)} function
% Starting with the accumulator registers $\ACCU^k_t$, each thread's update expression is built by introducing an \lstCPP{ite} for every program statement altering the accumulator, returning the corresponding successor according to the previously defined implementations of \lstinline[style=c++, basicstyle=\normalsize\ttfamily]{Instruction::encode(uint, uint, State)}.
% Starting with the base case of preserving the previous value, each accumulator register's state update expression is initialized with it's predecessor $\ACCU^{k - 1}_t$.
% Every program statement altering the accumulator then adds an \lstCPP{ite}, returning the according successor state depending on it's execution given by the corresponding execution variable $\EXEC^{k - 1}_{t, pc}$ and...
% >>>
% Starting with the accumulator registers $\ACCU^k_t$, each update expression is initialized with it's predecessor $\ACCU^{k - 1}_t$, forming the base case of preserving the previous state.
% >>>
% Every program statement altering the accumulator then adds an \lstCPP{ite} for choosing the successor given by the instruction's \lstinline[style=c++, basicstyle=\normalsize\ttfamily]{Instruction::encode(uint, uint, State)} depending on the corresponding execution variable $\EXEC^{k - 1}_{t, pc}$,   , embedding the current update expression in the \emph{else} branch.
% >>>
% We then add an \lstCPP{ite} for every program statement altering the accumulator to choose the successor state depending on the corresponding execution variable $\EXEC^{k - 1}_{t, pc}$ %, given by the instruction's \lstinline[style=c++, basicstyle=\normalsize\ttfamily]{Instruction::encode(uint, uint, State)}
% , utilizing the \lstCPP{Instruction} class's dynamic dispatch mechanism
% and embedding the current update expression in the \emph{else} branch as a fall back in case the statement has not been executed.
% >>>
% ~\\
% ~\\
% $\ACCU^{k - 1}_t$
Starting with the accumulator registers $\ACCU^k_t$, each update expression is initialized with its predecessor $\ACCU^{k - 1}_t$, forming the base case of preserving the previous state.
% The final expression assigned to $\ACCU^k_t$ is then built by adding an $\ITE$ for every program statement altering the accumulator, choosing the successor state depending on the corresponding execution variable $\EXEC^{k - 1}_{t, pc}$.
% >>>
% Starting with the accumulator registers $\ACCU^k_t$, each update expression is initialized with it's predecessor to preserve the previous state.
% Any program statement altering the accumulator then embeds the expression's current value in the \emph{else} branch of an $\ITE$ setting the successor state by encoding the corresponding instruction, depending on the execution variable $\EXEC^{k - 1}_{t, pc}$.
% Any program statement altering the accumulator then extends the update expression by embedding it's current value in the \emph{else} branch of an $\ITE$, setting the successor state according to the encoded instruction, depending on the corresponding execution variable $\EXEC^{k - 1}_{t, pc}$.
% Any program statement altering the accumulator then extends the expression by embedding it's current value in the \emph{else} branch of an $\ITE$, using the \lstCPP{encode} implementation of the corresponding instruction to set the successor state depending on the execution variable $\EXEC^{k - 1}_{t, pc}$.
Any program statement altering the accumulator then extends the expression by embedding its current value in the \emph{else} branch of an $\ITE$, using the corresponding instruction's \lstCPP{encode} implementation to set the successor state
depending on the
% according to it's
execution variable $\EXEC^{k - 1}_{t, pc}$.
% execution signalled by $\EXEC^{k - 1}_{t, pc}$.
% in case it has been executed signalled by $\EXEC^{k - 1}_{t, pc}$.
% in case of it's execution signalled by $\EXEC^{k - 1}_{t, pc}$.
% >>>
% To increase readability of the generated \SMTLIB expressions, we use a reverse iterator to create the $\ite$ cascade in program order.

% \begin{lstlisting}[style=c++]
% void define_accu (uint k)
% {
  % for (uint t = 0; t < programs.size(); t++)
    % {
      % std::string expr = accu_var(k - 1, t);
      % const auto & stmts = updates[State::accu][t];
      % for (auto pc = stmts.rbegin(); pc != stmts.rend(); ++pc)
        % expr = ite(exec_var(k - 1, t, *pc),
                   % program[t][*pc]->encode(k - 1, t, State::accu),
                   % expr);
      % formula << assign(accu_var(k, t), expr) << '\n';
    % }
% }
% \end{lstlisting}
%
% \noindent
% \emph{Example:} \lstCPP{define_accu(1)}
%
% \begin{lstlisting}[style=smtlib]
% (assert (= accu_1_0
           % (ite exec_0_0_0
                % (bvadd accu_0_0 #x0001)
                % (ite exec_0_0_2
                     % (ite (and full_0_0 (= adr_0_0 #x0001))
                          % val_0_0
                          % (select heap_0 #x0001))
                     % accu_0_0))))
% (assert (= accu_1_1
           % (ite exec_0_1_0
                % (bvadd accu_0_1 #x0001)
                % (ite exec_0_1_2
                     % (ite (and full_0_1 (= adr_0_1 #x0000))
                          % val_0_1
                          % (select heap_0 #x0000))
                     % accu_0_1))))
% (assert (= accu_1_2
           % (ite exec_0_2_1
                % (bvadd accu_0_2
                       % (ite (and full_0_2 (= adr_0_2 #x0000))
                            % val_0_2
                            % (select heap_0 #x0000)))
                % (ite exec_0_2_2
                     % (bvadd accu_0_2
                            % (ite (and full_0_2
                                      % (= adr_0_2 #x0001))
                                 % val_0_2
                                 % (select heap_0 #x0001)))
                     % accu_0_2))))
% \end{lstlisting}

\newpage

\begin{lstlisting}[style=c++]
void define_accu (uint k)
{
  for (uint t = 0; t < programs.size(); t++)
    {
      std::string next = accu(k - 1, t);
      const auto & stmts = updates[State::accu][t];
      for (auto pc = stmts.rbegin(); pc != stmts.rend(); ++pc)
        next = ite(exec(k - 1, t, *pc),
                   program[t][*pc]->encode(k - 1, t, State::accu),
                   next);
      formula << assign(accu(k, t), next) << '\n';
    }
}
\end{lstlisting}

\noindent
\emph{Example:} \lstCPP{define_accu(1)}

\begin{lstlisting}[style=smtlib]
(assert (= accu_1_0
           (ite exec_0_0_0
                (bvadd accu_0_0 #x0001)
                (ite exec_0_0_2
                     (ite (and full_0_0 (= adr_0_0 #x0001))
                          val_0_0
                          (select heap_0 #x0001))
                     accu_0_0))))
(assert (= accu_1_1 ... ; same as before but reading address 0
(assert (= accu_1_2
           (ite exec_0_2_1
                (bvadd accu_0_2
                       (ite (and full_0_2 (= adr_0_2 #x0000))
                            val_0_2
                            (select heap_0 #x0000)))
                (ite exec_0_2_2
                     (bvadd accu_0_2
                            (ite (and full_0_2
                                      (= adr_0_2 #x0001))
                                 val_0_2
                                 (select heap_0 #x0001)))
                     accu_0_2))))
\end{lstlisting}

% \todo[inline]{define mem} %%%%%%%%%%%%%%%%%%%%%%%%%%%%%%%%%%%%%%%%%%%%%%%%%%%%%%

\noindent
% State update functions of the remaining register states $\MEM^k_t$, $\SBADR^k_t$ and $\SBVAL^k_t$ are again defined almost identically and only differ by using the appropriate \lstCPP{State} entries and variable generators.
We omit the implementation of state update functions for the remaining register states $\MEM^k_t$, $\SBADR^k_t$ and $\SBVAL^k_t$ as they only differ \CHANGE{in} using the appropriate \lstCPP{State} entries and variable generators, but include their output \CHANGE{to show} the state updates generated for our demo example.

% \begin{lstlisting}[style=c++]
% void define_mem (uint k)
% {
  % for (uint t = 0; t < programs.size(); t++)
    % {
      % std::string val = mem_var(k - 1, t);
      % const auto & stmts = updates[State::mem][t];
      % for (auto pc = stmts.rbegin(); pc != stmts.rend(); ++pc)
        % val = ite(exec_var(k - 1, t, *pc),
                  % program[t][*pc].encode(k - 1, t, State::mem),
                  % val);
      % formula << assign(mem_var(k, t), val) << '\n';
    % }
% }
% \end{lstlisting}

~\\
\noindent
\emph{Example:} \lstCPP{define_mem(1)}

\vspace{-\baselineskip/2}

\begin{lstlisting}[style=smtlib]
(assert (= mem_1_0 mem_0_0))
(assert (= mem_1_1 mem_0_1))
(assert (= mem_1_2 mem_0_2))
\end{lstlisting}

% \todo[inline]{define adr} %%%%%%%%%%%%%%%%%%%%%%%%%%%%%%%%%%%%%%%%%%%%%%%%%%%

% \begin{lstlisting}[style=c++]
% void define_adr (uint k)
% {
  % for (const auto & [t, stmts] : alters[State::adr])
    % {
      % std::string val = adr_var(k - 1, t);
      % for (auto pc = stmts.rbegin(); pc != stmts.rend(); ++pc)
        % val = ite(exec_var(k - 1, t, *pc),
                  % program[t][*pc].encode(k, t, State::adr),
                  % val);
      % formula << assign(adr_var(k, t), val) << eol;
    % }
  % formula << eol;
% }
% \end{lstlisting}

\noindent
\emph{Example:} \lstCPP{define_adr(1)}

\vspace{-\baselineskip/2}

\begin{lstlisting}[style=smtlib]
(assert (= adr_1_0 (ite exec_0_0_1 #x0000 adr_0_0)))
(assert (= adr_1_1 (ite exec_0_1_1 #x0001 adr_0_1)))
(assert (= adr_1_2 adr_0_2))
\end{lstlisting}

% \todo[inline]{define val} %%%%%%%%%%%%%%%%%%%%%%%%%%%%%%%%%%%%%%%%%%%%%%%%%%%

% \begin{lstlisting}[style=c++]
% void define_val (uint k)
% {
  % for (const auto & [t, stmts] : alters[State::val])
    % {
      % std::string val = val_var(k - 1, t);
      % for (auto pc = stmts.rbegin(); pc != stmts.rend(); ++pc)
        % val = ite(exec_var(k - 1, t, *pc),
                  % program[t][*pc].encode(k, t, State::val),
                  % val);
      % formula << assign(val_var(k, t), val) << eol;
    % }
  % formula << eol;
% }
% \end{lstlisting}

\noindent
\emph{Example:} \lstCPP{define_val(1)}

\vspace{-\baselineskip/2}

\begin{lstlisting}[style=smtlib]
(assert (= val_1_0 (ite exec_0_0_1 accu_0_0 val_0_0)))
(assert (= val_1_1 (ite exec_0_1_1 accu_0_1 val_0_1)))
(assert (= val_1_2 val_0_2))
\end{lstlisting}

% \todo[inline]{define full} %%%%%%%%%%%%%%%%%%%%%%%%%%%%%%%%%%%%%%%%%%%%%%%%%%

\noindent
% collection the execution variables $\EXEC^{k - 1}_{t, pc}$ of \lstASM{STORE} instructions and the predecessor $\SBFULL^{k - 1}_t$ to preserve it's state.
% Successors of store buffer full flags are defined by a single
Store buffer full flags $\SBFULL^k_t$ are defined by a single
% Store buffer full states $\SBFULL^k_t$ are defined by a single
% $\ITE$, either resetting the state in case the store buffer has been flushed,
$\ITE$, either falsifying the state in case the store buffer has been flushed,
% or setting it by building a conjunction over the previously collected variables.
or assigning a conjunction over all execution variables $\EXEC^{k - 1}_{t, pc}$ of \lstASM{STORE} instructions to set it if an entry was added, as well as the predecessor $\SBFULL^{k - 1}_t$ for preserving its state otherwise.

\begin{lstlisting}[style=c++]
void define_full (uint k)
{
  for (uint t = 0; t < programs.size(); t++)
    {
      std::vector<std::string> args {full(k - 1, t)};
      for (uint pc : updates[State::full][t])
        args.push_back(exec(k - 1, t, pc));
      formula << assign(full(k, t),
                        ite(flush(k - 1, t),
                            "false",
                            lor(args))) << '\n';
    }
}
\end{lstlisting}

\noindent
\emph{Example:} \lstCPP{define_full(1)}

\begin{lstlisting}[style=smtlib]
(assert (= full_1_0 (ite flush_0_0
                            false
                            (or full_0_0 exec_0_0_1))))
(assert (= full_1_1 (ite flush_0_1
                            false
                            (or full_0_1 exec_0_1_1))))
(assert (= full_1_2 (ite flush_0_2 false full_0_2)))
\end{lstlisting}

\noindent
% \todo[inline]{define stmt} %%%%%%%%%%%%%%%%%%%%%%%%%%%%%%%%%%%%%%%%%%%%%%%%%%%%%
% \newpage
% \noindent
% Definition of statement activation flags $\STMT^k_{t, pc}$ relies on the predecessors identified during preprocessing
% Definition of statement activation flags $\STMT^k_{t, pc}$ is based on the execution of preceding instructions, given in the \lstCPP{predecessors} map identified during preprocessing,
% Definition of statement activation flags $\STMT^k_{t, pc}$ relies on the following activation rules, based on the execution of preceding instructions, given in the \lstCPP{predecessors} map identified during preprocessing.
% \begin{itemize}
  % \item if the statement has no predecessor, or none of them is activated (\emph{reactivation}), assign it’s previous value if it’s hasn’t been executed
    % \[
      % \STMT^k_{t, pc} = \STMT^{k - 1}_{t, pc} \land \lnot \EXEC^{k - 1}_{t, pc}
    % \]
  % \item if any predecessor is activated, assign it’s execution variable’s value, otherwise follow the rule for statements without a predecessors
    % \[
      % \STMT^k_{t, pc} = \ITE(\STMT^{k - 1}_{t, pre}, \EXEC^{k - 1}_{t, pre}, \STMT^{k - 1}_{t, pc} \land \lnot \EXEC^{k - 1}_{t, pc})
    % \]
  % \item in the case of a jump, use the conjunction of the jump’s execution variable exec_k_t_pc and it’s condition to activate it’s target
    % \[
    % \]
% \end{itemize}
% Definition of statement activation flags $\STMT^k_{t, pc}$ is based on the execution of preceding instructions, given in the \lstCPP{predecessors} map identified during preprocessing.
% In order to either explicitly deactivate it if or otherwise preserve it's previous state (\emph{reactivation}), the corresponding update expression
% In order to either preserve the activation flag's previous state (\emph{reactivation}), or explicit deactivation in case of its execution, the corresponding update expression is initialized with a conjunction containing or deavtivate it in case of it's execution, the corresponding update expression is initialized with a conjunction
% If the corresponding statement is the initial program instruction, or none of it's predecessors was activated, assign
% Instead of another costly cardinality constraint, definition of statement activation variables $\STMT^k_{t, pc}$ is based on the execution of the corresponding instruction or any of it's preceding program statements to determine their successor by
% Instead of another costly cardinality constraint, definition of statement activation variables $\STMT^k_{t, pc}$ is based on the previous execution of the corresponding instruction and preceding program statements to either deactivate or activate them in the current step.
% To do so, each statement activation variables update expression is initialized with a conjunction containing it's previous state $\STMT^{k - 1}_{t, pc}$ together with the negation of the corresponding execution variable $\EXEC^{k - 1}_{t, pc}$ in order to deactivate it in case of it's execution or preserve it's state otherwise.
% >>>
% \todo[noline]{refer to interleaving assumption}
In order to correctly model symbolic program counters \CHANGE{correctly}, individual state updates must ensure that at most one statement is activated in any step.
% >>>
% In order to correctly model our symbolic program counters, individual state updates are based on the previous execution of the corresponding statement as well as preceding program instructions to ensure that at most one statement is activated in every step.
% In order to correctly model our symbolic program counters, individual updates of statement activation flags $\STMT^k_{t, pc}$ are based on their previous execution or preceding program instructions to ensure that at most one statement is activated in every step.
% In order to correctly model our symbolic program counters, individual updates of statement activation flags $\STMT^k_{t, pc}$ are based on the previous execution of the corresponding instruction as well as preceding program instructions to ensure that at most one statement is activated in every step.
% In order to correctly model our symbolic program counters, individual state updates of statement activation variables $\STMT^k_{t, pc}$ must ensure that at most one is activated in every step.
% \newpage
% \noindent
% $\STMT^{k - 1}_{t, pc}$t
% Instead of another costly cardinality constraint, the update expression of each statement activation variable $\STMT^k_{t, pc}$ is initialized with a conjunction containing it's previous activation and the negation of the corresponding execution variable $\EXEC^{k - 1}_{t, pc}$ for explicitly deactivating the statement in case of it's execution, while preserving it's state otherwise.
% Instead of another costly cardinality constraint, the update expression of each statement activation variable $\STMT^k_{t, pc}$ is initialized with a conjunction containing it's previous activation and the negation of the corresponding execution variable $\EXEC^{k - 1}_{t, pc}$ for explicitly deactivating the statement if it has been executed or preserving it's state otherwise.
% >>>
We do so by initializing the update expression of each statement activation variable $\STMT^k_{t, pc}$ with a conjunction containing its previous activation and the negation of the corresponding execution variable $\EXEC^{k - 1}_{t, pc}$ for explicitly deactivating the statement if it has been executed while preserving its state otherwise.
% >>>
% Each preceding instruction in the thread's program then adds an $\ITE$, setting the successor state to the predecessor's execution
% Similar to defining the successor of register states, each preceding instruction in the thread's program then embeds the expression's current value in the \emph{else} branch of an $\ITE$, activating the statement depending on the particular predecessor's execution.
% Similar to defining the successor of register states, each preceding program statement then embeds the expression's current value in the \emph{else} branch of an $\ITE$, either enabling or disabling the statement activation variable according to the predecessors execution in case it was activated.
% Similar to defining the successor of register states, each preceding program statement then embeds the expression's current value in the \emph{else} branch of an $\ITE$, either enabling or disabling the statement activation variable according to the predecessor's execution.
% Similar to defining the successor of register states, each preceding program statement then embeds the expression's current value in the \emph{else} branch of an $\ITE$, either enabling the statement activation variable depending to the predecessor's execution or keeping it disabled.
\CHANGE{Similarly} to defining the successor of register states, each preceding program statement then embeds the expression's current value in the \emph{else} branch of an $\ITE$, activating the statement depending on the predecessor's execution.
% In case of the predecessor being a conditional jump, identified by \lstCPP{bool is_jump (Instruction *)},
% In case of the statement being the target of a jump,
% Conditional jumps, detected by \lstCPP{is_jump}
% In case of the predecessor being a conditional jump, identified by \lstCPP{is_jump},
% Special care must be taken in case of the predecessor being a conditional jump, identified by \lstCPP{is_jump}, to either activate the target of a successfull jump using a conjunction of it's execution variable and condition, generated by the relevant \lstCPP{encode} implementation, or the negation thereof if it failed.
Special care must be taken in case of the predecessor being a conditional jump, identified by \lstCPP{is_jump}, to either activate the target using a conjunction of its execution variable and condition, generated by the relevant \lstCPP{encode} implementation, or the next statement by negating the condition in case of a failed jump.

% using a conjunction of the jump's execution variable together with it's condition, generated by the relevant \lstCPP{encode} implementation, or it's negation.
% >>>
% >>>

% preserve state
% deactivate in case of it's execution

% \todo[noline]{introduce anonymous is_jump function}

% \begin{lstlisting}[style=c++]
% void define_stmt (uint k)
% {
  % for (uint t = 0; t < programs.size(); t++)
    % for (uint pc = 0; pc < programs[t].size(); pc++)
      % {
        % // statement reactivation
        % std::string expr = land(stmt_var(k - 1, t, pc),
                              % lnot(exec_var(k - 1, t, pc)));
        % // activation by predecessor
        % const auto & stmts = predecessors[t][pc];
        % for (auto pre = stmts.rbegin(); pre != stmts.rend(); ++pre)
          % {
            % std::string exec = exec_var(k - 1, t, *pre);
            % Instruction * op = programs[t][*pre];
            % // add condition if predecessor is a jump
            % if (is_jump(op))
              % {
                % std::string cond = op->encode(k - 1, t, State::stmt);
                % if (!cond.empty())
                  % exec = land(exec,
                              % // negate condition for failed jumps
                              % *pre == pc - 1 && op.arg != pc
                                % ? lnot(cond)
                                % : cond);
              % }
            % expr = ite(stmt_var(k - 1, t, *pre), exec, expr);
          % }
        % formula << assign(stmt_var(k, t, pc), expr) << '\n';
      % }
% }
% \end{lstlisting}

\begin{lstlisting}[style=c++]
void define_stmt (uint k)
{
  for (uint t = 0; t < programs.size(); t++)
    for (uint pc = 0; pc < programs[t].size(); pc++)
      {
        // statement reactivation
        std::string next = land(stmt(k - 1, t, pc),
                              lnot(exec(k - 1, t, pc)));
        // activation by predecessor
        const auto & stmts = predecessors[t][pc];
        for (auto pre = stmts.rbegin(); pre != stmts.rend(); ++pre)
          {
            std::string val = exec(k - 1, t, *pre);
            Instruction * op = programs[t][*pre];
            // add condition if predecessor is a jump
            if (is_jump(op))
              {
                std::string cond = op->encode(k - 1, t, State::stmt);
                val = land(val, op->arg == pc ? cond : lnot(cond));
              }
            next = ite(stmt(k - 1, t, *pre), val, next);
          }
        formula << assign(stmt(k, t, pc), next) << '\n';
      }
}
\end{lstlisting}

\noindent
\emph{Example:} \lstCPP{define_stmt(1)}

\begin{lstlisting}[style=smtlib]
(assert (= stmt_1_0_0 (and stmt_0_0_0 (not exec_0_0_0))))
(assert (= stmt_1_0_1 (ite stmt_0_0_0
                           exec_0_0_0
                           (and stmt_0_0_1 (not exec_0_0_1)))))
(assert (= stmt_1_0_2 (ite stmt_0_0_1
                           exec_0_0_1
                           (and stmt_0_0_2 (not exec_0_0_2)))))
(assert (= stmt_1_0_3 (ite stmt_0_0_2
                           exec_0_0_2
                           (and stmt_0_0_3 (not exec_0_0_3)))))
(assert (= stmt_1_0_4 (ite stmt_0_0_3
                           exec_0_0_3
                           (and stmt_0_0_4 (not exec_0_0_4)))))
(assert (= stmt_1_1_0 (and stmt_0_1_0 (not exec_0_1_0))))
(assert (= stmt_1_1_1 (ite stmt_0_1_0
                           exec_0_1_0
                           (and stmt_0_1_1 (not exec_0_1_1)))))
(assert (= stmt_1_1_2 (ite stmt_0_1_1
                           exec_0_1_1
                           (and stmt_0_1_2 (not exec_0_1_2)))))
(assert (= stmt_1_1_3 (ite stmt_0_1_2
                           exec_0_1_2
                           (and stmt_0_1_3 (not exec_0_1_3)))))
(assert (= stmt_1_1_4 (ite stmt_0_1_3
                           exec_0_1_3
                           (and stmt_0_1_4 (not exec_0_1_4)))))
(assert (= stmt_1_2_0 (and stmt_0_2_0 (not exec_0_2_0))))
(assert (= stmt_1_2_1 (ite stmt_0_2_0
                           exec_0_2_0
                           (and stmt_0_2_1 (not exec_0_2_1)))))
(assert (= stmt_1_2_2 (ite stmt_0_2_1
                           exec_0_2_1
                           (and stmt_0_2_2 (not exec_0_2_2)))))
(assert (= stmt_1_2_3 (ite stmt_0_2_2
                           exec_0_2_2
                           (and stmt_0_2_3 (not exec_0_2_3)))))
(assert (= stmt_1_2_4 (ite stmt_0_2_3
                           (and exec_0_2_3 (not (= accu_0_2
                                                   #x0000)))
                           (and stmt_0_2_4 (not exec_0_2_4)))))
(assert (= stmt_1_2_5 (ite stmt_0_2_3
                           (and exec_0_2_3 (= accu_0_2 #x0000))
                           (and stmt_0_2_5 (not exec_0_2_5)))))
\end{lstlisting}

% \begin{lstlisting}[style=smtlib]
% (assert (= stmt_1_0_0 (and stmt_0_0_0 (not exec_0_0_0))))
% (assert (= stmt_1_0_1 (ite stmt_0_0_0
                           % exec_0_0_0
                           % (and stmt_0_0_1 (not exec_0_0_1)))))
% (assert (= stmt_1_0_2 (ite stmt_0_0_1
                           % exec_0_0_1
                           % (and stmt_0_0_2 (not exec_0_0_2)))))
% (assert (= stmt_1_0_3 (ite stmt_0_0_2
                           % exec_0_0_2
                           % (and stmt_0_0_3 (not exec_0_0_3)))))
% (assert (= stmt_1_0_4 (ite stmt_0_0_3
                           % exec_0_0_3
                           % (and stmt_0_0_4 (not exec_0_0_4)))))
% ...
% (assert (= stmt_1_2_0 (and stmt_0_2_0 (not exec_0_2_0))))
% (assert (= stmt_1_2_1 (ite stmt_0_2_0
                           % exec_0_2_0
                           % (and stmt_0_2_1 (not exec_0_2_1)))))
% (assert (= stmt_1_2_2 (ite stmt_0_2_1
                           % exec_0_2_1
                           % (and stmt_0_2_2 (not exec_0_2_2)))))
% (assert (= stmt_1_2_3 (ite stmt_0_2_2
                           % exec_0_2_2
                           % (and stmt_0_2_3 (not exec_0_2_3)))))
% (assert (= stmt_1_2_4 (ite stmt_0_2_3
                           % (and exec_0_2_3 (not (= accu_0_2
                                                   % #x0000)))
                           % (and stmt_0_2_4 (not exec_0_2_4)))))
% (assert (= stmt_1_2_5 (ite stmt_0_2_3
                           % (and exec_0_2_3 (= accu_0_2 #x0000))
                           % (and stmt_0_2_5 (not exec_0_2_5)))))
% \end{lstlisting}

% \todo[inline]{define block} %%%%%%%%%%%%%%%%%%%%%%%%%%%%%%%%%%%%%%%%%%%%%%%%%%%%

\noindent
% Block flags $\BLOCK^k_{id, t}$ are defined by a single $\ITE$, either unblocking the thread by resetting the state if all threads synchronized upon the corresponding checkpoint $id$ in the previous step, or assigning a conjunction including all execution variables $\EXEC^{k - 1}_{t, pc}$ of related \lstASM{CHECK id} instructions to enable it, as well as the predecessor $\BLOCK^{k - 1}_{id, t}$ for preserving it's state.
Block flags $\BLOCK^k_{id, t}$ are defined by a single $\ITE$, unblocking the thread by resetting the state if all threads synchronized upon the corresponding checkpoint $id$ in the previous step. %$k - 1$.
Otherwise, \CHANGE{we assign} a disjunction including all execution variables $\EXEC^{k - 1}_{t, pc}$ of related \lstASM{CHECK id} instructions enabling the state as well as the predecessor $\BLOCK^{k - 1}_{id, t}$ to presere its value.
%signalled by $\CHECK^{k - 1}_{id}$

\begin{lstlisting}[style=c++]
void define_block (uint k)
{
  for (const auto & [id, threads] : checkpoints)
    for (const auto & [t, stmts] : threads)
      {
        std::vector<std::string> args {block(k - 1, id, t)};
        for (uint pc : stmts)
          args.push_back(exec(k - 1, t, pc));
        formula << assign(block(k, id, t),
                          ite(check(k - 1, id),
                              "false",
                              lor(args))) << '\n';
      }
}
\end{lstlisting}

\noindent
\emph{Example:} \lstCPP{define_block(1)}

\begin{lstlisting}[style=smtlib]
(assert (= block_1_0_0 (ite check_0_0
                            false
                            (or block_0_0_0 exec_0_0_3))))
(assert (= block_1_0_1 (ite check_0_0
                            false
                            (or block_0_0_1 exec_0_1_3))))
(assert (= block_1_0_2 (ite check_0_0
                            false
                            (or block_0_0_2 exec_0_2_0))))
\end{lstlisting}

% \todo[inline]{define halt} %%%%%%%%%%%%%%%%%%%%%%%%%%%%%%%%%%%%%%%%%%%%%%%%%%%%%

% \begin{lstlisting}[style=c++]
% void define_halt (uint k)
% {
  % for (uint t = 0; t < programs.size(); t++)
    % {
      % if (halts.find(thread) != halts.end())
        % {
          % std::vector<std::string> args;
          % for (const auto pc : halts[t])
            % args.push_back(exec_var(k - 1, t, pc));
          % args.push_back(halt_var(k - 1, t));
          % formula << assign(halt_var(k, t), lor(args)) << eol;
        % }
      % else
        % formula << assign(halt_var(k, t), "false") << eol;
    % }
  % formula << eol;
% }
% \end{lstlisting}

\noindent
% Since there is no way of restarting a halted thread, halt flags $\Halt^k_t$ are enabled by a disjunction including the execution variables $\EXEC^{k - 1}_{t, pc}$ of \lstASM{HALT} instructions.
Since halt flags $\HALT^k_t$ cannot be reset once a thread has been stopped, they are defined by a simple disjunction including the execution variables $\EXEC^{k - 1}_{t, pc}$ of \lstASM{HALT} instructions and the previous state $\HALT^{k - 1}_t$ to preserve their value.

\newpage

\begin{lstlisting}[style=c++]
void define_halt (uint k)
{
  for (uint t = 0; t < programs.size(); t++)
    {
      std::vector<std::string> args {halt(k - 1, t)};
      for (uint pc : updates[State::halt][t])
        args.push_back(exec(k - 1, t, pc));
      formula << assign(halt(k, t), lor(args)) << '\n';
    }
}
\end{lstlisting}

\noindent
\emph{Example:} \lstCPP{define_halt(1)}

\begin{lstlisting}[style=smtlib]
(assert (= halt_1_0 (or halt_0_0 exec_0_0_4)))
(assert (= halt_1_1 (or halt_0_1 exec_0_1_4)))
(assert (= halt_1_2 halt_0_2))
\end{lstlisting}

% \todo[inline]{define heap} %%%%%%%%%%%%%%%%%%%%%%%%%%%%%%%%%%%%%%%%%%%%%%%%%%%%%

\noindent
% Definition of the global heap array $\HEAP^k$ again starts by initializing the update expression with it's predecessor $\HEAP^{k - 1}$, preserving the previous state.
Definition of our shared memory array $\HEAP^k$ again starts by initializing the update expression with its previous state $\HEAP^{k - 1}$.
% Each thread then embeds the expression in the \emph{else} branch an $\ITE$, storing the store buffer's entry in case it was flushed.
% Each thread then extends the expression by embedding it's current value in the \emph{else} branch an $\ITE$, storing the store buffer's entry in case it was flushed.
% Each thread then extends the expression by embedding it in an $\ITE$, storing the store buffer's entry in case it was flushed.
% Each thread then extends the expression by embedding it in an $\ITE$ for every atomic write, modifying the array according to the instruction's \lstCPP{encode} implementation in case of its execution, as well as storing the store buffer entry if it was flushed.
% Each thread then extends the expression by embedding it in an $\ITE$ for every atomic write, modifying the array according to the instruction's \lstCPP{encode} implementation in case of its execution and storing the store buffer entry if it was flushed.
Each thread then extends the expression by embedding it in an $\ITE$ for every atomic write, modifying the array according to the instruction's \lstCPP{encode} implementation, followed by a final \CHANGE{if-then-else} to store the store buffer entry in case it was flushed.
% Each thread then extends the expression by embedding it in an $\ITE$ for every atomic write, modifying the array according to the instruction's \lstCPP{encode} implementation in case of its execution and another one to store the store buffer entry if it was flushed.
%
% Furthermore, if the contains a statement atomically writing to memory, another $\ITE$, modifying array according to the instruction's \lstCPP{encode} implementation is added.
% Furthermore, if the contains a statement atomically writing to memory, another $\ITE$, modifying array according to the instruction's \lstCPP{encode} implementation is added.
% The actual update expression is then built by iteratively embeding it's previous value in the \emph{else} branch of an $\ITE$, storing the

\begin{lstlisting}[style=c++]
void define_heap (uint k)
{
  std::string next = heap(k - 1);
  for (int t = programs.size() - 1; t >= 0; t--)
    {
      const auto & stmts = updates[State::heap][t];
      for (auto pc = stmts.rbegin(); pc != stmts.rend(); ++pc)
        next = ite(exec(k - 1, t, *pc),
                   programs[t][*pc].encode(k - 1, t, State::heap),
                   next);
      next = ite(flush(k - 1, t),
                 store(heap(k - 1),
                       adr(k - 1, t),
                       val(k - 1, t)),
                 next);
    }
  formula << assign(heap(k, t), next) << '\n';
}
\end{lstlisting}

\newpage

\noindent
\emph{Example:} \lstCPP{define_heap(1)}

\begin{lstlisting}[style=smtlib]
(assert (= heap_1
           (ite flush_0_0
                (store heap_0 adr_0_0 val_0_0)
                (ite flush_0_1
                     (store heap_0 adr_0_1 val_0_1)
                     (ite flush_0_2
                          (store heap_0 adr_0_2 val_0_2)
                          heap_0)))))
\end{lstlisting}

% \todo[inline]{define exit-flag} %%%%%%%%%%%%%%%%%%%%%%%%%%%%%%%%%%%%%%%%%%%%%%%%

% \begin{lstlisting}[style=c++]
% void define_exit_flag (uint k)
% {
  % if (halts.empty() && exits.empty())
    % return;
%
  % std::vector<std::string> args {exit_flag_var(k - 1)};
%
  % if (!halts.empty())
    % {
      % std::vector<std::string> halt;
      % for (uint t = 0; t < programs.size(); t++)
        % halt.push_back(halt_var(k, t));
      % args.push_back(land(halt));
    % }
%
  % if (!exits.empty())
    % for (const auto & [t, stmts] : exits)
      % for (const auto pc : stmts)
        % args.push_back(exec_var(k - 1, t, pc));
%
  % formula << assign(exit_flag_var(k), lor(args)) << eol << eol;
% }
% \end{lstlisting}

\noindent
% In order to enable the exit flag $\EXIT^k$ according to our termination criteria, we define it as a disjunction containing the execution variables $\EXEC^{k - 1}_{t, pc}$ of every \lstASM{EXIT} statement, a conjunction over all halt variables $\HALT^k_t$ to stop the machine if no more thread is running and the previous state $\EXIT^{k - 1}$, preserving it's value.
In order to enable the exit flag $\EXIT^k$ according to our termination criteria, we define it as a disjunction containing the execution variables $\EXEC^{k - 1}_{t, pc}$ of every \lstASM{EXIT} statement, a conjunction over \CHANGE{the} halt variables $\HALT^k_t$ \CHANGE{of threads containing a call to \lstASM{HALT} for} stopping the machine if no more thread\CHANGE{s are} running and the previous state $\EXIT^{k - 1}$ to preserve its value.

\begin{lstlisting}[style=c++]
void define_exit(uint k)
{
  std::vector<std::string> halts;
  for (const auto & stmts : updates[State::halt])
    for (uint pc : stmts)
      halts.push_back(halt(k, t));
  std::vector<std::string> args {exit(k - 1), land(halts)};
  for (const auto & stmts : updates[State::exit])
    for (uint pc : stmts)
      args.push_back(exec(k - 1, t, pc));
  formula << assign(exit(k), lor(args)) << '\n';
}
\end{lstlisting}

\noindent
\emph{Example:} \lstCPP{define_exit(1)}

\begin{lstlisting}[style=smtlib]
(assert (= exit_1
           (or exit_0
               (and halt_1_0 halt_1_1)
               exec_0_2_4
               exec_0_2_5)))
\end{lstlisting}

% \todo[inline]{define exit-code} %%%%%%%%%%%%%%%%%%%%%%%%%%%%%%%%%%%%%%%%%%%%%%%%

% \begin{lstlisting}[style=c++]
% void define_exit_code ()
% {
  % std::string expr = consth(0);
  % if (!exits.empty())
    % for (uint k = 0; k <= bound; k++)
      % for (int t = programs.size() - 1; t >= 0; t--)
          % for (const auto & pc : exits[t])
            % expr = ite(exec_var(k, t, pc),
                      % programs[t][pc].encode(k, t, State::exit),
                      % expr);
  % formula << assign(exit_code_var, expr) << '\n';
% }
% \end{lstlisting}

\noindent
% Our functional encoding scheme is concluded by the definition of the machine's exit code $\EXITCODE$.
% Our functional encoding scheme is concluded by defining the machine's exit code $\EXITCODE$.
Our functional encoding scheme is concluded by setting the machine's exit code during the final step.
% >>>
% Since it is only assigned once, \emph{static single assignments} allow us to reduce the total number of variables by introducing only a single state
%, defined by an $\ITE$ cascade over all steps $k \in [0, bound]$
% for all steps $k \in [0, bound]$
% >>>
% Since it is unique in every execution, \emph{static single assignments} allow us to reduce the total number of variables by introducing only a single state for all $k \in [0, bound]$.
Since it is unique in every execution, explicit unrolling allows us to reduce the total number of variables by introducing only a single state $\EXITCODE$, %for all steps $k \in [0, bound]$.
% Since it is unique in every execution, the total number of variables can be reduced by introducing only a single state in our unrolled model checking problem,
% Since it is unique in every execution, unrolling our model checking problem for each step $k \in [0, bound]$ allows us to reduce the total number of variables by introducing only a single state for all $k \in [0, bound]$,
% assigned during the final step $k = bound$.
% The defining expression is initialized with zero and extended by embedding it in an $\ITE$
% embedding it in an $\ITE$, covering the execution of any \lstASM{EXIT} statement by assigning the according value, depending on
defined by an expression initialized with zero and embedded in an $\ITE$
for each step $k \in [0, bound]$ and \lstASM{EXIT} statement, assigning the corresponding exit code in case of its execution.
% >>>
% In order to reduce the total number of variables, only a single state is used in
% update expression initialized with the default exit code.

\begin{lstlisting}[style=c++]
void define_exit_code ()
{
  std::string next = consth(0);
  for (uint k = 0; k <= bound; k++)
    for (const auto & stmts : updates[State::exit])
      for (uint pc : stmts)
        next = ite(exec(k, t, pc),
                   programs[t][pc].encode(k, t, State::exit),
                   next);
  formula << assign(exit_code, next) << '\n';
}
\end{lstlisting}

\begin{lstlisting}[style=smtlib]
(assert (= exit-code
           (ite exec_17_2_5
                #x0001
                (ite exec_17_2_4
                     #x0000
                     ...
                     (ite exec_0_2_5
                          #x0001
                          (ite exec_0_2_4
                               #x0000
                               #x0000))))))
\end{lstlisting}

\newpage



\newpage

\section{Relational Next State Logic}

% > show difference in terms of runtime
% State of the art SAT and SMT solvers are able to exploit certain structures in the problem formulation, increasing performance in terms of runtime.
% Different encodings may therefore exhibit a rather drastical impact on the relative hardness of the resulting formula.
Since state of the art SAT and SMT solvers are able to exploit certain structures in the problem formulation, the relative hardness of different semantically equivalent encodings may vary quite drastically.
% State of the art SAT and SMT solvers  their performance partly by exploiting the structure of a given problem.
% structural changes in the problem formulation
% In order to show the differences of structural
% semantically equivalent encodings in terms of the
% generated by using the \texttt{-e smtlib-relational} command line parameter.
% better for combinatorial problems and small domains
% To highlight these differences, we include another encoding using a relational next state logic, sought to be more suitable for large combinatorial problems and small domains, generated by using the \texttt{-e smtlib-relational} command line parameter.
% relational next state logic
% requireing order of magnitude more time to solve than our previous functional approach, leading to a tree-like/hierarchical structure
% > highlight differences
% > next states fully implied by each transition
To highlight these differences,
we included another variant using a relational next state logic,
generated with the \texttt{-e smtlib-relational} command line parameter, % and
% based on states being fully implied by each transition.
% based on implying the complete set of successors for each transition.
based on implying the successors for each possible transition.
% > flat encoding sought of being more suitable for combinatorial problems and small domains
% > order of magnitude more time to solve compared to functional approach, resulting in a hierarchical structure | expression tree
% > simplified by term rewriting during the solver's preprocessing stage
Although this flat encoding scheme tends to work well with large combinatorial problems and small domains, \CHANGE{our experiments show that} it requires \CHANGE{considerably} more time to solve compared to our functional approach. %, which can be radically simplified by clever term rewriting due to its strictly hierarchical nature. %structure.

% \todo[inline]{define states (relational)} %%%%%%%%%%%%%%%%%%%%%%%%%%%%%%%%%%%%%%%

\begin{lstlisting}[style=c++]
void define_states (uint k)
{
  for (uint t = 0; t < programs.size(); t++)
    {
      imply_thread_executed(k, t);
      imply_thread_not_executed(k, t);
      imply_thread_flushed(k, t);
    }
  imply_machine_exited(k);
}
\end{lstlisting}

\noindent
We begin by defining variadic template functions for simplifying the generation of frequently reoccurring expressions, used to restore a state's previous value and additionally reset flags depending on a given condition.
% conditionally reset flags.
%for generating a frequently reoccuring pattern

% \begin{lstlisting}[style=c++, belowskip=\medskipamount]
\begin{lstlisting}[style=c++]
template <class R, class ... T>
R restore (R (*var) (uint k, T ... args), uint k, T ... args)
{
  return equal(var(k, args...), var(k - 1, args...));
}
\end{lstlisting}

\begin{lstlisting}[style=c++, aboveskip=0pt]
template <class R, class ... T>
R reset (R (*var) (uint k, uint x, T ... args),
         R (*cond) (uint k, uint x),
         uint k, uint x, T ... args)
{
  return equal(var(k, x, args...),
               ite(cond(k - 1, x),
                   "false",
                   var(k - 1, x, args...)));
}
\end{lstlisting}

\noindent
% Definition of successor states depending on the execution of a certain thread is simplified by introducing the following custom map,
% consolidating the next state expressions
% initialized with the common next state expressions
% Implying the successor states according to the execution of a certain thread requires m
% In order to minimize the effort when defining the successors for each possible instruction, we introduce the following custom map
% >>>
% Implying the whole set of states influenced by executing any instruction leads to most of them being unchanged and required to be preserved in the next step.
% Implying the whole set of states that may be influenced by executing any kind of instruction leads to most of them being unchanged and required to be preserved in the next step.
% Implying all states that might be influenced by executing a certain thread, leads to most of them being unchanged by any specific instruction.
Implying a common set of states that might be influenced by executing a certain thread
% leads to most of them being unchanged by any specific statement.
results in most of them \CHANGE{remaining} unchanged. % by any specific statement.
% Implying a common set of states for each statement that might be executed by a certain thread
% Since most states remain unchanged when executing a certain instruction, implying
% >>>
% To reduce the effort when implementing the next state functions of every instruction, we introduce the following custom map,
% To reduce the implementation effort for a particular instruction's encoding function, we introduce the following custom map,
% >>>
% To reduce the effort during the implementation of each instruction's encoding function, we introduce the following custom map,
% To reduce the implementation effort for each instruction's encoding function, we introduce the following custom map,
% To reduce the effort involved in implementing the next state function of each instruction, we introduce the following map, % !
% containing the default expressions for restoring the previous value of this common set of states,
% containing the common set of states, initialized with the default expression for restoring their value,
% initialized with the default expressions for restoring this common set of states in the current step,
% aggregating all relevant states and initializing them with the corresponding state-preserving expression, % !
% containing the common set of states, initialized with the default expression for restoring their value,
% >>>
% initialized the update expression of all common states
% thus requiring only the entries of states actually being altered to be replaced.
% in order to set only the states that actually change.
% such that only the ones set by the specific instruction needs to be overwritten
% thus requiring only the entries of states which are actually altered to be replaced. % !
% >>>
% We introduce the following next state map to reduce the effort involved in implementing the encoding function of each instruction by aggregating all relevant states and initializing them with the corresponding state-preserving expression, thus requiring only the entries of states which are actually altered to be replaced.
% >>>
% Furthermore, it offers an implicit \lstCPP{std::string} conversion operator, returning a conjunction to define the successors of all contained states at once.
% Furthermore, it offers an implicit \lstCPP{std::string} conversion operator, returning a conjunction to define the successors of all contained states at once.
% >>>
% In order to reduce the effort involved in implementing the encoding function of each instruction we introduce the following map, aggregating all relevant states and provide an implicit \lstCPP{std::string} conversion operator returning a conjunction to define the successors of all contained states at once.
In order to reduce the effort involved in implementing the encoding function of each instruction, we introduce the following map,
using the previously defined template functions to initialize it with the corresponding state-preserving expressions,
% and provide an implicit \lstCPP{std::string} conversion operator returning a conjunction to define the successors of all contained states at once.
thus requiring only the entries of states which are actually altered to be replaced
and provide an implicit \lstCPP{std::string} conversion operator returning a conjunction to define all successors at once.
% By using the previously defined template functions, we are able to add a constructor to initialize them with the corresponding state-preserving expressions for a given step and thread, thus requiring only the entries of states which are actually altered to be replaced.
% Furthermore, they are initialized with corresponding state-preserving expressions, thus requiring only the entries of states which are actually altered to be replaced.

% \todo[inline]{map containing next state expressions to simplify definition of successor states that depend on the execution of a certain thread}
% \todo[inline]{replace entries with the required next state expressions}
% \todo[inline]{implicit \lstCPP{std::string} conversion operator: returning a conjunction defining the successors for all states contained}
% \todo[inline]{used by \lstCPP{imply_thread_*executed}}

% // restore block
% if (!checkpoints.empty())
%   {
%     std::vector<std::string> block_vars;
%     for (const auto & [id, threads] : checkpoints)
%       if (threads.find(t) != threads.end())
%         block_vars.push_back(reset(&block_var,
%                                    &check_var,
%                                    k, id, t));
%     (*this)[State::block] = land(block_vars);
%   }
\begin{lstlisting}[style=c++]
struct Next : std::map<State, std::string>
{
  Next (k, t)
    {
      (*this)[State::accu] = restore(&accu, k, t);
      (*this)[State::mem] = restore(&mem, k, t);
      (*this)[State::adr] = restore(&adr, k, t);
      (*this)[State::val] = restore(&val, k, t);
      (*this)[State::full] = restore(&full, k, t);
      (*this)[State::halt] = restore(&halt, k, t);
      (*this)[State::heap] = restore(&heap, k);
      (*this)[State::exit] = lnot(exit(k, t));
      std::vector<std::string> block_vars;
      for (const auto & [id, threads] : checkpoints)
        if (threads.find(t) != threads.end())
          block_vars.push_back(reset(&block,
                                     &check,
                                     k, id, t));
      if (!block_vars.empty())
        (*this)[State::block] = land(block_vars);
    }

  operator std::string () const
    {
      std::vector<std::string> args;
      for (const auto & [_, expr] : *this)
        if (!expr.empty())
          args.push_back(expr);
      return land(args);
    }
};
\end{lstlisting}

% \todo[inline]{changed by every instruction and therefore not included in the next state map}
% \todo[inline]{requires setting all statement activation variables of a given thread}
% \todo[inline]{helper functions for activating a certain statement}

\newpage
\noindent
Since the statement activation variables $\STMT^k_{t, pc}$ of each thread must be fully defined for every possible transition, we include helper functions to generate the required conjunctions.

\begin{lstlisting}[style=c++]
std::string activate (uint k, uint t, uint pc)
{
  std::vector<std::string> stmts;
  for (uint i = 0; i < programs[t].size(); i++)
    stmts.push_back(i == pc ? stmt(k, t, pc)
                            : lnot(stmt(k, t, i)));
  return land(stmts);
}
\end{lstlisting}

% \todo[inline]{overload to be used by jump instructions, either activating the successor statement or jump target depending on its condition}
\noindent
% In terms of jump instructions, we define another overload for either activating the successor statement or the target, depending on the corresponding condition.
In case of jump instructions, another overload either activates the target or successor statement, depending on the corresponding condition.
% Another overload handles jump instructions by either activating the successor statement or the target, depending on the corresponding condition.

% \begin{lstlisting}[style=c++, aboveskip=0pt]
\begin{lstlisting}[style=c++]
std::string activate (uint k, uint t, uint pc, Instruction * jmp)
{
  std::vector<std::string> stmts;
  std::string condition = jmp->encode(k - 1, t, State::stmt);
  for (uint i = 0; i < programs[t].size(); i++)
    {
      std::string stmt_var = stmt(k, t, i);
      if (i == jmp->arg)
        stmts.push_back(ite(condition, stmt_var, lnot(stmt_var)));
      else if (i == pc + 1)
        stmts.push_back(ite(condition, lnot(stmt_var), stmt_var));
      else
        stmts.push_back(lnot(stmt_var));
    }
  return land(stmts);
}
\end{lstlisting}

% \todo[inline]{add another virtual member to \lstCPP{Instruction} class to dynamically dispatch encoding functions}

\noindent
% The next state functions, responsible for defining the successors when executing a certain instruction, are implemented by including yet another virtual member function to the \lstCPP{Instruction} class.
% Analogously to the encoding functions responsible for defining a specific next state
% by dynamically dispatching the according encoding function.
% Generating the next state for each possible instruction
% Encoding the common set of next states for every possible instruction
% Generating the common set of successor states for every possible instruction according to our relational encoding scheme
Generating a common set of successor states for every possible instruction %according to our relational encoding scheme
% according to our relational logic
% again relies on
is again based on
dynamically dispatched virtual member functions of concrete \lstCPP{Instruction} class instantiations.
% dispatched member functions achieved by adding the following virtual member function to the \lstCPP{Instruction} class
% We therefore introduce
Therefore, we introduce the following overload including the current program counter for activating the next statement. % and their actual implementations given below.

% \todo[inline]{adding the following virtual member function to the \lstCPP{Instruction} class and their concrete implementations given below}
% \todo[inline]{including its program counter to activate successor statements}

% \lstCPP{virtual std::string Instruction::encode (uint t, uint k, uint pc) = 0;}
\begin{lstlisting}[style=c++]
virtual std::string Instruction::encode (uint t, uint k, uint pc) = 0;
\end{lstlisting}

% \todo[inline]{encode LOAD (relational)} %%%%%%%%%%%%%%%%%%%%%%%%%%%%%%%%%%%%%%%%

\noindent
By using the previously defined helper constructs, the actual implementations can be \CHANGE{simplified} quite drastically and % are given below. % of every instruction's encoding function
% The actual implementations are given below and can be simplified quite drastically by using the previously defined helper constructs.
% The actual implementations can be simplified quite drastically by using the previously defined helper constructs and are given below.
\CHANGE{even though they therefore almost look the same, we include the full definitions for completeness.}

\newpage

\defencode{\lstASM{LOAD arg}}
\begin{lstlisting}[style=c++, style=encode]
std::string Load::encode (uint t, uint k, uint pc)
{
  Next next(k, t);
  next[State::accu] = equal(accu(k, t),
                            encode(k - 1, t, State::accu));
  next[State::stmt] = activate(k, t, pc + 1);
  return next;
}
\end{lstlisting}

% \todo[inline]{encode STORE (relational)} %%%%%%%%%%%%%%%%%%%%%%%%%%%%%%%%%%%%%%%

\defencode{\lstASM{STORE arg}}
\begin{lstlisting}[style=c++, style=encode]
std::string Store::encode (uint t, uint k, uint pc)
{
  Next next(k, t);
  next[State::adr] = equal(adr(k, t), encode(k - 1, t, State::adr));
  next[State::val] = equal(val(k, t), encode(k - 1, t, State::val));
  next[State::full] = full(k, t);
  next[State::stmt] = activate(k, t, pc + 1);
  return next;
}
\end{lstlisting}

% \todo[inline]{encode FENCE (relational)} %%%%%%%%%%%%%%%%%%%%%%%%%%%%%%%%%%%%%%%

\defencode{\lstASM{FENCE}}
\begin{lstlisting}[style=c++, style=encode]
std::string Fence::encode (uint t, uint k, uint pc)
{
  Next next(k, t);
  next[State::stmt] = activate(k, t, pc + 1);
  return next;
}
\end{lstlisting}

% \todo[inline]{encode ADD (relational)} %%%%%%%%%%%%%%%%%%%%%%%%%%%%%%%%%%%%%%%%%

\defencode{\lstASM{ADD arg}}
\begin{lstlisting}[style=c++, style=encode]
std::string Add::encode (uint t, uint k, uint pc)
{
  Next next(k, t);
  next[State::accu] = equal(accu(k, t),
                            encode(k - 1, t, State::accu));
  next[State::stmt] = activate(k, t, pc + 1);
  return next;
}
\end{lstlisting}

% \todo[inline]{encode ADDI (relational)} %%%%%%%%%%%%%%%%%%%%%%%%%%%%%%%%%%%%%%%%

\newpage

\defencode{\lstASM{ADDI arg}}
\begin{lstlisting}[style=c++, style=encode]
std::string Addi::encode (uint t, uint k, uint pc)
{
  Next next(k, t);
  next[State::accu] = equal(accu(k, t),
                            encode(k - 1, t, State::accu));
  next[State::stmt] = activate(k, t, pc + 1);
  return next;
}
\end{lstlisting}

% \todo[inline]{encode SUB (relational)} %%%%%%%%%%%%%%%%%%%%%%%%%%%%%%%%%%%%%%%%%

\defencode{\lstASM{SUB arg}}
\begin{lstlisting}[style=c++, style=encode]
std::string Sub::encode (uint t, uint k, uint pc)
{
  Next next(k, t);
  next[State::accu] = equal(accu(k, t),
                            encode(k - 1, t, State::accu));
  next[State::stmt] = activate(k, t, pc + 1);
  return next;
}
\end{lstlisting}

% \todo[inline]{encode SUBI (relational)} %%%%%%%%%%%%%%%%%%%%%%%%%%%%%%%%%%%%%%%%

\defencode{\lstASM{SUBI arg}}
\begin{lstlisting}[style=c++, style=encode]
std::string Subi::encode (uint t, uint k, uint pc)
{
  Next next(k, t);
  next[State::accu] = equal(accu(k, t),
                            encode(k - 1, t, State::accu));
  next[State::stmt] = activate(k, t, pc + 1);
  return next;
}
\end{lstlisting}

% \todo[inline]{encode MUL (relational)} %%%%%%%%%%%%%%%%%%%%%%%%%%%%%%%%%%%%%%%%%

\defencode{\lstASM{MUL arg}}
\begin{lstlisting}[style=c++, style=encode]
std::string Mul::encode (uint t, uint k, uint pc)
{
  Next next(k, t);
  next[State::accu] = equal(accu(k, t),
                            encode(k - 1, t, State::accu));
  next[State::stmt] = activate(k, t, pc + 1);
  return next;
}
\end{lstlisting}

% \todo[inline]{encode MULI (relational)} %%%%%%%%%%%%%%%%%%%%%%%%%%%%%%%%%%%%%%%%

\newpage

\defencode{\lstASM{MULI arg}}
\begin{lstlisting}[style=c++, style=encode]
std::string Muli::encode (uint t, uint k, uint pc)
{
  Next next(k, t);
  next[State::accu] = equal(accu(k, t),
                            encode(k - 1, t, State::accu));
  next[State::stmt] = activate(k, t, pc + 1);
  return next;
}
\end{lstlisting}

% \todo[inline]{encode CMP (relational)} %%%%%%%%%%%%%%%%%%%%%%%%%%%%%%%%%%%%%%%%%

\defencode{\lstASM{CMP arg}}
\begin{lstlisting}[style=c++, style=encode]
std::string Cmp::encode (uint t, uint k, uint pc)
{
  Next next(k, t);
  next[State::accu] = equal(accu(k, t),
                            encode(k - 1, t, State::accu));
  next[State::stmt] = activate(k, t, pc + 1);
  return next;
}
\end{lstlisting}

% \todo[inline]{encode JMP (relational)} %%%%%%%%%%%%%%%%%%%%%%%%%%%%%%%%%%%%%%%%%

\defencode{\lstASM{JMP arg}}
\begin{lstlisting}[style=c++, style=encode]
std::string Jmp::encode (uint t, uint k, uint pc)
{
  Next next(k, t);
  next[State::stmt] = activate(k, t, arg);
  return next;
}
\end{lstlisting}

% \todo[inline]{encode JZ (relational)} %%%%%%%%%%%%%%%%%%%%%%%%%%%%%%%%%%%%%%%%%%

\defencode{\lstASM{JZ arg}}
\begin{lstlisting}[style=c++, style=encode]
std::string Jz::encode (uint t, uint k, uint pc)
{
  Next next(k, t);
  next[State::stmt] = activate(k, t, pc, this);
  return next;
}
\end{lstlisting}

% \todo[inline]{encode JNZ (relational)} %%%%%%%%%%%%%%%%%%%%%%%%%%%%%%%%%%%%%%%%%

\newpage

\defencode{\lstASM{JNZ arg}}
\begin{lstlisting}[style=c++, style=encode]
std::string Jnz::encode (uint t, uint k, uint pc)
{
  Next next(k, t);
  next[State::stmt] = activate(k, t, pc, this);
  return next;
}
\end{lstlisting}

% \todo[inline]{encode JS (relational)} %%%%%%%%%%%%%%%%%%%%%%%%%%%%%%%%%%%%%%%%%%

\defencode{\lstASM{JS arg}}
\begin{lstlisting}[style=c++, style=encode]
std::string Js::encode (uint t, uint k, uint pc)
{
  Next next(k, t);
  next[State::stmt] = activate(k, t, pc, this);
  return next;
}
\end{lstlisting}

% \todo[inline]{encode JNS (relational)} %%%%%%%%%%%%%%%%%%%%%%%%%%%%%%%%%%%%%%%%%

\defencode{\lstASM{JNS arg}}
\begin{lstlisting}[style=c++, style=encode]
std::string Jns::encode (uint t, uint k, uint pc)
{
  Next next(k, t);
  next[State::stmt] = activate(k, t, pc, this);
  return next;
}
\end{lstlisting}

% \todo[inline]{encode JNZNS (relational)} %%%%%%%%%%%%%%%%%%%%%%%%%%%%%%%%%%%%%%%

\defencode{\lstASM{JNZNS arg}}
\begin{lstlisting}[style=c++, style=encode]
std::string Jnzns::encode (uint t, uint k, uint pc)
{
  Next next(k, t);
  next[State::stmt] = activate(k, t, pc, this);
  return next;
}
\end{lstlisting}

% \todo[inline]{encode MEM (relational)} %%%%%%%%%%%%%%%%%%%%%%%%%%%%%%%%%%%%%%%%%

\newpage

\defencode{\lstASM{MEM arg}}
\begin{lstlisting}[style=c++, style=encode]
std::string Mem::encode (uint t, uint k, uint pc)
{
  Next next(k, t);
  next[State::accu] = equal(accu(k, t),
                            encode(k - 1, t, State::accu));
  next[State::mem] = equal(mem(k, t), encode(k - 1, t, State::mem));
  next[State::stmt] = activate(k, t, pc + 1);
  return next;
}
\end{lstlisting}

% \todo[inline]{encode CAS (relational)} %%%%%%%%%%%%%%%%%%%%%%%%%%%%%%%%%%%%%%%%%

\defencode{\lstASM{CAS arg}}
\begin{lstlisting}[style=c++, style=encode]
std::string Cas::encode (uint t, uint k, uint pc)
{
  Next next(k, t);
  next[State::accu] = equal(accu(k, t),
                            encode(k - 1, t, State::accu));
  next[State::stmt] = activate(k, t, pc + 1);
  next[State::heap] = equal(heap(k),
                            encode(k - 1, t, State::heap));
  return next;
}
\end{lstlisting}

% \todo[inline]{encode CHECK (relational)} %%%%%%%%%%%%%%%%%%%%%%%%%%%%%%%%%%%%%%%

\defencode{\lstASM{CHECK arg}}
\begin{lstlisting}[style=c++, style=encode]
std::string Check::encode (uint t, uint k, uint pc)
{
  Next next(k, t);
  next[State::stmt] = activate(k, t, pc + 1);
  std::vector<std::string> blocks;
  for (const auto & [id, threads] : checkpoints)
    if (threads.find(t) != threads.end())
      blocks.push_back(id == op.arg
                         ? block(k, id, t)
                         : reset(&block, &check, k, id, t));
  next[State::block] = land(blocks);
  return state;
}
\end{lstlisting}

\noindent
\CHANGE{Since a thread might wait for different checkpoints},
% Since the might be multiple checkpoints per thread,
the encode function has to enable the corresponding blocking flag $\BLOCK^k_{id, t}$ while preserving or resetting any other.

% \todo[inline]{encode HALT (relational)} %%%%%%%%%%%%%%%%%%%%%%%%%%%%%%%%%%%%%%%%

\newpage

\defencode{\lstASM{HALT}}
\begin{lstlisting}[style=c++, style=encode]
std::string Halt::encode (uint t, uint k, uint pc)
{
  Next next(k, t);
  next[State::stmt] = activate(k, t, -1);
  if (programs.size() > 1)
    {
      std::vector<std::string> args;
      for (uint thread = 0; thread < programs.size(); thread++)
        if (thread != t)
          args.push_back(halt(k, thread));
      next[State::halt] =
        land(halt(k, t),
             ite(land(args),
                 land(exit(k), equal(exit_code, consth(0))),
                 lnot(exit(k))));
    }
  else
    next[State::halt] = land(halt(k, t),
                             exit(k),
                             equal(exit_code, consth(0))));
  state.erase(state.find(State::exit));
  return state;
}
\end{lstlisting}

\noindent
% Setting a thread's halt flag $\HALT^k_t$
% Stopping the machine if all halt flags $\HALT^k_t$ are set, we have to distinguish between single and multi threaded programs.
Stopping the machine if all threads halted is achieved by an $\ITE$, enabling the exit flag $\EXIT^k$ and assigning the default exit code if all halt flags $\HALT^k_t$ are set.
Otherwise, the exit flag is disabled to continue execution.
% In case only a single thread is running, the exit variables can be
Since the exit flag is now already contained in the \lstCPP{State::halt} entry, we have to erase \lstCPP{State::exit} to prevent conflicts due to its falsification. %, leading to a conflict.

% \todo[inline]{encode EXIT (relational)} %%%%%%%%%%%%%%%%%%%%%%%%%%%%%%%%%%%%%%%%

\defencode{\lstASM{EXIT arg}}
\begin{lstlisting}[style=c++, style=encode]
std::string Exit::encode (uint t, uint k, uint pc)
{
  Next next(k, t);
  next[State::stmt] = activate(k, t, -1);
  next[State::exit] = exit(k);
  next[State::exit_code] = equal(exit_code,
                                 encode(t, k, State::exit_code));
  return state;
}
\end{lstlisting}

\newpage

% \todo[inline]{imply thread executed} %%%%%%%%%%%%%%%%%%%%%%%%%%%%%%%%%%%%%%%%%%%

% \noindent
% Before we continue with the main functions, implying the next states for each transition,
% wrapper used to generate the assertions for implying the next states of a given transition.

% \begin{lstlisting}[style=c++]
% std::string imply_next (std::string var, std::string state)
% {
  % return assertion(imply(var, state));
% }
% \end{lstlisting}

\noindent
% With all utility constructs and instruction related encoding functions being defined, we are now able to generate the actual implications, starting with the successor states after executing any given statement.
With all utility constructs and instruction related encoding functions defined, we are now able to imply the next states for all transitions, starting with the execution of any given statement.
% Execution of a certain statement, signalled by $\EXEC^{k - 1}_{t, pc}$,
% Starting with the implication of successor states of
% using the previously defined enconding functions.

% // smtlib::Relational::imply_thread_executed -----------------------------------
% //
% // thread t executed an instruction (exec_k_t_pc):
% // * update thread state accordingly
% // * restore heap (or update upon a successful CAS)
% // * unset exit flag iff neither an EXIT, nor a HALT
% // * set exit code iff the instruction was an EXIT
\begin{lstlisting}[style=c++]
void imply_thread_executed (uint k, uint t)
{
  for (uint pc = 0; pc < programs[t].size(); pc++)
    formula << assertion(imply(exec(k - 1, t, pc),
                               programs[t][pc]->encode(k, t, pc)))
            << '\n';
}
\end{lstlisting}

\noindent
\emph{Example:} \lstCPP{imply_thread_executed(1, 0)}

\begin{lstlisting}[style=smtlib]
(assert (=> exec_0_0_0 ; ADDI 1
            (and (= accu_1_0 (bvadd accu_0_0 #x0001))
                 (= mem_1_0 mem_0_0)
                 (= adr_1_0 adr_0_0)
                 (= val_1_0 val_0_0)
                 (= full_1_0 full_0_0)
                 (and (not stmt_1_0_0)
                           stmt_1_0_1
                      (not stmt_1_0_2)
                      (not stmt_1_0_3)
                      (not stmt_1_0_4))
                 (= block_1_0_0 (ite check_0_0 false block_0_0_0))
                 (= halt_1_0 halt_0_0)
                 (= heap_1 heap_0)
                 (not exit_1))))
(assert (=> exec_0_0_1 ; STORE 0
            (and (= accu_1_0 accu_0_0)
                 (= mem_1_0 mem_0_0)
                 (= adr_1_0 #x0000)
                 (= val_1_0 accu_0_0)
                 full_1_0
                 (and (not stmt_1_0_0)
                      (not stmt_1_0_1)
                      stmt_1_0_2
                      (not stmt_1_0_3)
                      (not stmt_1_0_4))
                 (= block_1_0_0 (ite check_0_0 false block_0_0_0))
                 (= halt_1_0 halt_0_0)
                 (= heap_1 heap_0)
                 (not exit_1)))) ...
\end{lstlisting}

% \todo[inline]{imply thread not executed} %%%%%%%%%%%%%%%%%%%%%%%%%%%%%%%%%%%%%%%

\noindent
% In case a thread was not executed, all local states must be preserved, except the store
% In case a thread is not executed, it still might flush it's store buffer, % or have a block flag reset by another thread executing the corresponding checkpoint's final \lstASM{CHECK} instruction.
% or is released from a checkpoint if another thread just executed the corresponding \lstASM{CHECK} instruction.
% All local thread states are therefore preserved, while the store buffer full flag $\SBFULL^k_t$ and any block flag $\BLOCK^k_{id, t}$ might be reset.
% while preserving all local states.
% >>>
% In case a thread was not executed in the previous step, all local states must be preserved, while the store buffer full flag $\SBFULL^k_t$ and block flags $\BLOCK^k_{id, t}$ might be reset by a store buffer flush or due to being released from a checkpoint if another thread executed the corresponds \lstASM{CHECK} instruction.
In case a thread was not executed in the previous step, all local states must be preserved except the store buffer full flag $\SBFULL^k_t$ and block flags $\BLOCK^k_{id, t}$, as they might be reset by a store buffer flush, or due to being released from a checkpoint if another thread executed the corresponds \lstASM{CHECK} instruction.

% // imply_thread_not_executed
% //
% // thread t didn't execute an instruction (not thread_k_t):
% // * preserve thread state
% // * reset full iff the thread's store buffer has been flushed
\begin{lstlisting}[style=c++]
void imply_thread_not_executed (uint k, uint t)
{
  Next next(k, t);
  next[State::full] = reset(&full, &flush, k, t);
  std::vector<std::string> stmts;
  for (uint pc = 0; pc < programs.size(); pc++)
    stmts.push_back(restore(&stmt, k, t, pc));
  next[State::stmt] = land(stmts);
  next.erase(next.find(State::heap));
  next.erase(next.find(State::exit));

  formula << assertion(imply(lnot(thread(k - 1, t)), next)) << '\n';
}
\end{lstlisting}

\noindent
\emph{Example:} \lstCPP{imply_thread_not_executed(1, 0)}

\begin{lstlisting}[style=smtlib]
(assert (=> (not thread_0_0)
            (and (= accu_1_0 accu_0_0)
                 (= mem_1_0 mem_0_0)
                 (= adr_1_0 adr_0_0)
                 (= val_1_0 val_0_0)
                 (= full_1_0 (ite flush_0_0 false full_0_0))
                 (and (= stmt_1_0_0 stmt_0_0_0)
                      (= stmt_1_0_1 stmt_0_0_1)
                      (= stmt_1_0_2 stmt_0_0_2)
                      (= stmt_1_0_3 stmt_0_0_3)
                      (= stmt_1_0_4 stmt_0_0_4))
                 (= block_1_0_0 (ite check_0_0 false block_0_0_0))
                 (= halt_1_0 halt_0_0))))
\end{lstlisting}

% \todo[inline]{imply thread flushed} %%%%%%%%%%%%%%%%%%%%%%%%%%%%%%%%%%%%%%%%%%%%

\newpage

\noindent
% Since all thread states are already handeled,
With all local thread states already being set,
% Since resetting the store buffer full flag $\SBFULL^k_t$ is already handeled by the previous function,
% a thread flushing it's store buffer only needs to set the relevant global machine states, namely writing the heap array $\HEAP^k$ and disabling the exit flag $\EXIT^k$ to continue execution.
% a thread flushing it's store buffer
a store buffer flush
% only needs to write the heap array $\HEAP^k$ and disable the exit flag $\EXIT^k$ to continue executing.
% just writes the heap array $\HEAP^k$ and disables the exit flag $\EXIT^k$ to continue execution.
just needs to write the heap array $\HEAP^k$ and disable the exit flag $\EXIT^k$ to continue execution.
% continue execution by disabling the exit flag $\EXIT^k$.

% // imply_thread_flushed
% //
% // thread t flushed its store buffer (flush_k_t):
% // * update heap
% // * unset exit flag
\begin{lstlisting}[style=c++]
void imply_thread_flushed (uint k, uint t)
{
  formula << assertion(imply(flush(k - 1, t),
                             land(equal(heap(k),
                                        store(heap(k - 1),
                                              adr(k - 1, t),
                                              val(k - 1, t)))
                                  lnot(exit(k))))) << '\n';
}
\end{lstlisting}

\noindent
\emph{Example:} \lstCPP{imply_thread_flushed(1, 0)}

\begin{lstlisting}[style=smtlib]
(assert (=> flush_0_0
            (and (= heap_1 (store heap_0 adr_0_0 val_0_0))
                 (not exit_1))))
\end{lstlisting}

% \todo[inline]{imply machine exited} %%%%%%%%%%%%%%%%%%%%%%%%%%%%%%%%%%%%%%%%%%%%

\noindent
Finally, to keep the model consistent if the machine terminated in a step $k < bound$, the heap array $\HEAP^{k - 1}$ must be preserved.
% On the other hand, if the machine still has not exited in the final step, the exit code needs to be set since it would remain undefined.
On the other hand, if the machine is still running in the final step,
% the exit code needs to be set since it would remain undefined.
we need to set the default exit code as it would otherwise remain undefined.

% // imply_machine_exited
% //
% // machine exited (exit_k):
% // * preserve exit flag
% // * preserve heap
% // * set exit code to zero iff machine never exited (step == bound)
\begin{lstlisting}[style=c++]
void imply_machine_exited (uint k)
{
  formula << assertion(imply(exit(k - 1),
                             land(restore(&heap, k),
                                  exit(k)))) << '\n';
  if (k == bound)
    formula << assertion(imply(lnot(exit(k)),
                               equal(exit_code, consth(0))))
            << '\n';
}
\end{lstlisting}

\noindent
\emph{Example:} \lstCPP{imply_machine_exited(1, 0)}

\begin{lstlisting}[style=smtlib]
(assert (=> exit_0 (and (= heap_1 heap_0) exit_1)))
\end{lstlisting}

\noindent
\emph{Example:} \lstCPP{imply_machine_exited(17, 0)}

\begin{lstlisting}[style=smtlib]
(assert (=> exit_16 (and (= heap_17 heap_16) exit_17)))
(assert (=> (not exit_17) (= exit-code #x0000)))
\end{lstlisting}



\newpage

\section{BTOR2}

\newcommand{\LPREV}{\rotatebox[origin=c]{180}{$\Lsh$}}
\newcommand{\RPREV}{\rotatebox[origin=c]{180}{$\Rsh$}}

% We also included the possibility to generate the novel word level model checking format \BTOR \cite{ref:BTOR2}, providing a sequential extensions for specifying word-level model checking problems with registers and memories.
We also included the possibility to generate the word level model checking format {\BTOR} \cite{ref:BTOR2}, using the \texttt{-e btor2} command line parameter.
% providing a \emph{sequential extension} for specifying % register and array states in combination with a transition function,
% By providing a \emph{sequential extension} for specifying % register and array states in combination with a transition function,
It provides a \emph{sequential extension} for specifying % register and array states in combination with a transition function,
states in combination with their transition functions,
which are automatically unrolled via symbolic substitution by the accompanying bounded model checker BtorMC.
% the generated formulas size remains constant with respect to the upper bound, compared to the linear growth of our {\SMTLIB} encodings.
% it produces constant sized formulas with respect to the upper bound
% In contrast to \SMTLIB, which requires manual unrolling, \BTOR produces constant sized formulas with respect to the upper bound.
% Automatic unrolling via symbolic substitution by the accompanying bounded model checker BtorMC
% >>>
In contrast to the linear growth of our {\SMTLIB} encodings, the generated formulas size therefore remains constant for any given bound. % with respect to the upper bound.
% The generated formulas size therefore remains constant with respect to the upper bound, compared to the linear growth of our {\SMTLIB} encodings.
% In contrast to {\SMTLIB}, which requires manual unrolling, it produces constant sized formulas with respect to the upper bound by providing a \emph{sequential extension} for specifying states in combination with their transition functions.
% states in combination with their transition functions.
% symbolicly substituting current state expressions into next state functions and incremental SMT solving.
% Its line based syntax on the other hand requires more housekeeping, leading to a slightly more complex encoding process.
% We omit an outline of the encoding process, as it is semantically equivalent to the {\SMTLIB} formulas using our functional next state logic and a bit more tedious due to the line based syntax.
% We omit an outline of the encoding process, as it is semantically equivalent to the previous one using our functional next state logic and a bit more tedious due to the line based syntax.
We omit further details about the encoding process, as it is more or less identical to the previous using our functional next state logic, but a bit more tedious due to the line based syntax.
% Instead, the encoding of our store buffer litmus test example is reviewed st
Instead, we review the encoding of our store buffer litmus test example to highlight the pros and cons of {\BTOR}.

% \bigskip
% \noindent
{\BTOR} formulas consist of nodes, one per line, each prefixed with a strictly increasing numeric identifier serving as the input to subsequent nodes.
% > harder to read | decrease readability
% > easier to parse | implicitly cycle free
% > direct access to every part of the expression tree
By requiring every node to be defined before being used, the resulting direct acyclic graph is easy to parse
% and allows transparent access to every part of the expression, thus avoiding the need for additional auxiliary variables.
% and no additional auxiliary variables are needed since every part of the expression may be referenced at a later point.
and since every part of the expression may be referenced at a later point,
% the need for additional auxiliary variables is avoided.
no additional auxiliary variables are needed.
% While this decreases readability, the resulting direct acyclic graph is easy to parse and allow direct reference to any part of without the need of any additional variable.
% These main restrictions slightly decrease readability but easier to pass since everything needs to be defined before
% By using these strictly increasing identifiers as inputs to subsequent nodes, no other
% By requiring these strictly increasing identifiers to be the only form of input for subsequent nodes,
% no other variables are needed to build the expression tree
% These strictly increasing identifiers are used as a reference during the definition of subsequent nodes for
% used as a reference in the definition of subsequent nodes
% Primary inputs to the formula are declared by the \lstBTOR{input} keyword for free variables and \lstBTOR{state} for flags, registers and arrays.
See \cite{ref:BTOR2} for a detailed format description and the list of available nodes.

The encoding of our store buffer litmus test example given in Table \ref{tbl:encoding:programs} starts by defining the required boolean, bit-vector and array sorts.
\begin{lstlisting}[style=btor2]
1 sort bitvec 1   `\lstCMT{$\BVSORT[1]$}`
2 sort bitvec 16  `\lstCMT{$\BVSORT$}`
3 sort array 2 2  `\lstCMT{$\ASORT$}`
\end{lstlisting}
Next, a special set of input nodes is used to define the required boolean
% Because operands are limited to identifiers, the required boolean and
% Followed by input nodes for the required boolean
\begin{lstlisting}[style=btor2]
4 zero 1
5 one 1
\end{lstlisting}
% bit-vector constants have to be defined through a special set of input nodes.
and bit-vector constants.
\begin{lstlisting}[style=btor2]
6 zero 2
7 one 2
8 constd 2 5
9 constd 2 17
\end{lstlisting}

\noindent
% Register and array states are declared by the keyword \lstBTOR{state}, may be initialized using \lstBTOR{init} and their transition function defined via \lstBTOR{next}.
% In order to initialize our $\HEAP$ array according to example's memory map given in Listing \ref{lst:mmap:init.mmap},
% we introduce a \lstBTOR{state} node.
% Initialize temporary memory map array state.
% the keyword \lstBTOR{state} to declare registers and memories in combination with their transition function through \lstBTOR{next}.
% The main advantage of {\BTOR} is the ability to declare states via the keyword \lstBTOR{state} and their transition function using \lstBTOR{next}.
% In order to initialize our $\HEAP$ array, a constant array \lstBTOR{state} is introduced and set according to the example's memory map given in Listing \ref{lst:mmap:init.mmap}.
In order to initialize our $\HEAP$ array with the contents of this example's memory map given in Listing \ref{lst:mmap:init.mmap}, an uninitialized array state is declared and the corresponding elements written accordingly.
% By omitting it's transition function, it is treated as a constant.
By omitting its transition function, it is treated as a primary input only used during initialization and can safely be ignored in later steps.
\begin{lstlisting}[style=btor2]
10 state 3        `\lstCMT{mmap}`
11 write 3 10 6 6 `\lstCMT{$\texttt{mmap}[0] \leftarrow 0$}`
12 write 3 11 7 6 `\lstCMT{$\texttt{mmap}[1] \leftarrow 0$}`
\end{lstlisting}

\noindent
% Next, all states are declared:
Since our actual machine states are used in the definition of constraints and transition functions, they must be declared in advance.
% The encoding continues with the declaration of machine states
\bigbreak
\noindent
$\bullet \;$ Accumulator registers: % $\ACCU_t$:
\begin{lstlisting}[style=btor2]
13 state 2 `\lstCMT{$\ACCU_0$}`
14 state 2 `\lstCMT{$\ACCU_1$}`
15 state 2 `\lstCMT{$\ACCU_2$}`
\end{lstlisting}

\noindent
$\bullet \;$ CAS memory registers: % $\MEM_t$:
\begin{lstlisting}[style=btor2]
16 state 2 `\lstCMT{$\MEM_0$}`
17 state 2 `\lstCMT{$\MEM_1$}`
18 state 2 `\lstCMT{$\MEM_2$}`
\end{lstlisting}

\noindent
$\bullet \;$ Store buffer address registers: % $\SBADR_t$:
\begin{lstlisting}[style=btor2]
19 state 2 `\lstCMT{$\SBADR_0$}`
20 state 2 `\lstCMT{$\SBADR_1$}`
21 state 2 `\lstCMT{$\SBADR_2$}`
\end{lstlisting}

\noindent
$\bullet \;$ Store buffer value registers: % $\SBVAL_t$:
\begin{lstlisting}[style=btor2]
22 state 2 `\lstCMT{$\SBVAL_0$}`
23 state 2 `\lstCMT{$\SBVAL_1$}`
24 state 2 `\lstCMT{$\SBVAL_2$}`
\end{lstlisting}

\noindent
$\bullet \;$ Store buffer full flags: % $\SBFULL_t$:
\begin{lstlisting}[style=btor2]
25 state 1 `\lstCMT{$\SBFULL_0$}`
26 state 1 `\lstCMT{$\SBFULL_1$}`
27 state 1 `\lstCMT{$\SBFULL_2$}`
\end{lstlisting}

\noindent
$\bullet \;$ Statement activation flags: % $\STMT_{t, pc}$:
\begin{lstlisting}[style=btor2]
28 state 1 `\lstCMT{$\STMT_{0, 0}$}`
29 state 1 `\lstCMT{$\STMT_{0, 1}$}`
30 state 1 `\lstCMT{$\STMT_{0, 2}$}`
31 state 1 `\lstCMT{$\STMT_{0, 3}$}`
32 state 1 `\lstCMT{$\STMT_{0, 4}$}`

33 state 1 `\lstCMT{$\STMT_{1, 0}$}`
34 state 1 `\lstCMT{$\STMT_{1, 1}$}`
35 state 1 `\lstCMT{$\STMT_{1, 2}$}`
36 state 1 `\lstCMT{$\STMT_{1, 3}$}`
37 state 1 `\lstCMT{$\STMT_{1, 4}$}`

38 state 1 `\lstCMT{$\STMT_{2, 0}$}`
39 state 1 `\lstCMT{$\STMT_{2, 1}$}`
40 state 1 `\lstCMT{$\STMT_{2, 2}$}`
41 state 1 `\lstCMT{$\STMT_{2, 3}$}`
42 state 1 `\lstCMT{$\STMT_{2, 4}$}`
43 state 1 `\lstCMT{$\STMT_{2, 5}$}`
\end{lstlisting}

\noindent
$\bullet \;$ Block flags: % $\BLOCK_{id, t}$:
\begin{lstlisting}[style=btor2]
44 state 1 `\lstCMT{$\BLOCK_{0, 0}$}`
45 state 1 `\lstCMT{$\BLOCK_{0, 1}$}`
46 state 1 `\lstCMT{$\BLOCK_{0, 2}$}`
\end{lstlisting}

\noindent
$\bullet \;$ Halt flags: % $\HALT_t$:
\begin{lstlisting}[style=btor2]
47 state 1 `\lstCMT{$\HALT_0$}`
48 state 1 `\lstCMT{$\HALT_1$}`
49 state 1 `\lstCMT{$\HALT_2$}`
\end{lstlisting}

\noindent
$\bullet \;$ Shared memory: % $\HEAP$:
\begin{lstlisting}[style=btor2]
50 state 3 `\lstCMT{$\HEAP$}`
\end{lstlisting}

\noindent
$\bullet \;$ Exit flag: % $\EXIT$:
\begin{lstlisting}[style=btor2]
51 state 1 `\lstCMT{$\EXIT$}`
\end{lstlisting}

\noindent
$\bullet \;$ Exit code: % $\EXITCODE$:
\begin{lstlisting}[style=btor2]
52 state 2 `\lstCMT{$\EXITCODE$}`
\end{lstlisting}

\noindent
Same applies to the free transition variables $\THREAD_t$ and $\FLUSH_t$ serving as the model's input.
% serving as the model's inputs declared as free variables through the keyword \lstBTOR{input}.
% Because they serve as the model's input, we declare them as free variables through the keyword \lstBTOR{input}.
\begin{lstlisting}[style=btor2]
53 input 1 `\lstCMT{$\THREAD_0$}`
54 input 1 `\lstCMT{$\THREAD_1$}`
55 input 1 `\lstCMT{$\THREAD_2$}`

56 input 1 `\lstCMT{$\FLUSH_0$}`
57 input 1 `\lstCMT{$\FLUSH_1$}`
58 input 1 `\lstCMT{$\FLUSH_2$}`
\end{lstlisting}

\noindent
% In contrast to our {\SMTLIB} encoding, no additional declarations of helper variables are needed and our execution variables $\EXEC_{t, pc}$ can be defined by simple \lstBTOR{and} nodes.
In contrast to the {\SMTLIB} encoding, no additional declarations are needed for the introduction of helper variables and our execution variables $\EXEC_{t, pc}$ therefore defined as simple \lstBTOR{and} nodes.
% Execution variable definitions.
\begin{lstlisting}[style=btor2]
59 and 1 28 53        `\lstCMT{$\STMT_{0, 0} \land \THREAD_0 \equiv \EXEC_{0, 0}$}`
60 and 1 29 53        `\lstCMT{$\STMT_{0, 1} \land \THREAD_0 \equiv \EXEC_{0, 1}$}`
61 and 1 30 53        `\lstCMT{$\STMT_{0, 2} \land \THREAD_0 \equiv \EXEC_{0, 2}$}`
62 and 1 31 53        `\lstCMT{$\STMT_{0, 3} \land \THREAD_0 \equiv \EXEC_{0, 3}$}`
63 and 1 32 53        `\lstCMT{$\STMT_{0, 4} \land \THREAD_0 \equiv \EXEC_{0, 4}$}`

64 and 1 33 54        `\lstCMT{$\STMT_{1, 0} \land \THREAD_1 \equiv \EXEC_{1, 0}$}`
65 and 1 34 54        `\lstCMT{$\STMT_{1, 1} \land \THREAD_1 \equiv \EXEC_{1, 1}$}`
66 and 1 35 54        `\lstCMT{$\STMT_{1, 2} \land \THREAD_1 \equiv \EXEC_{1, 2}$}`
67 and 1 36 54        `\lstCMT{$\STMT_{1, 3} \land \THREAD_1 \equiv \EXEC_{1, 3}$}`
68 and 1 37 54        `\lstCMT{$\STMT_{1, 4} \land \THREAD_1 \equiv \EXEC_{1, 4}$}`

69 and 1 38 55        `\lstCMT{$\STMT_{2, 0} \land \THREAD_2 \equiv \EXEC_{2, 0}$}`
70 and 1 39 55        `\lstCMT{$\STMT_{2, 1} \land \THREAD_2 \equiv \EXEC_{2, 1}$}`
71 and 1 40 55        `\lstCMT{$\STMT_{2, 2} \land \THREAD_2 \equiv \EXEC_{2, 2}$}`
72 and 1 41 55        `\lstCMT{$\STMT_{2, 3} \land \THREAD_2 \equiv \EXEC_{2, 3}$}`
73 and 1 42 55        `\lstCMT{$\STMT_{2, 4} \land \THREAD_2 \equiv \EXEC_{2, 4}$}`
74 and 1 43 55        `\lstCMT{$\STMT_{2, 5} \land \THREAD_2 \equiv \EXEC_{2, 5}$}`
\end{lstlisting}

\noindent
% Checkpoint variable definitions.
% Similarly, the checkpoint variable $\CHECK_0$ is defined through a conjunction over the corresponding block flags $\BLOCK_{0, t}$.
% Due to the fixed arity of {\BTOR} nodes,
The fixed arity of {\BTOR} operators imposes a minor inconvenience
% The absence of variadic versions for associative operators in {\BTOR} imposes a minor inconvenience
% when building expressions over a larger number of variables.
% if a connective has to be applied over a large number of variables.
if the same connective has to be applied over a larger number of variables.
% applying a connective over a larger number of variables.
% as it requires a cascade of \lstBTOR{and} nodes to define the checkpoint variable $\CHECK_0$ as a conjunction over the corresponding block flags $\BLOCK_{0, t}$.
For example, a cascade of \lstBTOR{and} nodes is required to build the conjunction over all block flags $\BLOCK_{0, t}$ representing the checkpoint variable $\CHECK_0$.
% For example, the conjunction over all block flags $\BLOCK_{0, t}$ representing the checkpoint variable $\CHECK_0$ requires a cascade of \lstBTOR{and} nodes.
% $ \CHECK_0 \equiv \bigwedge_{0 \leq t \leq 3} \BLOCK_{0, t}$.
\begin{lstlisting}[style=btor2]
75 and 1 44 45        `\lstCMT{$\BLOCK_{0, 0} \land \BLOCK_{0, 1}$}`
76 and 1 46 75        `\lstCMT{$\LPREV \land \BLOCK_{0, 2} \equiv \CHECK_0$}`
\end{lstlisting}

\noindent
% Cardinality constraints.
% Similarly, the clause representing the required greater than one constraint
% Similarly, the disjunction serving as the \emph{greater than zero} part of our \emph{exactly one} cardinality constraint is defined by a cascade of \lstBTOR{or} nodes.
Similarly, the disjunction serving as the \emph{at-least-one} predicate of our \emph{exactly-one} constraint is defined by a cascade of \lstBTOR{or} nodes
% $>^1_7\!(\THREAD_0, \THREAD_1, \THREAD_2, \FLUSH_0, \FLUSH_1, \FLUSH_2, \EXIT)$
% \[
  % \EXIT \lor (\bigvee_{0 \leq t \leq 3} \THREAD_t \lor \FLUSH_t)
% \]
% needs to be defined through a cascade of disjunctions.
\begin{lstlisting}[style=btor2]
77 or 1 53 54         `\lstCMT{$\THREAD_0 \lor \THREAD_1$}`
78 or 1 55 77         `\lstCMT{$\LPREV \lor \THREAD_2$}`
79 or 1 56 78         `\lstCMT{$\LPREV \lor \FLUSH_0$}`
80 or 1 57 79         `\lstCMT{$\LPREV \lor \FLUSH_1$}`
81 or 1 58 80         `\lstCMT{$\LPREV \lor \FLUSH_2$}`
82 or 1 51 81         `\lstCMT{$\LPREV \lor \EXIT$}`
\end{lstlisting}
% The corresponding invariant is then added through the keyword \lstBTOR{constraint}.
% and the corresponding invariant then added through the keyword \lstBTOR{constraint}.
and added as an invariant through the keyword \lstBTOR{constraint}.
% `\lstCMT{$>^1_7\!(\THREAD_0, \ldots, \FLUSH_0, \ldots, \EXIT)$}`
\begin{lstlisting}[style=btor2]
83 constraint 82
\end{lstlisting}
% Sequential counter constraint
% Definition of the \emph{at-most-one} predicate, implemented by Carsten Sinz's sequential counter constraint $\text{LT}^{7, 1}_{\text{SEQ}}$ \cite{ref:Sinz} starts with the declaration of the required auxiliary input variables.
Definition of the remaining \emph{at-most-one} predicate starts with the declaration of auxiliary input variables required by Carsten Sinz's sequential counter constraint $\text{LT}^{7, 1}_{\text{SEQ}}$ \cite{ref:Sinz},
% The \emph{exactly-one} constraint is then completed by adding the \emph{at-most-one} predicate implemented through Carsten Sinz's sequential counter constraint $\text{LT}^{7, 1}_{\text{SEQ}}$.
\begin{lstlisting}[style=btor2]
84 input 1            `\lstCMT{$\THREAD^{\text{aux}}_0$}`
85 input 1            `\lstCMT{$\THREAD^{\text{aux}}_1$}`
86 input 1            `\lstCMT{$\THREAD^{\text{aux}}_2$}`
87 input 1            `\lstCMT{$\FLUSH^{\text{aux}}_0$}`
88 input 1            `\lstCMT{$\FLUSH^{\text{aux}}_1$}`
89 input 1            `\lstCMT{$\FLUSH^{\text{aux}}_2$}`
\end{lstlisting}
% followed by the definition of clauses.
followed by the corresponding set of clauses.
\begin{lstlisting}[style=btor2]
90 or 1 -53 84        `\lstCMT{$\lnot \THREAD_0 \lor \THREAD_0^{\text{aux}}$}`
91 or 1 -54 85        `\lstCMT{$\lnot \THREAD_1 \lor \THREAD_1^{\text{aux}}$}`
92 or 1 -84 85        `\lstCMT{$\lnot \THREAD_0^{\text{aux}} \lor \THREAD_1^{\text{aux}}$}`
93 or 1 -54 -84       `\lstCMT{$\lnot \THREAD_1 \lor \lnot \THREAD_0^{\text{aux}}$}`
94 or 1 -55 86        `\lstCMT{$\lnot \THREAD_2 \lor \THREAD_2^{\text{aux}}$}`
95 or 1 -85 86        `\lstCMT{$\lnot \THREAD_1^{\text{aux}} \lor \THREAD_2^{\text{aux}}$}`
96 or 1 -55 -85       `\lstCMT{$\lnot \THREAD_2 \lor \lnot \THREAD_1^{\text{aux}}$}`
97 or 1 -56 87        `\lstCMT{$\lnot \FLUSH_0 \lor \FLUSH_0^{\text{aux}}$}`
98 or 1 -86 87        `\lstCMT{$\lnot \THREAD_2^{\text{aux}} \lor \FLUSH_0^{\text{aux}}$}`
99 or 1 -56 -86       `\lstCMT{$\lnot \FLUSH_0 \lor \THREAD_2^{\text{aux}}$}`
100 or 1 -57 88       `\lstCMT{$\lnot \FLUSH_1 \lor \FLUSH_0^{\text{aux}}$}`
101 or 1 -87 88       `\lstCMT{$\lnot \FLUSH_0^{\text{aux}} \lor \FLUSH_1^{\text{aux}}$}`
102 or 1 -57 -87      `\lstCMT{$\lnot \FLUSH_1 \lor \lnot \FLUSH_0^{\text{aux}}$}`
103 or 1 -58 89       `\lstCMT{$\lnot \FLUSH_2 \lor \FLUSH_1^{\text{aux}}$}`
104 or 1 -88 89       `\lstCMT{$\lnot \FLUSH_1^{\text{aux}} \lor \FLUSH_2^{\text{aux}}$}`
105 or 1 -58 -88      `\lstCMT{$\lnot \FLUSH_2 \lor \lnot \FLUSH_1^{\text{aux}}$}`
106 or 1 -51 -89      `\lstCMT{$\lnot \EXIT \lor \lnot \FLUSH_2^{\text{aux}}$}`
\end{lstlisting}
% The final $\leq^1_7$ \lstBTOR{constraint} can then be built by a conjunction over the previously defined clauses.
The \emph{exactly-one} constraint can now be completed by building a conjunction over aforementioned clauses
\begin{lstlisting}[style=btor2]
107 and 1 90 91       `\lstCMT{$(\lnot \THREAD_0 \lor \THREAD_0^{\text{aux}}) \land (\lnot \THREAD_1 \lor \THREAD_1^{\text{aux}})$}`
108 and 1 92 107      `\lstCMT{$\LPREV \land (\lnot \THREAD_0^{\text{aux}} \lor \THREAD_1^{\text{aux}})$}`
109 and 1 93 108      `\lstCMT{$\LPREV \land (\lnot \THREAD_1 \lor \lnot \THREAD_0^{\text{aux}})$}`
110 and 1 94 109      `\lstCMT{$\LPREV \land (\lnot \THREAD_2 \lor \THREAD_2^{\text{aux}})$}`
111 and 1 95 110      `\lstCMT{$\LPREV \land (\lnot \THREAD_1^{\text{aux}} \lor \THREAD_2^{\text{aux}})$}`
112 and 1 96 111      `\lstCMT{$\LPREV \land (\lnot \THREAD_2 \lor \lnot \THREAD_1^{\text{aux}})$}`
113 and 1 97 112      `\lstCMT{$\LPREV \land (\lnot \FLUSH_0 \lor \FLUSH_0^{\text{aux}})$}`
114 and 1 98 113      `\lstCMT{$\LPREV \land (\lnot \THREAD_2^{\text{aux}} \lor \FLUSH_0^{\text{aux}})$}`
115 and 1 99 114      `\lstCMT{$\LPREV \land (\lnot \FLUSH_0 \lor \THREAD_2^{\text{aux}})$}`
116 and 1 100 115     `\lstCMT{$\LPREV \land (\lnot \FLUSH_1 \lor \FLUSH_0^{\text{aux}})$}`
117 and 1 101 116     `\lstCMT{$\LPREV \land (\lnot \FLUSH_0^{\text{aux}} \lor \FLUSH_1^{\text{aux}})$}`
118 and 1 102 117     `\lstCMT{$\LPREV \land (\lnot \FLUSH_1 \lor \lnot \FLUSH_0^{\text{aux}})$}`
119 and 1 103 118     `\lstCMT{$\LPREV \land (\lnot \FLUSH_2 \lor \FLUSH_1^{\text{aux}})$}`
120 and 1 104 119     `\lstCMT{$\LPREV \land (\lnot \FLUSH_1^{\text{aux}} \lor \FLUSH_2^{\text{aux}})$}`
121 and 1 105 120     `\lstCMT{$\LPREV \land (\lnot \FLUSH_2 \lor \lnot \FLUSH_1^{\text{aux}})$}`
122 and 1 106 121     `\lstCMT{$\LPREV \land (\lnot \EXIT \lor \lnot \FLUSH_2^{\text{aux}})$}`
\end{lstlisting}
and adding yet another invariant.
\begin{lstlisting}[style=btor2]
123 constraint 122
\end{lstlisting}
% `\lstCMT{$\leq^1_7\!(\THREAD_0, \ldots, \FLUSH_0, \ldots, \EXIT)$}`

\noindent
% Store buffer constraints.
% The other store buffer,
% The remaining scheduling constraints can be defined rather effortlessly.
% \bigbreak
% \noindent
% $\bullet \;$ Disable flushing empty store buffer and executing barrier instructions while it is full:
% Transitions are then further restricted by the remaining scheduling constraints like
% The number of valid transitions is then further restricted by our remaining scheduling constraints like
This cardinality constraint is then again further influenced
by prohibiting certain transitions like
% Possible transitions are then further restricted by our common set of scheduling constraints like
% Again, execution is further restricted,
% Next, we also introduce our common scheduling constraints like
flushing an empty store buffer or executing barrier instructions if it is full,
\begin{lstlisting}[style=btor2]
124 or 1 29 32        `\lstCMT{$\STMT_{0, 1} \lor \STMT_{0, 4}$}`
125 implies 1 124 -53 `\lstCMT{$\LPREV \implies \lnot \THREAD_0$}`
126 ite 1 25 125 -56  `\lstCMT{$\ITE(\SBFULL_0, \;\LPREV\:, \lnot \FLUSH_0)$}`
127 constraint 126

128 or 1 34 37        `\lstCMT{$\STMT_{1, 1} \lor \STMT_{1, 4}$}`
129 implies 1 128 -54 `\lstCMT{$\LPREV \implies \lnot \THREAD_1$}`
130 ite 1 26 129 -57  `\lstCMT{$\ITE(\SBFULL_1, \;\LPREV\:, \lnot \FLUSH_1)$}`
131 constraint 130

132 implies 1 58 27   `\lstCMT{$\FLUSH_2 \implies \SBFULL_2$}`
133 constraint 132
\end{lstlisting}

\noindent
% Checkpoint constraints.
executing a thread while it is waiting for a checkpoint,
\begin{lstlisting}[style=btor2]
134 and 1 44 -76      `\lstCMT{$\BLOCK_{0, 0} \land \lnot \CHECK_0$}`
135 implies 1 134 -53 `\lstCMT{$\LPREV \implies \lnot \THREAD_0$}`
136 constraint 135

137 and 1 45 -76      `\lstCMT{$\BLOCK_{0, 1} \land \lnot \CHECK_0$}`
138 implies 1 137 -54 `\lstCMT{$\LPREV \implies \lnot \THREAD_1$}`
139 constraint 138

140 and 1 46 -76      `\lstCMT{$\BLOCK_{0, 2} \land \lnot \CHECK_0$}`
141 implies 1 140 -55 `\lstCMT{$\LPREV \implies \lnot \THREAD_2$}`
142 constraint 141
\end{lstlisting}

\noindent
% Halt constraints.
and stopping halted threads from being scheduled.
% and halt related constraints ...
\begin{lstlisting}[style=btor2]
143 implies 1 47 -53  `\lstCMT{$\HALT_0 \implies \lnot \THREAD_0$}`
144 constraint 143

145 implies 1 48 -54  `\lstCMT{$\HALT_1 \implies \lnot \THREAD_1$}`
146 constraint 145

147 implies 1 49 -55  `\lstCMT{$\HALT_2 \implies \lnot \THREAD_2$}`
148 constraint 147
\end{lstlisting}

\noindent
% Accumulator state definitions.
As mentioned earlier, the main advantage of the {\BTOR} format lies in the ability to define states in combination with their transition function.
% In case of our machine states therefore follows a pretty basic pattern:
In our case, definition of machine states therefore follows a pretty basic pattern:
% Our machine state can therefore be defined by following a pretty basic pattern:
% This feature allows us to define our machine states
% after initializing the state by using the keyword \lstBTOR{init},
after initializing the state with the keyword \lstBTOR{init},
% and specifying its successor through the
% and adding the set of nodes for defining its successor,
% its transition function is implemented by the according set of nodes and
% its transition function is defined through the appropriate set of nodes and registered with the corresponding keyword \lstBTOR{next}.
its successor is defined through the appropriate set of nodes and
the resulting expression
registered as the given state's transition function by using the corresponding keyword \lstBTOR{next}.
% and defining its successor through the appropriate set of nodes,
% and adding the set of nodes yielding its successor,
% and specified to be the specific states transition function through the keyword \lstBTOR{next}.
% the specific state's transition function can be registered with accompanying keyword \lstBTOR{next}.
% Due to the reuse of nodes, the encoding is much more compact and given below for completeness.
\bigbreak
\noindent
$\bullet\ $ Accumulator registers:
\begin{lstlisting}[style=btor2]
149 init 2 13 6       `\lstCMT{$\ACCU^0_0 \leftarrow 0$}`
150 add 2 13 7        `\lstCMT{$\ACCU_0 + 1$}`
151 ite 2 59 150 13   `\lstCMT{$\ITE(\EXEC_{0, 0}, \,\LPREV\,, \ACCU_0) \equiv \texttt{ADDI}_0$}`
152 read 2 50 7       `\lstCMT{$\READ(1)$}`
153 eq 1 19 7         `\lstCMT{$\SBADR_0 = 1$}`
154 and 1 25 153      `\lstCMT{$\LPREV \land \SBFULL_0$}`
155 ite 2 154 22 152  `\lstCMT{$\ITE(\,\LPREV\,, \SBVAL_0, \READ(1))$}`
156 ite 2 61 155 151  `\lstCMT{$\ITE(\EXEC_{0, 2}, \,\LPREV\,, \, \texttt{ADDI}_0)$}`
157 next 2 13 156     `\lstCMT{$\ACCU'_0 \leftarrow \RPREV$}`

158 init 2 14 6       `\lstCMT{$\ACCU^0_1 \leftarrow 0$}`
159 add 2 14 7        `\lstCMT{$\ACCU_1 + 1$}`
160 ite 2 64 159 14   `\lstCMT{$\ITE(\EXEC_{1, 0}, \,\LPREV\,, \ACCU_1) \equiv \texttt{ADDI}_1$}`
161 read 2 50 6       `\lstCMT{$\READ(0)$}`
162 eq 1 20 6         `\lstCMT{$\SBADR_1 = 0$}`
163 and 1 26 162      `\lstCMT{$\LPREV \land \SBFULL_1$}`
164 ite 2 163 23 161  `\lstCMT{$\ITE(\,\LPREV\,, \SBVAL_1, \READ(0))$}`
165 ite 2 66 164 160  `\lstCMT{$\ITE(\EXEC_{1, 2}, \,\LPREV\,, \, \texttt{ADDI}_1)$}`
166 next 2 14 165     `\lstCMT{$\ACCU'_1 \leftarrow \RPREV$}`

167 init 2 15 6       `\lstCMT{$\ACCU^0_2 \leftarrow 0$}`
168 eq 1 21 6         `\lstCMT{$\SBADR_2 = 0$}`
169 and 1 27 168      `\lstCMT{$\LPREV \land \SBFULL_2$}`
170 ite 2 169 24 161  `\lstCMT{$\ITE(\,\LPREV\,, \SBVAL_2, \READ(0))$}`
171 add 2 15 170      `\lstCMT{$\LPREV + \ACCU_2$}`
172 ite 2 70 171 15   `\lstCMT{$\ITE(\EXEC_{2, 1}, \,\LPREV\,, \ACCU_2) \equiv \texttt{ADD}$}`
173 eq 1 21 7         `\lstCMT{$\SBADR_2 = 1$}`
174 and 1 27 173      `\lstCMT{$\LPREV \land \SBFULL_2$}`
175 ite 2 174 24 152  `\lstCMT{$\ITE(\,\LPREV\,, \SBVAL_2, \READ(1))$}`
176 add 2 15 175      `\lstCMT{$\LPREV + \ACCU_2$}`
177 ite 2 71 176 172  `\lstCMT{$\ITE(\EXEC_{2, 2}, \,\LPREV\,, \texttt{ADD})$}`
178 next 2 15 177     `\lstCMT{$\ACCU'_2 \leftarrow \RPREV$}`
\end{lstlisting}
                      % `\lstCMT{$\ITE(\EXEC_{2, 2}, \,\LPREV\,, \ITE(\EXEC_{2, 1}, \ITE(\,\LPREV\,, \SBVAL_2, \HEAP[0]), \ACCU_2))$}`
                      % `\lstCMT{\hspace{0.8cm}$\LPREV,$}`
                      % `\lstCMT{\hspace{0.8cm}$\ITE(\EXEC_{2, 1},$}`

\noindent
$\bullet\ $ CAS memory registers:
\begin{lstlisting}[style=btor2]
179 init 2 16 6       `\lstCMT{$\MEM^0_0 \leftarrow 0$}`
180 next 2 16 16      `\lstCMT{$\MEM'_0 \leftarrow \MEM_0$}`

181 init 2 17 6       `\lstCMT{$\MEM^0_1 \leftarrow 0$}`
182 next 2 17 17      `\lstCMT{$\MEM'_1 \leftarrow \MEM_1$}`

183 init 2 18 6       `\lstCMT{$\MEM^0_2 \leftarrow 0$}`
184 next 2 18 18      `\lstCMT{$\MEM'_2 \leftarrow \MEM_2$}`
\end{lstlisting}

\noindent
$\bullet\ $ Store buffer address registers:
\begin{lstlisting}[style=btor2]
185 init 2 19 6       `\lstCMT{$\SBADR^0_0 \leftarrow 0$}`
186 ite 2 60 6 19     `\lstCMT{$\ITE(\EXEC_{0, 1}, 0, \SBADR_0)$}`
187 next 2 19 186     `\lstCMT{$\SBADR'_0 \leftarrow \RPREV$}`

188 init 2 20 6       `\lstCMT{$\SBADR^0_1 \leftarrow 0$}`
189 ite 2 65 7 20     `\lstCMT{$\ITE(\EXEC_{1, 1}, 0, \SBADR_1)$}`
190 next 2 20 189     `\lstCMT{$\SBADR'_1 \leftarrow \RPREV$}`

191 init 2 21 6       `\lstCMT{$\SBADR^0_2 \leftarrow 0$}`
192 next 2 21 21      `\lstCMT{$\SBADR'_2 \leftarrow \SBADR_2$}`
\end{lstlisting}

\noindent
$\bullet\ $ Store buffer value registers:
\begin{lstlisting}[style=btor2]
193 init 2 22 6       `\lstCMT{$\SBVAL^0_0 \leftarrow 0$}`
194 ite 2 60 13 22    `\lstCMT{$\ITE(\EXEC_{0, 1}, \ACCU_0, \SBVAL_0)$}`
195 next 2 22 194     `\lstCMT{$\SBVAL'_0 \leftarrow \RPREV$}`

196 init 2 23 6       `\lstCMT{$\SBVAL^0_1 \leftarrow 0$}`
197 ite 2 65 14 23    `\lstCMT{$\ITE(\EXEC_{1, 1}, \ACCU_1, \SBVAL_1)$}`
198 next 2 23 197     `\lstCMT{$\SBVAL'_1 \leftarrow \RPREV$}`

199 init 2 24 6       `\lstCMT{$\SBVAL^0_1 \leftarrow 0$}`
200 next 2 24 24      `\lstCMT{$\SBVAL'_2 \leftarrow \SBVAL_2$}`
\end{lstlisting}

\noindent
$\bullet\ $ Store buffer full flags:
\begin{lstlisting}[style=btor2]
201 init 1 25 4       `\lstCMT{$\SBFULL^0_0 \leftarrow 0$}`
202 or 1 60 25        `\lstCMT{$\EXEC_{0, 1} \lor \SBFULL_0$}`
203 ite 1 56 4 202    `\lstCMT{$\ITE(\FLUSH_0, 0, \,\LPREV\,)$}`
204 next 1 25 203     `\lstCMT{$\SBFULL'_0 \leftarrow \RPREV$}`

205 init 1 26 4       `\lstCMT{$\SBFULL^0_1 \leftarrow 0$}`
206 or 1 65 26        `\lstCMT{$\EXEC_{1, 1} \lor \SBFULL_1$}`
207 ite 1 57 4 206    `\lstCMT{$\ITE(\FLUSH_1, 0, \,\LPREV\,)$}`
208 next 1 26 207     `\lstCMT{$\SBFULL'_1 \leftarrow \RPREV$}`

209 init 1 27 4       `\lstCMT{$\SBFULL^0_2 \leftarrow 0$}`
210 ite 1 58 4 27     `\lstCMT{$\ITE(\FLUSH_2, 0, \SBFULL_2)$}`
211 next 1 27 210     `\lstCMT{$\SBFULL'_2 \leftarrow \RPREV$}`
\end{lstlisting}

\noindent
$\bullet\ $ Statement activation flags:
\begin{lstlisting}[style=btor2]
212 init 1 28 5       `\lstCMT{$\STMT^0_{0, 0} \leftarrow 1$}`
213 and 1 28 -59      `\lstCMT{$\STMT_{0, 0} \land \lnot \EXEC_{0, 0}$}`
214 next 1 28 213     `\lstCMT{$\STMT'_{0, 0} \leftarrow \RPREV$}`

215 init 1 29 4       `\lstCMT{$\STMT^0_{0, 1} \leftarrow 0$}`
216 and 1 29 -60      `\lstCMT{$\STMT_{0, 1} \land \lnot \EXEC_{0, 1}$}`
217 ite 1 28 59 216   `\lstCMT{$\ITE(\STMT_{0, 0}, \EXEC_{0, 0}, \,\RPREV\,)$}`
218 next 1 29 217     `\lstCMT{$\STMT'_{0, 1} \leftarrow \RPREV$}`

219 init 1 30 4       `\lstCMT{$\STMT^0_{0, 2} \leftarrow 0$}`
220 and 1 30 -61      `\lstCMT{$\STMT_{0, 2} \land \lnot \EXEC_{0, 2}$}`
221 ite 1 29 60 220   `\lstCMT{$\ITE(\STMT_{0, 1}, \EXEC_{0, 1}, \,\RPREV\,)$}`
222 next 1 30 221     `\lstCMT{$\STMT'_{0, 2} \leftarrow \RPREV$}`

223 init 1 31 4       `\lstCMT{$\STMT^0_{0, 3} \leftarrow 0$}`
224 and 1 31 -62      `\lstCMT{$\STMT_{0, 3} \land \lnot \EXEC_{0, 3}$}`
225 ite 1 30 61 224   `\lstCMT{$\ITE(\STMT_{0, 2}, \EXEC_{0, 2}, \,\RPREV\,)$}`
226 next 1 31 225     `\lstCMT{$\STMT'_{0, 3} \leftarrow \RPREV$}`

227 init 1 32 4       `\lstCMT{$\STMT^0_{0, 4} \leftarrow 0$}`
228 and 1 32 -63      `\lstCMT{$\STMT_{0, 4} \land \lnot \EXEC_{0, 4}$}`
229 ite 1 31 62       `\lstCMT{$\ITE(\STMT_{0, 3}, \EXEC_{0, 3}, \,\RPREV\,)$}`
230 next 1 32 229     `\lstCMT{$\STMT'_{0, 4} \leftarrow \RPREV$}`

231 init 1 33 5       `\lstCMT{$\STMT^0_{1, 0} \leftarrow 1$}`
232 and 1 33 -64      `\lstCMT{$\STMT_{1, 0} \land \lnot \EXEC_{1, 0}$}`
233 next 1 33 232     `\lstCMT{$\STMT'_{1, 0} \leftarrow \RPREV$}`

234 init 1 34 4       `\lstCMT{$\STMT^0_{1, 1} \leftarrow 0$}`
235 and 1 34 -65      `\lstCMT{$\STMT_{1, 1} \land \lnot \EXEC_{1, 1}$}`
236 ite 1 33 64 235   `\lstCMT{$\ITE(\STMT_{1, 0}, \EXEC_{1, 0}, \,\RPREV\,)$}`
237 next 1 34 236     `\lstCMT{$\STMT'_{1, 1} \leftarrow \RPREV$}`

238 init 1 35 4       `\lstCMT{$\STMT^0_{1, 2} \leftarrow 0$}`
239 and 1 35 -66      `\lstCMT{$\STMT_{1, 2} \land \lnot \EXEC_{1, 2}$}`
240 ite 1 34 65 239   `\lstCMT{$\ITE(\STMT_{1, 1}, \EXEC_{1, 1}, \,\RPREV\,)$}`
241 next 1 35 240     `\lstCMT{$\STMT'_{1, 2} \leftarrow \RPREV$}`

242 init 1 36 4       `\lstCMT{$\STMT^0_{1, 3} \leftarrow 0$}`
243 and 1 36 -67      `\lstCMT{$\STMT_{1, 3} \land \lnot \EXEC_{1, 3}$}`
244 ite 1 35 66 243   `\lstCMT{$\ITE(\STMT_{1, 2}, \EXEC_{1, 2}, \,\RPREV\,)$}`
245 next 1 36 244     `\lstCMT{$\STMT'_{1, 3} \leftarrow \RPREV$}`

246 init 1 37 4       `\lstCMT{$\STMT^0_{1, 4} \leftarrow 0$}`
247 and 1 37 -68      `\lstCMT{$\STMT_{1, 4} \land \lnot \EXEC_{0, 4}$}`
248 ite 1 36 67 247   `\lstCMT{$\ITE(\STMT_{1, 3}, \EXEC_{1, 3}, \,\RPREV\,)$}`
249 next 1 37 248     `\lstCMT{$\STMT'_{1, 4} \leftarrow \RPREV$}`

250 init 1 38 5       `\lstCMT{$\STMT^0_{2, 0} \leftarrow 1$}`
251 and 1 38 -69      `\lstCMT{$\STMT_{2, 0} \land \lnot \EXEC_{2, 0}$}`
252 next 1 38 251     `\lstCMT{$\STMT'_{2, 0} \leftarrow \RPREV$}`

253 init 1 39 4       `\lstCMT{$\STMT^0_{2, 1} \leftarrow 0$}`
254 and 1 39 -70      `\lstCMT{$\STMT_{2, 1} \land \lnot \EXEC_{2, 1}$}`
255 ite 1 38 69 254   `\lstCMT{$\ITE(\STMT_{2, 0}, \EXEC_{2, 0}, \,\RPREV\,)$}`
256 next 1 39 255     `\lstCMT{$\STMT'_{2, 1} \leftarrow \RPREV$}`

257 init 1 40 4       `\lstCMT{$\STMT^0_{2, 2} \leftarrow 0$}`
258 and 1 40 -71      `\lstCMT{$\STMT_{2, 2} \land \lnot \EXEC_{2, 2}$}`
259 ite 1 39 70 258   `\lstCMT{$\ITE(\STMT_{2, 1}, \EXEC_{2, 1}, \,\RPREV\,)$}`
260 next 1 40 259     `\lstCMT{$\STMT'_{2, 2} \leftarrow \RPREV$}`

261 init 1 41 4       `\lstCMT{$\STMT^0_{2, 3} \leftarrow 0$}`
262 and 1 41 -72      `\lstCMT{$\STMT_{2, 3} \land \lnot \EXEC_{2, 3}$}`
263 ite 1 40 71 262   `\lstCMT{$\ITE(\STMT_{2, 2}, \EXEC_{2, 2}, \,\RPREV\,)$}`
264 next 1 41 263     `\lstCMT{$\STMT'_{2, 3} \leftarrow \RPREV$}`

265 init 1 42 4       `\lstCMT{$\STMT^0_{2, 4} \leftarrow 0$}`
266 and 1 42 -73      `\lstCMT{$\STMT_{2, 4} \land \lnot \EXEC_{2, 4}$}`
267 eq 1 15 6         `\lstCMT{$\ACCU_2 = 0$}`
268 and 1 72 -267     `\lstCMT{$\EXEC_{2, 3} \land \ACCU_2 \neq 0$}`
269 ite 1 41 268 266  `\lstCMT{$\ITE(\STMT_{2, 3}, \,\LPREV\,, \STMT_{2, 4} \land \lnot \EXEC_{2, 4})$}`
270 next 1 42 269     `\lstCMT{$\STMT'_{2, 4} \leftarrow \RPREV$}`

271 init 1 43 4       `\lstCMT{$\STMT^0_{2, 5} \leftarrow 0$}`
272 and 1 43 -74      `\lstCMT{$\STMT_{2, 5} \land \lnot \EXEC_{2, 5}$}`
273 and 1 72 267      `\lstCMT{$\EXEC_{2, 3} \land \ACCU_2 = 0$}`
274 ite 1 41 273 272  `\lstCMT{$\ITE(\STMT_{2, 3}, \,\LPREV\,, \STMT_{2, 5} \land \lnot \EXEC_{2, 5})$}`
275 next 1 43 274     `\lstCMT{$\STMT'_{2, 5} \leftarrow \RPREV$}`
\end{lstlisting}

\noindent
$\bullet\ $ Block flags:
\begin{lstlisting}[style=btor2]
276 init 1 44 4       `\lstCMT{$\BLOCK^0_{0, 0} \leftarrow 0$}`
277 or 1 62 44        `\lstCMT{$\EXEC_{0, 3} \lor \BLOCK_{0, 0}$}`
278 ite 1 76 4 277    `\lstCMT{$\ITE(\CHECK_0, 0, \,\RPREV\,)$}`
279 next 1 44 278     `\lstCMT{$\BLOCK'_{0, 0} \leftarrow \RPREV$}`

280 init 1 45 4       `\lstCMT{$\BLOCK^0_{0, 1} \leftarrow 0$}`
281 or 1 67 45        `\lstCMT{$\EXEC_{1, 3} \lor \BLOCK_{0, 1}$}`
282 ite 1 76 4 281    `\lstCMT{$\ITE(\CHECK_0, 0, \,\RPREV\,)$}`
283 next 1 45 282     `\lstCMT{$\BLOCK'_{0, 1} \leftarrow \RPREV$}`

284 init 1 46 4       `\lstCMT{$\BLOCK^0_{0, 2} \leftarrow 0$}`
285 or 1 69 46        `\lstCMT{$\EXEC_{2, 0} \lor \BLOCK_{0, 2}$}`
286 ite 1 76 4 285    `\lstCMT{$\ITE(\CHECK_0, 0, \,\RPREV\,)$}`
287 next 1 46 286     `\lstCMT{$\BLOCK'_{0, 2} \leftarrow \RPREV$}`
\end{lstlisting}

\noindent
$\bullet\ $ Halt flags:
\begin{lstlisting}[style=btor2]
288 init 1 47 4       `\lstCMT{$\HALT^0_0 \leftarrow 0$}`
289 or 1 63 47        `\lstCMT{$\EXEC_{0, 4} \lor \HALT_0$}`
290 next 1 47 289     `\lstCMT{$\HALT'_0 \leftarrow \RPREV$}`

291 init 1 48 4       `\lstCMT{$\HALT^0_1 \leftarrow 0$}`
292 or 1 68 48        `\lstCMT{$\EXEC_{1, 4} \lor \HALT_1$}`
293 next 1 48 292     `\lstCMT{$\HALT'_1 \leftarrow \RPREV$}`

294 init 1 49 4       `\lstCMT{$\HALT^0_2 \leftarrow 0$}`
295 next 1 49 49      `\lstCMT{$\HALT'_2 \leftarrow \HALT_2$}`
\end{lstlisting}

\noindent
$\bullet\ $ Shared memory:
\begin{lstlisting}[style=btor2]
296 init 3 50 12      `\lstCMT{$\HEAP^0 \leftarrow \texttt{mmap}$}`
297 write 3 50 19 22  `\lstCMT{$\WRITE(\SBADR_0, \SBVAL_0)$}`
298 ite 3 56 297 50   `\lstCMT{$\ITE(\FLUSH_0, \,\LPREV\,, \HEAP) \equiv \texttt{FLUSH}_0$}`
299 write 3 50 20 23  `\lstCMT{$\WRITE(\SBADR_1, \SBVAL_1)$}`
300 ite 3 57 299 298  `\lstCMT{$\ITE(\FLUSH_1, \,\LPREV\,, \texttt{FLUSH}_0) \equiv \texttt{FLUSH}_1$}`
301 write 3 50 21 24  `\lstCMT{$\WRITE(\SBADR_2, \SBVAL_2)$}`
302 ite 3 58 301 300  `\lstCMT{$\ITE(\FLUSH_2, \,\LPREV\,, \texttt{FLUSH}_1)$}`
303 next 3 50 302     `\lstCMT{$\HEAP' \leftarrow \RPREV$}`
\end{lstlisting}

\noindent
$\bullet\ $ Exit flag:
\begin{lstlisting}[style=btor2]
304 init 1 51 4       `\lstCMT{$\EXIT^0 \leftarrow 0$}`
305 and 1 289 292     `\lstCMT{$\HALT'_0 \land \HALT'_1$}`
306 and 1 49 305      `\lstCMT{$\LPREV \land \HALT_2$}`
307 or 1 51 306       `\lstCMT{$\LPREV \lor \EXIT$}`
308 or 1 73 307       `\lstCMT{$\LPREV \lor \EXEC_{2, 4}$}`
309 or 1 74 308       `\lstCMT{$\LPREV \lor \EXEC_{2, 5}$}`
310 next 1 51 309     `\lstCMT{$\EXIT' \leftarrow \RPREV$}`
\end{lstlisting}

\noindent
$\bullet\ $ Exit code:
\begin{lstlisting}[style=btor2]
311 init 2 52 6       `\lstCMT{$\EXITCODE^0 \leftarrow 0$}`
312 ite 2 73 6 52     `\lstCMT{$\ITE(\EXEC_{2, 4}, 0, \EXITCODE)$}`
313 ite 2 74 7 312    `\lstCMT{$\ITE(\EXEC_{2, 5}, 1, \,\RPREV\,)$}`
314 next 2 52 313     `\lstCMT{$\EXITCODE' \leftarrow \RPREV$}`
\end{lstlisting}

\noindent
Finally, we declare an exit code greater than zero as the only \lstBTOR{bad} state.
\begin{lstlisting}[style=btor2]
315 neq 1 6 52
316 bad 315
\end{lstlisting}

% \newpage
%
% \begin{minipage}[t]{0.45\textwidth}
% \begin{lstlisting}[style=btor2]
% 1 sort bitvec 1
% 2 sort bitvec 16
% 3 sort array 2 2
%
% 4 zero 1
% 5 one 1
%
% 6 zero 2
% 7 one 2
% 8 constd 2 5
% 9 constd 2 17
%
% 10 state 3 mmap
% 11 write 3 10 6 6
% 12 write 3 11 7 6
%
% 13 state 2 accu_0
% 14 state 2 accu_1
% 15 state 2 accu_2
%
% 16 state 2 mem_0
% 17 state 2 mem_1
% 18 state 2 mem_2
%
% 19 state 2 sb-adr_0
% 20 state 2 sb-adr_1
% 21 state 2 sb-adr_2
%
% 22 state 2 sb-val_0
% 23 state 2 sb-val_1
% 24 state 2 sb-val_2
%
% 25 state 1 sb-full_0
% 26 state 1 sb-full_1
% 27 state 1 sb-full_2
%
% 28 state 1 stmt_0_0
% 29 state 1 stmt_0_1
% 30 state 1 stmt_0_2
% 31 state 1 stmt_0_3
% 32 state 1 stmt_0_4
% \end{lstlisting}
% \end{minipage}
% \begin{minipage}[t]{0.45\textwidth}
% \begin{lstlisting}[style=btor2]
% 33 state 1 stmt_1_0
% 34 state 1 stmt_1_1
% 35 state 1 stmt_1_2
% 36 state 1 stmt_1_3
% 37 state 1 stmt_1_4
%
% 38 state 1 stmt_2_0
% 39 state 1 stmt_2_1
% 40 state 1 stmt_2_2
% 41 state 1 stmt_2_3
% 42 state 1 stmt_2_4
% 43 state 1 stmt_2_5
%
% 44 state 1 block_0_0
% 45 state 1 block_0_1
% 46 state 1 block_0_2
%
% 47 state 1 halt_0
% 48 state 1 halt_1
% 49 state 1 halt_2
%
% 50 state 3 heap
%
% 51 state 1 exit
%
% 52 state 2 exit-code
%
% 53 input 1 thread_0
% 54 input 1 thread_1
% 55 input 1 thread_2
%
% 56 input 1 flush_0
% 57 input 1 flush_1
% 58 input 1 flush_2
%
% 59 and 1 28 53 exec_0_0
% 60 and 1 29 53 exec_0_1
% 61 and 1 30 53 exec_0_2
% 62 and 1 31 53 exec_0_3
% 63 and 1 32 53 exec_0_4
% \end{lstlisting}
% \end{minipage}

% \subsubsection{BTOR2 generator functions}
% \todo[inline]{BTOR2 generator functions} %%%%%%%%%%%%%%%%%%%%%%%%%%%%%%%%%%%%%%%
%
% \begin{lstlisting}[style=c++]
% template <class ... T>
% std::string expr (const std::string & nid, const T & ... args)
% {
  % std::string e = nid;
  % (((e += ' ') += args), ...);
  % auto & end = e.back();
  % if (end == ' ')
    % end = eol;
  % else
    % e += eol;
  % return e;
% }
% \end{lstlisting}
%
% \todo[inline]{main encoding function (smtlib)} %%%%%%%%%%%%%%%%%%%%%%%%%%%%%%%%%
%
% \begin{lstlisting}[style=c++]
% void encode ()
% {
  % declare_sorts();
  % declare_constants();
  % define_mmap();
  % declare_states();
  % declare_inputs();
  % define_transitions();
  % define_constraints();
  % define_states();
% }
% \end{lstlisting}
%
% \todo[inline]{declare sorts} %%%%%%%%%%%%%%%%%%%%%%%%%%%%%%%%%%%%%%%%%%%%%%%%%%%
%
% \begin{lstlisting}[style=c++]
% void declare_sorts ()
% {
  % formula <<
    % sort_bitvec(sid_bool = nid(), "1") <<
    % sort_bitvec(sid_bv = nid(), std::to_string(word_size)) <<
    % sort_array(sid_heap = nid(), "2", "2") <<
    % eol;
% }
% \end{lstlisting}



\newpage

\section{Testing and Debugging}

The basic functionality of our toolchain is validated by a total of 800 test cases.
While fixing {\CPP} related errors was relatively trivial due to the broad range of available tools,
debugging the generated SMT encodings
% on the other hand
% requires more sophisticated methods.
is a bit more tedious
since cause and effect have to be determined manually by inspecting rather comprehensive models and formulas.

In order to automatically validate the correctness of our encodings, we included a \texttt{replay} \CHANGE{mode} for reevaluating the execution sequence of a given trace by the virtual machine used for simulation.
% Therefore, we included a \texttt{replay} \CHANGE{mode} for reevaluating the execution sequence of a given trace by the virtual machine used for simulation in order to automatically validate the correctness of our encodings.
If a mismatch is found, it returns the first step in which the results start to differ and prints the corresponding states of both execution traces for comparison.
% prints the states of the first step in which the results start to differ and returns
% printing the first step in which the traces start to differ
% returning step in which the resulting models start to differ and print
This mode also simplified the encodings' development significantly,
% since the step in which the model starts to diverge from our machine model
% as it automatically finds the step in which the encoding starts to differ from our machine model and
since the condensed error trace helps to narrow down the potentially flawed SMT expressions.
% since it helps to narrow down the potentially erroneous SMT expressions by

Of course, only satisfiable formulas may be replayed.
% reevaluated by \texttt{replay}.
If the outcome of a valid execution sequence is accidentally unsatisfiable and therefore yields no trace, we still have to manually inspect the formula.
Unfortunately, no SMT solver supported the extraction of an \emph{unsatisfiable core} at the time we developed our encodings, but this feature starts to get implemented in recent solvers like Bitwuzla \cite{ref:Bitwuzla}.
However, without any additional constraints, the only reason our formulas can get unsatisfiable is due to a violation of the cardinality constraint and investigations may therefore be concentrated on expressions containing transition variables.

Another frequent cause for unintended results lies in the nature of bounded model checking.
Before expecting an error in the encoding, it is highly advisable to make sure that the targeted states are indeed reachable by the chosen bound.
For example, when a model for the \lstASM{CAS} variant of our concurrent counter experiments (described in the next section) was found during our initial trials, we immediately started to panic and questioned our encoding.
Since such situations already commonly occurred during development, we remembered the following proverb which emerged during that time and
soon realized that the problem indeed was yet another insufficient bound.
% later found the problem to be yet another insufficient bound.
% since we questioned our encoding which was previously believed to be corre
% only to find the problem being an insufficient bound.
% resulting in the following proverb:
\begin{center}
\itshape
If you ask yourself -- ``How can this be?'' -- maybe just the bound's too wee!
\end{center}


\newpage

\newcommand{\EncoderStatsTable}[3]{
  \begin{table}[!hbt]
  \noindent\makebox[\textwidth]{
  \footnotesize
  \centering
  \begin{tabu}{|l|X[r]|X[r]|X[r]|X[r]|X[r]|X[r]|}
    \tabucline{2-}
    \multicolumn{1}{c|}{}
      & \multicolumn{2}{c|}{\textbf{btor2}}
      & \multicolumn{2}{c|}{\textbf{functional}}
      & \multicolumn{2}{c|}{\textbf{relational}} \\
    \tabucline{2-}
    \multicolumn{1}{c|}{}
      & \multicolumn{1}{c|}{time}
      & \multicolumn{1}{c|}{size}
      & \multicolumn{1}{c|}{time}
      & \multicolumn{1}{c|}{size}
      & \multicolumn{1}{c|}{time}
      & \multicolumn{1}{c|}{size} \\
    \tabucline{2-}
    \firsthline
    \csvreader
      [late after line=\\, separator=pipe]
      {#1}
      {1=\A, 2=\B, 3=\C, 4=\D, 5=\E, 6=\F, 7=\G}
      {\ref{tbl:\A} & \B & \C & \D & \E & \F & \G}
    \lasthline
  \end{tabu}}
  \label{#2}
  \caption{#3}
  \end{table}
}

\newcommand{\SolverStatsTable}[3]{
  \begin{table}[!hbt]
  \noindent\makebox[\textwidth]{
  \footnotesize
  \centering
  \begin{tabu}{|l|X[r]|X[r]|X[r]|X[r]|X[r]|X[r]|X[r]|}
    \tabucline{2-}
    \multicolumn{1}{c|}{}
      & \multicolumn{1}{c|}{\textbf{BtorMC}}
      & \multicolumn{2}{c|}{\textbf{Boolector}}
      & \multicolumn{2}{c|}{\textbf{Z3}}
      & \multicolumn{2}{c|}{\textbf{CVC4}} \\
    \tabucline{2-}
    \multicolumn{1}{c|}{}
      & \multicolumn{1}{c|}{\scriptsize btor2}
      & \multicolumn{1}{c|}{\scriptsize functional}
      & \multicolumn{1}{c|}{\scriptsize relational}
      & \multicolumn{1}{c|}{\scriptsize functional}
      & \multicolumn{1}{c|}{\scriptsize relational}
      & \multicolumn{1}{c|}{\scriptsize functional}
      & \multicolumn{1}{c|}{\scriptsize relational} \\
    \tabucline{2-}
    \firsthline
    \csvreader
      [late after line=\\, separator=pipe]
      {#1}
      {1=\A, 2=\B, 3=\C, 4=\D, 5=\E, 6=\F, 7=\G, 8=\H}
      {\ref{tbl:\A} & \B & \C & \D & \E & \F & \G & \H}
    \lasthline
  \end{tabu}}
  \label{#2}
  \caption{#3}
  \end{table}
}

%------------------------------------------------------------------------------%

\section{Experiments}

\subsection{LitmusTests}

\todo[inline]{table - encoding stats - intel}

\EncoderStatsTable
  {figures/litmus-intel-encoder.dat}
  {tbl:litmus:intel:encoder}
  {Intel Litmus Test Econding Time [ms] and Size [\#expressions]}

\todo[inline]{table - solving times - intel}

\SolverStatsTable
  {figures/litmus-intel-solver.dat}
  {tbl:litmus:intel:solver}
  {Intel Litmus Test Solving Times in Seconds}

\newpage

\todo[inline]{table - encoding stats - amd}

\EncoderStatsTable
  {figures/litmus-amd-encoder.dat}
  {tbl:litmus:amd:encoder}
  {AMD Litmus Test Econding Time [ms] and Size [\#expressions]}

\todo[inline]{table - solving times - amd}

\SolverStatsTable
  {figures/litmus-amd-solver.dat}
  {tbl:litmus:amd:solver}
  {AMD Litmus Test Solving Times in Seconds}

\subsection{Statistical Counter}


\newpage

\section{Conclusion}

% haben gezeigt dass unser Ansatz funktioniert ...
% schön wäre, wenn man einen derartigen Verifikations Schritt in Compiler Toolchains integrieren könnte um damit Binaries automatisiert auf derartige Fehler zu überprüfen ...
% aber leider nicht praktikabel, da viel zu lange Laufzeit ...
% Future Work: Encoding optimieren, Laufzeit minimieren, evtl. auf "echte" Prozessoren anwenden ...

In this thesis we explored the potential of bounded model checking for verifying concurrent programs
including the target architecture's memory ordering habits.
Due to the complexity of actual processors, we devised a highly simplified virtual machine model
% to deliver a proof-of-concept
% We therefore devised a highly simplified virtual machine model
% imitating the behaviour of current x86 implementations.
and showed that the addition of store buffers is sufficient to imitate the behaviour of current x86 implementations.
% An encoding for this specific machine model was developed.
% Execution of arbitrary programs
% An encoding for the execution of arbitrary programs was developed

After designing a formal framework for the execution of arbitrary programs on this virtual machine model,
% encodings in two different SMT formats were developed and
encodings of the corresponding model checking problems in two different SMT formats were developed.
% we developed a tool for automating the encoding and evaluation of the resulting model checking problems in two different SMT formats
% To simplify experimentation
We then implemented ConcuBinE, a tool for automating the encoding and evaluation of generated SMT formulas using state-of-the-art solvers.

Even though experiments show the basic feasibility of our approach, the relatively larger runtimes for even small examples puts its practicability in question.
% However, we would like to see the inclusion of such a verification step in compiler toolchains
% However, we still believe it as an opportunity to integrate such a verification step into compiler toolchains to fully proof the
However, we still see it as an opportunity to integrate a similar checking procedure into compiler toolchains, automatically verifying the binaries of critical software with regards to the target architecture's memory ordering model and other specifics.

Future work could therefore involve optimizing the encoding,
% designing an intermediate language for specifying assertions,
devising a method for specifying assertions,
developing a procedure to additionally compute the WCET for automatically determining the upper bound and applying the presented principles to a real architecture like Intel's or AMD's x86 processors.


\newpage

\bibliographystyle{plain}
\bibliography{references}

\newpage

\appendix

\counterwithin{figure}{subsection}
\counterwithin{table}{subsection}
\counterwithin{lstlisting}{subsection}

\section*{Appendices}

% \section*{Appendix}

\addcontentsline{toc}{section}{Appendices}
\renewcommand{\thesubsection}{\Alph{subsection}}

\renewcommand{\tablename}{Example}
% \setcounter{table}{0}

\subsection{Intel Litmus Tests}

\label{appendix:litmus:intel}

Intel's memory-ordering litmus tests as seen in \cite[Section 8.2.3]{ref:Intel} ported to our virtual machine model.

% The x86-TSO memory-ordering model (used by AMD and Intel) is defined by a set of litmus tests.
% To show equivalence of ConcuBinE-TSO and x86-TSO, the tests have been ported to ConcuBinE.

% \subsubsection{Intel}

% {\small \textit{see Intel® 64 and IA-32 Architectures Software Developer’s Manual Volume 3A}}

% \textbf{\cite[8.2.3.2]{ref:Intel}}

\subsubsection*{Neither Loads Nor Stores Are Reordered with Like Operations}

The Intel-64 memory-ordering model allows neither loads nor stores to be reordered with the same kind of operation.
That is, it ensures that loads are seen in program order and that stores are seen in program order.
This is illustrated by the following example:

\begin{table}[!hbt]
\noindent\makebox[\textwidth]{
\centering
\begin{tabu}{c}
  \begin{tabu}{|l|l|}
    \firsthline
    \textbf{Thread 0} & \textbf{Thread 1} \\
    \hline \hline
    \lstASM{ADDI 1}   & \\
    \lstASM{STORE 0}  & \lstASM{MEM 1} \\
    \lstASM{STORE 1}  & \lstASM{LOAD 0} \\
    \lasthline
  \end{tabu} \\
  \begin{tabu}{|X[c]|X[c]|}
    \firsthline
    \textbf{Initial} & \textbf{Not Allowed} \hspace{1pt} \xmark \\
    \hline \hline
    $\HEAP[0] = \HEAP[1] = 0$ & $\MEM_1 = 1 \land \ACCU_1 = 0$ \\
    \lasthline
  \end{tabu} \\
\end{tabu}}
\caption{Stores Are Not Reordered with Other Stores \cite[Example 8-1]{ref:Intel}}
\label{tbl:litmus:intel:1}
\end{table}

\noindent
The disallowed return values could be exhibited only if thread 0’s two stores are reordered (with the two loads occurring between them) or if thread 1’s two loads are reordered (with the two stores occurring between them).
\bigbreak
\noindent
If $\MEM_1 = 1$, the store to $\HEAP[1]$ occurs before the load from $\HEAP[1]$.
Because the Intel-64 memory-ordering model does not allow stores to be reordered, the earlier store to $\HEAP[0]$ occurs before the load from $\HEAP[1]$.
Because the Intel-64 memory-ordering model does not allow loads to be reordered, the store to $\HEAP[0]$ also occurs before the later load from $\HEAP[0]$.
Thus $\ACCU_1 = 1$.

\newpage

\subsubsection*{Stores Are Not Reordered with Earlier Loads}

The Intel-64 memory-ordering model ensures that a store by a thread may not occur before a previous load by the same thread.
This is illustrated by the following example:

\begin{table}[!hbt]
\noindent\makebox[\textwidth]{
\centering
\begin{tabu}{c}
  \begin{tabu}{|l|l|}
    \firsthline
    \textbf{Thread 0} & \textbf{Thread 1} \\
    \hline \hline
    \lstASM{MEM 0}    & \lstASM{MEM 1} \\
    \lstASM{ADDI 1}   & \lstASM{ADDI 1} \\
    \lstASM{STORE 1}  & \lstASM{STORE 0} \\
    \lasthline
  \end{tabu} \\
  \begin{tabu}{|X[c]|X[c]|}
    \firsthline
    \textbf{Initial} & \textbf{Not Allowed} \hspace{1pt} \xmark \\
    \hline \hline
    $\HEAP[0] = \HEAP[1] = 0$ & $\MEM_0 = \MEM_1 = 1$ \\
    \lasthline
  \end{tabu} \\
\end{tabu}}
\caption{Stores Are Not Reordered with Older Loads \cite[Example 8-2]{ref:Intel}}
\label{tbl:litmus:intel:2}
\end{table}

\noindent
Assume $\MEM_0 = 1$.
\begin{itemize}
  \item Because $\MEM_0 = 1$, thread 1’s store to $\HEAP[0]$ occurs before thread 0’s load from $\HEAP[0]$.
  \item Because the Intel-64 memory-ordering model prevents each store from being reordered with the earlier load by the same thread, thread 1’s load from $\HEAP[1]$ occurs before its store to $\HEAP[0]$.
  \item Similarly, thread 0’s load from $\HEAP[0]$ occurs before its store to $\HEAP[1]$.
  \item Thus, thread 1’s load from $\HEAP[1]$ occurs before thread 0’s store to $\HEAP[1]$, implying $\MEM_1 = 0$.
\end{itemize}

\newpage

\subsubsection*{Loads May Be Reordered with Earlier Stores to Different Locations}

The Intel-64 memory-ordering model allows a load to be reordered with an earlier store to a different location.
However, loads are not reordered with stores to the same location.
\bigbreak
\noindent
The fact that a load may be reordered with an earlier store to a different location is illustrated by the following example:

\begin{table}[!hbt]
\noindent\makebox[\textwidth]{
\centering
\begin{tabu}{c}
  \begin{tabu}{|l|l|}
    \firsthline
    \textbf{Thread 0} & \textbf{Thread 1} \\
    \hline \hline
    \lstASM{ADDI 1}   & \lstASM{ADDI 1} \\
    \lstASM{STORE 0}  & \lstASM{STORE 1} \\
    \lstASM{LOAD 1}   & \lstASM{LOAD 0} \\
    \lasthline
  \end{tabu} \\
  \begin{tabu}{|X[c]|X[c]|}
    \firsthline
    \textbf{Initial} & \textbf{Allowed} \hspace{1pt} \cmark \\
    \hline \hline
    $\HEAP[0] = \HEAP[1] = 0$ & $\ACCU_0 = \ACCU_1 = 0$ \\
    \lasthline
  \end{tabu} \\
\end{tabu}}
\caption{Loads May be Reordered with Older Stores \cite[Example 8-3]{ref:Intel}}
\label{tbl:litmus:intel:3}
\end{table}

\noindent
At each thread, the load and the store are to different locations and hence may be reordered.
Any interleaving of the operations is thus allowed.
One such interleaving has the two loads occurring before the two stores.
This would result in each load returning value 0.

\newpage

\subsubsection*{Loads Are Not Reordered with Older Stores to the Same Location}

The Intel-64 memory-ordering model allows a load to be reordered with an earlier store to a different location.
However, loads are not reordered with stores to the same location.
\bigbreak
\noindent
The fact that a load may not be reordered with an earlier store to the same location is illustrated by the following
example:

\begin{table}[!hbt]
\noindent\makebox[\textwidth]{
\centering
\begin{tabu}{c}
  \begin{tabu}{|l|}
    \firsthline
    \textbf{Thread 0} \\
    \hline \hline
    \lstASM{ADDI 1} \\
    \lstASM{STORE 0} \\
    \lstASM{LOAD 0} \\
    \lasthline
  \end{tabu} \\
  \begin{tabu}{|X[c]|X[c]|}
    \firsthline
    \textbf{Initial} & \textbf{Not Allowed} \hspace{1pt} \xmark \\
    \hline \hline
    $\HEAP[0] = 0$ & $\ACCU_0 = 0$ \\
    \lasthline
  \end{tabu} \\
\end{tabu}}
\caption{Loads Are not Reordered with Older Stores to the Same Location \cite[Example 8-4]{ref:Intel}}
\label{tbl:litmus:intel:4}
\end{table}

\noindent
The Intel-64 memory-ordering model does not allow the load to be reordered with the earlier store because the accesses are to the same location.
Therefore, $\ACCU_0 = 1$ must hold.

\newpage

\subsubsection*{Intra-Processor Forwarding Is Allowed}

The memory-ordering model allows concurrent stores by two threads to be seen in different orders by those two threads; specifically, each thread may perceive its own store occurring before that of the other.
This is illustrated by the following example:

\begin{table}[!hbt]
\noindent\makebox[\textwidth]{
\centering
\begin{tabu}{c}
  \begin{tabu}{|l|l|}
    \firsthline
    \textbf{Thread 0} & \textbf{Thread 1} \\
    \hline \hline
    \lstASM{ADDI 1}   & \lstASM{ADDI 1} \\
    \lstASM{STORE 0}  & \lstASM{STORE 1} \\
    \lstASM{LOAD 0}   & \lstASM{LOAD 1} \\
    \lstASM{LOAD 1}   & \lstASM{LOAD 0} \\
    \lasthline
  \end{tabu} \\
  \begin{tabu}{|X[c]|X[c]|}
    \firsthline
    \textbf{Initial} & \textbf{Allowed} \hspace{1pt} \cmark \\
    \hline \hline
    $\HEAP[0] = \HEAP[1] = 0$ & $\ACCU_0 = \ACCU_1 = 0$ \\
    \lasthline
  \end{tabu} \\
\end{tabu}}
\caption{Intra-Processor Forwarding is Allowed \cite[Example 8-5]{ref:Intel}}
\label{tbl:litmus:intel:5}
\end{table}

\noindent
The memory-ordering model imposes no constraints on the order in which the two stores appear to execute by the two threads.
This fact allows thread 0 to see its store before seeing thread 1's, while thread 1 sees its store before seeing thread 0's.
(Each thread is self consistent.)
This allows $\ACCU_0 = \ACCU_1 = 0$.
\bigbreak
\noindent
In practice, the reordering in this example can arise as a result of store-buffer forwarding.
While a store is temporarily held in a thread's store buffer, it can satisfy the thread's own loads but is not visible to (and cannot satisfy) loads by other threads.

\newpage

\subsubsection*{Stores Are Transitively Visible}

The memory-ordering model ensures transitive visibility of stores; stores that are causally related appear to all threads to occur in an order consistent with the causality relation.
This is illustrated by the following example:

\begin{table}[!hbt]
\noindent\makebox[\textwidth]{
\centering
\begin{tabu}{c}
  \begin{tabu}{|l|l|l|}
    \firsthline
    \textbf{Thread 0} & \textbf{Thread 1} & \textbf{Thread 2} \\
    \hline \hline
    \lstASM{ADDI 1}   & \lstASM{MEM 0}    & \\
    \lstASM{STORE 0}  & \lstASM{JNZ 3}    & \\
                      & \lstASM{ADDI 1}   & \\
                      & \lstASM{STORE 1}  & \lstASM{MEM 1} \\
                      &                   & \lstASM{LOAD 0} \\
    \lasthline
  \end{tabu} \\
  \begin{tabu}{|X[c]|X[c]|}
    \firsthline
    \textbf{Initial} & \textbf{Not Allowed} \hspace{1pt} \xmark \\
    \hline \hline
    $\HEAP[0] = \HEAP[1] = 0$ & $\MEM_1 = \MEM_2 = 1 \land \ACCU_2 = 0$ \\
    \lasthline
  \end{tabu} \\
\end{tabu}}
\caption{Stores Are Transitively Visible \cite[Example 8-6]{ref:Intel}}
\label{tbl:litmus:intel:6}
\end{table}

\noindent
Assume that $\MEM_1 = 1$ and $\MEM_2 = 1$.
\begin{itemize}
  \item Because $\MEM_1 = 1$, thread 0’s store occurs before thread 1’s load.
  \item Because the memory-ordering model prevents a store from being reordered with an earlier load (see \cite[Section 8.2.3.3]{ref:Intel}), thread 1’s load occurs before its store. Thus, thread 0’s store causally precedes thread 1’s store.
  \item Because thread 0’s store causally precedes thread 1’s store, the memory-ordering model ensures that thread 0’s store appears to occur before thread 1’s store from the point of view of all threads.
  \item Because $\MEM_2 = 1$, thread 1’s store occurs before thread 2’s load.
  \item Because the Intel-64 memory-ordering model prevents loads from being reordered (see \cite[Section 8.2.3.2]{ref:Intel}), thread 2’s load occur in order.
  \item The above items imply that thread 0’s store to $\HEAP[0]$ occurs before thread 2’s load from $\HEAP[0]$. This implies that $\ACCU_2 = 1$.
\end{itemize}

\newpage

\subsubsection*{Stores Are Seen in a Consistent Order by Other Threads}

As noted in \cite[Section 8.2.3.5]{ref:Intel}, the memory-ordering model allows stores by two threads to be seen in different orders by those two threads.
However, any two stores must appear to execute in the same order to all threads other than those performing the stores.
This is illustrated by the following example:

\begin{table}[!hbt]
\noindent\makebox[\textwidth]{
\centering
\begin{tabu}{c}
  \begin{tabu}{|l|l|l|l|}
    \firsthline
    \textbf{Thread 0} & \textbf{Thread 1} & \textbf{Thread 2} & \textbf{Thread 3} \\
    \hline \hline
    \lstASM{ADDI 1}   & \lstASM{ADDI 1}   & \lstASM{MEM 0}    & \lstASM{MEM 1} \\
    \lstASM{STORE 0}  & \lstASM{STORE 1}  & \lstASM{LOAD 1}   & \lstASM{LOAD 0} \\
    \lasthline
  \end{tabu} \\
  \begin{tabu}{|X[c]|X[2.14,c]|}
    \firsthline
    \textbf{Initial} & \textbf{Not Allowed} \hspace{1pt} \xmark \\
    \hline \hline
    $\HEAP[0] = \HEAP[1] = 0$ & $\MEM_2 = \MEM_3 = 1 \land \ACCU_2 = \ACCU_3 = 0$ \\
    \lasthline
  \end{tabu} \\
\end{tabu}}
\caption{Stores Are Seen in a Consistent Order by Other Threads \cite[Example 8-7]{ref:Intel}}
\label{tbl:litmus:intel:7}
\end{table}

\noindent
By the principles discussed in \cite[Section 8.2.3.2]{ref:Intel},
\begin{itemize}
  \item thread 2’s first and second load cannot be reordered,
  \item thread 3’s first and second load cannot be reordered.
  \item If $\MEM_2 = 1$ and $\ACCU_2 = 0$, thread 0’s store appears to precede thread 1’s store with respect to thread 2.
  \item Similarly, $\MEM_3 = 1$ and $\ACCU_3 = 0$ imply that thread 1’s store appears to precede thread 0’s store with respect to thread 1.
\end{itemize}
Because the memory-ordering model ensures that any two stores appear to execute in the same order to all threads (other than those performing the stores), this set of return values is not allowed.

\newpage

\subsubsection*{Locked Instructions Have a Total Order}

The memory-ordering model ensures that all threads agree on a single execution order of all locked instructions.
This is illustrated by the following example:

\begin{table}[!hbt]
\noindent\makebox[\textwidth]{
\centering
\begin{tabu}{c}
  \begin{tabu}{|l|l|l|l|}
    \firsthline
    \textbf{Thread 0} & \textbf{Thread 1} & \textbf{Thread 2} & \textbf{Thread 3} \\
    \hline \hline
    \lstASM{ADDI 1}   & \lstASM{ADDI 1}   &                   & \\
    \lstASM{CAS 0}    & \lstASM{CAS 1}    &                   & \\
                      &                   & \lstASM{MEM 0}    & \lstASM{MEM 1} \\
                      &                   & \lstASM{LOAD 1}   & \lstASM{LOAD 0} \\
    \lasthline
  \end{tabu} \\
  \begin{tabu}{|X[c]|X[2.14,c]|}
    \firsthline
    \textbf{Initial} & \textbf{Not Allowed} \hspace{1pt} \xmark \\
    \hline \hline
    $\HEAP[0] = \HEAP[1] = 0$ & $\MEM_2 = \MEM_3 = 1 \land \ACCU_2 = \ACCU_3 = 0$ \\
    \lasthline
  \end{tabu} \\
\end{tabu}}
\caption{Locked Instructions Have a Total Order \cite[Example 8-8]{ref:Intel}}
\label{tbl:litmus:intel:8}
\end{table}

\noindent
Thread 2 and thread 3 must agree on the order of the two executions of \lstASM{CAS}.
Without loss of generality, suppose that thread 0’s \lstASM{CAS} occurs first.
\begin{itemize}
  \item If $\MEM_3 = 1$, thread 1’s \lstASM{CAS} into $\HEAP[1]$ occurs before thread 3’s load from $\HEAP[1]$.
  \item Because the Intel-64 memory-ordering model prevents loads from being reordered (see \cite[Section 8.2.3.2]{ref:Intel}), thread 3’s loads occur in order and, therefore, thread 1’s \lstASM{CAS} occurs before thread 3’s load from $\HEAP[0]$.
  \item Since thread 0’s \lstASM{CAS} into $\HEAP[0]$ occurs before thread 1’s \lstASM{CAS} (by assumption), it occurs before thread 3’s load from $\HEAP[0]$.
  Thus, $\ACCU_3 = 1$.
\end{itemize}
A similar argument (referring instead to thread 2’s loads) applies if thread 1’s \lstASM{CAS} occurs before thread 0’s \lstASM{CAS}.

\newpage

\subsubsection*{Loads Are Not Reordered with Locked Instructions}

The memory-ordering model prevents loads and stores from being reordered with locked instructions that execute earlier or later.
The examples in this section illustrate only cases in which a locked instruction is executed before a load or a store.
The reader should note that reordering is prevented also if the locked instruction is executed after a load or a store.
\bigbreak
\noindent
The first example illustrates that loads may not be reordered with earlier locked instructions:

\begin{table}[!hbt]
\noindent\makebox[\textwidth]{
\centering
\begin{tabu}{c}
  \begin{tabu}{|l|l|l|l|}
    \firsthline
    \textbf{Thread 0} & \textbf{Thread 1} \\
    \hline \hline
    \lstASM{ADDI 1}   & \lstASM{ADDI 1} \\
    \lstASM{CAS 0}    & \lstASM{CAS 1} \\
    \lstASM{LOAD 1}   & \lstASM{LOAD 0} \\
    \lasthline
  \end{tabu} \\
  \begin{tabu}{|X[c]|X[c]|}
    \firsthline
    \textbf{Initial} & \textbf{Not Allowed} \hspace{1pt} \xmark \\
    \hline \hline
    $\HEAP[0] = \HEAP[1] = 0$ & $\ACCU_0 = \ACCU_1 = 0$ \\
    \lasthline
  \end{tabu} \\
\end{tabu}}
\caption{Loads Are not Reordered with Locks \cite[Example 8-9]{ref:Intel}}
\label{tbl:litmus:intel:9}
\end{table}

\noindent
As explained in \cite[Section 8.2.3.8]{ref:Intel}, there is a total order of the executions of locked instructions.
Without loss of generality, suppose that thread 0’s \lstASM{CAS} occurs first.
\bigbreak
\noindent
Because the Intel-64 memory-ordering model prevents thread 1’s load from being reordered with its earlier \lstASM{CAS}, thread 0’s \lstASM{CAS} occurs before thread 1’s load.
This implies $\ACCU_1 = 1$.
\bigbreak
\noindent
A similar argument (referring instead to thread 2’s accesses) applies if thread 1’s \lstASM{CAS} occurs before thread 0’s \lstASM{CAS}.

\newpage

\subsubsection*{Stores Are Not Reordered with Locked Instructions}

The memory-ordering model prevents loads and stores from being reordered with locked instructions that execute earlier or later.
The examples in this section illustrate only cases in which a locked instruction is executed before a load or a store.
The reader should note that reordering is prevented also if the locked instruction is executed after a load or a store.
\bigbreak
\noindent
The second example illustrates that a store may not be reordered with an earlier locked instruction:

\begin{table}[!hbt]
\noindent\makebox[\textwidth]{
\centering
\begin{tabu}{c}
  \begin{tabu}{|l|l|l|l|}
    \firsthline
    \textbf{Thread 0} & \textbf{Thread 1} \\
    \hline \hline
    \lstASM{ADDI 1}   & \\
    \lstASM{CAS 0}    & \lstASM{MEM 1} \\
    \lstASM{STORE 1}  & \lstASM{LOAD 0} \\
    \lasthline
  \end{tabu} \\
  \begin{tabu}{|X[c]|X[c]|}
    \firsthline
    \textbf{Initial} & \textbf{Not Allowed} \hspace{1pt} \xmark \\
    \hline \hline
    $\HEAP[0] = \HEAP[1] = 0$ & $\ACCU_1 = 0 \land \MEM_1 = 1$ \\
    \lasthline
  \end{tabu} \\
\end{tabu}}
\caption{Stores Are not Reordered with Locks \cite[Example 8-10]{ref:Intel}}
\label{tbl:litmus:intel:10}
\end{table}

\noindent
Assume $\MEM_1 = 1$.
\begin{itemize}
  \item Because $\MEM_1 = 1$, thread 0’s store to $\HEAP[1]$ occurs before thread 1’s load from $\HEAP[1]$.
  \item Because the memory-ordering model prevents a store from being reordered with an earlier locked instruction, thread 0’s \lstASM{CAS} into $\HEAP[0]$ occurs before its store to $\HEAP[1]$.
  \item Thus, thread 0’s \lstASM{CAS} into $\HEAP[0]$ occurs before thread 1’s load from $\HEAP[1]$.
  \item Because the memory-ordering model prevents loads from being reordered (see \cite[Section 8.2.3.2]{ref:Intel}), thread 1’s loads occur in order and, therefore, thread 1’s \lstASM{CAS} into $\HEAP[0]$ occurs before thread 1’s load from $\HEAP[0]$. Thus, $\ACCU_1 = 1$.
\end{itemize}

\newpage

\subsection{AMD Litmus Tests}

\label{appendix:litmus:amd}

AMD's memory-ordering litmus tests as seen in \cite[Section 7.2]{ref:AMD} ported to our virtual machine model.

% \subsubsection*{Loads Do Not Pass Previous Loads, Stores Do Not Pass Previous Stores}
\subsubsection*{Neither Loads Nor Stores Are Reordered with Like Operations}

Successive stores from a single thread are committed to system memory and visible to other threads in program order.
A store by a thread cannot be committed to memory before a read appearing earlier in the program has captured its targeted data from memory.
In other words, stores from a thread cannot be reordered to occur prior to a load preceding it in program order.

\begin{table}[!hbt]
\noindent\makebox[\textwidth]{
\centering
\begin{tabu}{c}
  \begin{tabu}{|l|l|}
    \firsthline
    \textbf{Thread 0} & \textbf{Thread 1} \\
    \hline \hline
    \lstASM{ADDI 1}   & \\
    \lstASM{STORE 0}  & \lstASM{MEM 1} \\
    \lstASM{STORE 1}  & \lstASM{LOAD 0} \\
    \lasthline
  \end{tabu} \\
  \begin{tabu}{|X[c]|X[c]|}
    \firsthline
    \textbf{Initial} & \textbf{Not Allowed} \hspace{1pt} \xmark \\
    \hline \hline
    $\HEAP[0] = \HEAP[1] = 0$ & $\MEM_1 = 1 \land \ACCU_1 = 0$ \\
    \lasthline
  \end{tabu} \\
\end{tabu}}
\caption{Stores Are Not Reordered with Other Stores \cite[Example 1]{ref:AMD}}
\label{tbl:litmus:amd:1}
\end{table}

\noindent
\lstASM{LOAD 0} cannot read $0$ when \lstASM{LOAD 1} reads $1$.

\newpage

\subsubsection*{Stores Are Not Reordered with Earlier Loads}

Successive stores from a single thread are committed to system memory and visible to other threads in program order.
A store by a thread cannot be committed to memory before a read appearing earlier in the program has captured its targeted data from memory.
In other words, stores from a thread cannot be reordered to occur prior to a load preceding it in program order.

\begin{table}[!hbt]
\noindent\makebox[\textwidth]{
\centering
\begin{tabu}{c}
  \begin{tabu}{|l|l|}
    \firsthline
    \textbf{Thread 0} & \textbf{Thread 1} \\
    \hline \hline
    \lstASM{MEM 0}    & \lstASM{MEM 1} \\
    \lstASM{ADDI 1}   & \lstASM{ADDI 1} \\
    \lstASM{STORE 1}  & \lstASM{STORE 0} \\
    \lasthline
  \end{tabu} \\
  \begin{tabu}{|X[c]|X[c]|}
    \firsthline
    \textbf{Initial} & \textbf{Not Allowed} \hspace{1pt} \xmark \\
    \hline \hline
    $\HEAP[0] = \HEAP[1] = 0$ & $\MEM_0 = \MEM_1 = 1$ \\
    \lasthline
  \end{tabu} \\
\end{tabu}}
\caption{Stores Are Not Reordered with Older Loads \cite[Example 2]{ref:AMD}}
\label{tbl:litmus:amd:2}
\end{table}

\noindent
\lstASM{LOAD 0} and \lstASM{LOAD 1} cannot both read $1$.

\subsubsection*{Stores Can Be Arbitrarily Delayed}

Stores from a thread appear to be committed to the memory system in program order; however, stores can be delayed arbitrarily by store buffering while the thread continues operation.
Therefore, stores from a thread may not appear to be sequentially consistent.

\begin{table}[!hbt]
\noindent\makebox[\textwidth]{
\centering
\begin{tabu}{c}
  \begin{tabu}{|l|l|}
    \firsthline
    \textbf{Thread 0} & \textbf{Thread 1} \\
    \hline \hline
    \lstASM{ADDI 1}   & \lstASM{ADDI 1} \\
    \lstASM{STORE 0}  & \lstASM{STORE 1} \\
    \lstASM{ADDI 1}   & \lstASM{ADDI 1} \\
    \lstASM{STORE 0}  & \lstASM{STORE 1} \\
    \lstASM{LOAD 0}   & \lstASM{LOAD 1} \\
    \lasthline
  \end{tabu} \\
  \begin{tabu}{|X[c]|}
    \firsthline
    \textbf{Allowed} \hspace{1pt} \cmark \\
    \hline \hline
    $\ACCU_0 = \ACCU_1 = 1$ \\
    \lasthline
  \end{tabu} \\
\end{tabu}}
\caption{Stores Can Be Arbitrarily Delayed \cite[Example 3]{ref:AMD}}
\label{tbl:litmus:amd:3}
\end{table}

\noindent
Both \lstASM{LOAD 0} and \lstASM{LOAD 1} may read $1$.

\subsubsection*{Loads May Be Reordered with Earlier Stores to Different Locations}

Non-overlapping Loads may pass stores.

\begin{table}[!hbt]
\noindent\makebox[\textwidth]{
\centering
\begin{tabu}{c}
  \begin{tabu}{|l|l|}
    \firsthline
    \textbf{Thread 0} & \textbf{Thread 1} \\
    \hline \hline
    \lstASM{ADDI 1}   & \lstASM{ADDI 1} \\
    \lstASM{STORE 0}  & \lstASM{STORE 1} \\
    \lstASM{LOAD 1}   & \lstASM{LOAD 0} \\
    \lasthline
  \end{tabu} \\
  \begin{tabu}{|X[c]|X[c]|}
    \firsthline
    \textbf{Initial} & \textbf{Allowed} \hspace{1pt} \cmark \\
    \hline \hline
    $\HEAP[0] = \HEAP[1] = 0$ & $\ACCU_0 = \ACCU_1 = 0$ \\
    \lasthline
  \end{tabu} \\
\end{tabu}}
\caption{Loads May be Reordered with Older Stores \cite[Example 4]{ref:AMD}}
\label{tbl:litmus:amd:4}
\end{table}

\noindent
All combinations of values ($00$, $01$, $10$, and $11$) may be observed by threads 0 and 1.

\subsubsection*{Sequential Consistency}

Where sequential consistency is needed (for example in Dekker’s algorithm for mutual exclusion), a \lstASM{FENCE} instruction should be used between the store and the subsequent load, or an atomic instruction, such as \lstASM{CAS}, should be used for the store.

\begin{table}[!hbt]
\noindent\makebox[\textwidth]{
\centering
\begin{tabu}{c}
  \begin{tabu}{|l|l|}
    \firsthline
    \textbf{Thread 0} & \textbf{Thread 1} \\
    \hline \hline
    \lstASM{ADDI 1}   & \lstASM{ADDI 1} \\
    \lstASM{STORE 0}  & \lstASM{STORE 1} \\
    \lstASM{FENCE}    & \lstASM{FENCE} \\
    \lstASM{LOAD 1}   & \lstASM{LOAD 0} \\
    \lasthline
  \end{tabu} \\
  \begin{tabu}{|X[c]|X[c]|}
    \firsthline
    \textbf{Initial} & \textbf{Not Allowed} \hspace{1pt} \xmark \\
    \hline \hline
    $\HEAP[0] = \HEAP[1] = 0$ & $\ACCU_0 = \ACCU_1 = 0$ \\
    \lasthline
  \end{tabu} \\
\end{tabu}}
\caption{Sequential Consistency \cite[Example 5]{ref:AMD}}
\label{tbl:litmus:amd:5}
\end{table}

\noindent
\lstASM{LOAD 0} and \lstASM{LOAD 1} cannot both read $0$.

\newpage

\subsubsection*{Stores Are Seen in a Consistent Order by Other Threads}

Stores to different locations in memory observed from two (or more) other threads will appear in the same order to all observers.
Behavior such as that is shown in this code example:

\begin{table}[!hbt]
\noindent\makebox[\textwidth]{
\centering
\begin{tabu}{c}
  \begin{tabu}{|l|l|l|l|}
    \firsthline
    \textbf{Thread 0} & \textbf{Thread 1} & \textbf{Thread 2} & \textbf{Thread 3} \\
    \hline \hline
    \lstASM{ADDI 1}   & \lstASM{ADDI 1}   & \lstASM{MEM 0}    & \lstASM{MEM 1} \\
    \lstASM{STORE 0}  & \lstASM{STORE 1}  & \lstASM{LOAD 1}   & \lstASM{LOAD 0} \\
    \lasthline
  \end{tabu} \\
  \begin{tabu}{|X[c]|X[2.14,c]|}
    \firsthline
    \textbf{Initial} & \textbf{Not Allowed} \hspace{1pt} \xmark \\
    \hline \hline
    $\HEAP[0] = \HEAP[1] = 0$ & $\MEM_2 = \MEM_3 = 1 \land \ACCU_2 = \ACCU_3 = 0$ \\
    \lasthline
  \end{tabu} \\
\end{tabu}}
\caption{Stores Are Seen in a Consistent Order by Other Threads \cite[Example 6]{ref:AMD}}
\label{tbl:litmus:amd:6}
\end{table}

\noindent
Thread 2 seeing \lstASM{STORE 0} from thread 0 before \lstASM{STORE 1} from thread 1, while thread 3 sees \lstASM{STORE 1} from thread 1 before \lstASM{STORE 0} from thread 0, is not allowed.

\subsubsection*{Stores Are Transitively Visible}

Dependent stores between different threads appear to occur in program order, as shown in the code example below.

\begin{table}[!hbt]
\noindent\makebox[\textwidth]{
\centering
\begin{tabu}{c}
  \begin{tabu}{|l|l|l|}
    \firsthline
    \textbf{Thread 0} & \textbf{Thread 1} & \textbf{Thread 2} \\
    \hline \hline
    \lstASM{ADDI 1}   & \lstASM{MEM 0}    & \\
    \lstASM{STORE 0}  & \lstASM{JNZ 3}    & \\
                      & \lstASM{ADDI 1}   & \\
                      & \lstASM{STORE 1}  & \lstASM{MEM 1} \\
                      &                   & \lstASM{LOAD 0} \\
    \lasthline
  \end{tabu} \\
  \begin{tabu}{|X[c]|X[c]|}
    \firsthline
    \textbf{Initial} & \textbf{Not Allowed} \hspace{1pt} \xmark \\
    \hline \hline
    $\HEAP[0] = \HEAP[1] = 0$ & $\MEM_1 = \MEM_2 = 1 \land \ACCU_2 = 0$ \\
    \lasthline
  \end{tabu} \\
\end{tabu}}
\caption{Stores Are Transitively Visible \cite[Example 7]{ref:AMD}}
\label{tbl:litmus:amd:7}
\end{table}

\noindent
If thread 1 reads a value from $\HEAP[0]$ (written by thread 0) before carrying out a store to $\HEAP[1]$, and if thread 2 reads the updated value from $\HEAP[1]$, a subsequent read of $\HEAP[0]$ must also be the updated value.

\subsubsection*{Intra-Processor Forwarding Is Allowed}

The local visibility (within a thread) for a memory operation may differ from the global visibility (from another thread).
Using a data bypass, a local load can read the result of a local store in a store buffer, before the store becomes globally visible.
Program order is still maintained when using such bypasses.

\begin{table}[!hbt]
\noindent\makebox[\textwidth]{
\centering
\begin{tabu}{c}
  \begin{tabu}{|l|l|}
    \firsthline
    \textbf{Thread 0} & \textbf{Thread 1} \\
    \hline \hline
    \lstASM{ADDI 1}   & \lstASM{ADDI 1} \\
    \lstASM{STORE 0}  & \lstASM{STORE 1} \\
    \lstASM{MEM 0}    & \lstASM{MEM 1} \\
    \lstASM{LOAD 1}   & \lstASM{LOAD 0} \\
    \lasthline
  \end{tabu} \\
  \begin{tabu}{|X[c]|X[2.14,c]|}
    \firsthline
    \textbf{Initial} & \textbf{Allowed} \hspace{1pt} \cmark \\
    \hline \hline
    $\HEAP[0] = \HEAP[1] = 0$ & $\MEM_0 = \MEM_1 = 1 \land \ACCU_0 = \ACCU_1 = 0$ \\
    \lasthline
  \end{tabu} \\
\end{tabu}}
\caption{Intra-Processor Forwarding is Allowed \cite[Example 8]{ref:AMD}}
\label{tbl:litmus:amd:8}
\end{table}

\noindent
\lstASM{LOAD 0} in thread 0 can read $1$ using the data bypass, while \lstASM{LOAD 0} in thread 1 can read $0$.
Similarly, \lstASM{LOAD 1} in thread 1 can read $1$ while \lstASM{LOAD 1} in thread 0 can read $0$.
Therefore, the result $\MEM_0 = 1$, $\ACCU_0 = 0$, $\MEM_1 = 1$ and $\ACCU_1 = 0$ may occur.
There are no constraints on the relative order of when the \lstASM{STORE 0} of thread 0 is visible to thread 1 relative to when the \lstASM{STORE 1} of thread 1 is visible to thread 0.

\newpage

\subsubsection*{Global Visibility}

If a very strong memory ordering model is required that does not allow local store-load bypasses, a \lstASM{FENCE} instruction or an atomic instruction such as \lstASM{CAS} should be used between the store and the subsequent load.
This enforces a memory ordering stronger than total store ordering.

\begin{table}[!hbt]
\noindent\makebox[\textwidth]{
\centering
\begin{tabu}{c}
  \begin{tabu}{|l|l|}
    \firsthline
    \textbf{Thread 0} & \textbf{Thread 1} \\
    \hline \hline
    \lstASM{ADDI 1}   & \lstASM{ADDI 1} \\
    \lstASM{STORE 0}  & \lstASM{STORE 1} \\
    \lstASM{FENCE}    & \lstASM{FENCE} \\
    \lstASM{MEM 0}    & \lstASM{MEM 1} \\
    \lstASM{LOAD 1}   & \lstASM{LOAD 0} \\
    \lasthline
  \end{tabu} \\
  \begin{tabu}{|X[c]|X[2.14,c]|}
    \firsthline
    \textbf{Initial} & \textbf{Not Allowed} \hspace{1pt} \xmark \\
    \hline \hline
    $\HEAP[0] = \HEAP[1] = 0$ & $\MEM_0 = \MEM_1 = 1 \land \ACCU_0 = \ACCU_1 = 0$ \\
    \lasthline
  \end{tabu} \\
\end{tabu}}
\caption{Global Visibility \cite[Example 9]{ref:AMD}}
\label{tbl:litmus:amd:9}
\end{table}

\noindent
In this example, the \lstASM{FENCE} instruction ensures that any buffered stores are globally visible before the loads are allowed to execute, so the result $\MEM_0 = 1$, $\ACCU_0 = 0$, $\MEM_1 = 1$ and $\ACCU_1 = 0$ will not occur.


% \newpage

% \listoftodos

\end{document}
